% 本文件由 LaTeX 排版 Agent 根据有效 translation.json 自动生成。
% author=Ambrose; work_id=3401; title=论圣职人员的职责

\chapter{论圣职人员的职责}

% structure: p0001 source=h1 target=omitted reason=page-title-already-rendered-by-parent

卷一\newline{}卷二\newline{}卷三\par

\SourceRecord{Translated by H. de Romestin, E. de Romestin and H.T.F. Duckworth.；Nicene and Post-Nicene Fathers, Second Series；Vol. 10.；Edited by Philip Schaff and Henry Wace.；Buffalo, NY: Christian Literature Publishing Co.,；1896.；<http://www.newadvent.org/fathers/3401.htm>.}

% structure: p0001 source=h1 target=section reason=page-title

\section{\WorkTitle{论圣职人员的责任}}

\Person{圣安博罗削}极重视圣职的尊严，切望其教区的圣职人员活出与其崇高圣召相称的生活，成为信友善良且有益的榜样。因此他撰写了本论著，阐述圣职人员的责任，并以\Person{西塞罗}的\WorkTitle{论责任}为蓝本。\par

作者说他的目的在于，将先前教导过的道理深印于他所祝圣者的心中。一如\Person{西塞罗}，他论述了正直、合宜或尊贵的事（即\OriginalTerm{合宜}{decorum}），以及有益的事（即\OriginalTerm{利益}{utile}）；但并非针对此世，而是针对来世。在第一卷中他教导合宜或尊贵之事；在第二卷中讲论有益之事；在第三卷中则综合考量两者。\par

在第一卷中，他将责任分为\Quote{普通的}，即诫命的途径，对所有人都同样有约束力；以及\Quote{完美的}，在于遵行劝谕。继而讨论了一些基本的责任，如对父母和长者的责任之后，他触及了引领心灵的两项原则，即理性和欲望，并指明合宜在于思考良善和正直的事，以及使欲望服从理性，并提供了若干规则和榜样，最后以讨论四枢德，即智德、义德、勇德和节德作结。\par

在第二卷中，从合宜之事转向有益之事，他指出我们只能参照永生来衡量何为真正有益的，以此反驳异教哲学家的谬误，并表明有益之事在于认识天主及善度生活。其间他还指出，合宜之事其实就是有益之事，并以数章实际考量结束本卷。\par

在第三卷中，他论述了成全的责任，并定下规则：凡事必须询问什么是有益的，不是为个人，而是为众人或为所有人。凡不合宜之事，断不可追求；为此，不仅利益，甚至友谊和生命本身都必须退让。随后他以众多例子证明，圣人们如何追求合宜之事，并从而获得了有益之事。\par

\Person{圣安博罗削}以\Person{西塞罗}的论著为蓝本撰写本论的目的，似乎是为了驳斥异教的一些错误原则，并显示基督教道德远高于外邦人的道德。本论大约撰于公元391年。\par

% structure: p0008 source=h2 target=subsection reason=relative-heading-rank; base=h2

\subsection{第一章}

\Emph{主教的特殊职务是施教；然而\Person{圣安博罗削}为能施教，必须自己先学习；或者不如说，他必须教导他未曾学过的；无论如何，他的学与教必须同时并进。}\par

1. 我想，若我在自己的子女中间，顺从我施教的愿望，这不算过分自大，因为谦逊的导师自己曾说过：\BibleQuote{来，孩子们，听我训示：我要教导你们敬畏主。} 在此，既可见他对天主敬畏的谦逊，又可见其恩宠。因为他说\Quote{敬畏主}，这本是众人共有的，他以此描述了敬畏天主的主要标记。然而，敬畏本身是智慧的开端和真福的源泉——因为敬畏主的人是有福的——他明确地表明自己是教导智慧的导师，是达致真福的向导。\par

2. 因此，我们切望效法他对天主的敬畏，并且在分施恩宠上并非无据，就将智慧之神所传授于他，并通过他使我们明了，且借着眼见和榜样所学到的那些事，如同交付给子女一般交付给你们。因为如今我们再也不能逃避司铎职分的需要所加给我们的施教责任了，尽管我们曾试图回避：\BibleQuote{因为天主赐予一些人作宗徒，一些人作先知，一些人作传福音者，一些人作牧者和教师。}\BibleRef{弗 4:11}\par

3. 因此，我不为自己贪图宗徒的光荣（因为除了天主圣子亲自拣选的人，谁能有此光荣呢？）；也不贪图先知的恩宠、传福音者的德能，或牧者的审慎照顾。我只渴望在圣经上达到宗徒在圣者的职务中列于末位的那种关切与勤勉；\BibleRef{格前 12:10} 我正是渴望此事，好使我在努力施教时，能够学习。因为只有一位真正的导师，唯独祂教导一切，自己却从未学习；但人却必须先学习而后施教，并从祂接受，以传授于他人。\par

4. 但连这个情形也不适用于我。因为我是在审判席和\OriginalTerm{职服}{infulis}上被强拉下来，进入司铎职分，并开始教导你们我自己尚未学过的东西。如此，我是在开始学习之前就开始施教了。因此，我必须同时学习与施教，因为我此前没有闲暇学习。\par

% structure: p0014 source=h2 target=subsection reason=relative-heading-rank; base=h2

\subsection{第二章}

\Emph{说话会招致许多危险；圣经指示的补救方法在于缄默。}\par

5. 那么，我们在一切之前，首先应当学习的，岂不就是静默，好使我们能够说话吗？免得我的声音先定了我的罪，别人的声音才宣判我无罪；因为经上记载着：\BibleQuote{凭你的话，你将受到定罪。} \BibleRef{玛 12:37} 那么，你有什么必要急忙通过说话去冒受定罪的危险呢？当你保持静默更能安全时，你又何必呢？我曾见过多少人由于说话而陷入罪过，但几乎没见过一人因静默而如此；因此，懂得如何静默比懂得如何说话更难。我知道大多数人说话是因为他们不懂得如何静默。即使说话对人毫无益处，也很少有人保持静默。那么，晓得如何静默的人才是明智的。最后，天主的智慧说：\BibleQuote{主赐给了我受教的舌头，使我知道何时该说话。} 那么，从主那里领受了知道何时该说话的人，才真正称得上明智。因此圣经说得很好：\BibleQuote{明智人在时机来到之前保持静默。} \BibleRef{德 20:7}\par

6. 因此，主的圣人们喜爱静默，因为他们知道人的声音往往就是罪的表达，人的言谈便是人类过错的开始。最后，主的圣者说：\BibleQuote{我说过：我要谨守我的道路，免得我以口舌犯罪。} 因为他知道，并且也读过：一个人能免于自己舌头的鞭笞，以及自己良心的见证，乃是天主保护的标记；\BibleRef{约 5:21} 我们受到自己思想无声的谴责和良心的审判所责罚。我们也受到自己声音的鞭打，当我们说出使灵魂身受重伤、使心灵剧烈伤痛的话时。可是，有谁的心是洁净无罪的玷污，而在口舌上也不犯罪呢？因此，当他看到没有人能使自己的口不吐恶言时，他便通过静默的规则给自己制定了无辜的法律，目的是用静默来避免那种他难以在言谈中逃脱的过失。\par

7. 那么，让我们听从谨慎之师的教导吧：\BibleQuote{我说过：我要谨守我的道路；} 那也就是说，\BibleQuote{我自言自语：在我思想的无声命令中，我已嘱咐我自己，要谨守我的道路。} 有些道路是我们应当依循的，另有些道路则是我们应当留意的。我们必须依循主的道路，而对于自己的道路则要留意，免得它们引我们陷入罪过。一个人若不急于说话，就能留意。法律上说：\BibleQuote{以色列，你要听：主你的天主。} \BibleRef{申 6:4} 它没有说：\BibleQuote{说话，} 而是说：\BibleQuote{听。} \Person{厄娃}跌倒，是因为她对男人说了她未曾从主她的天主那里听到的话。天主对你说的第一句话就是：听吧！你若听了，就当留意你的道路；你若跌倒了，就赶快改正你的行径。因为：\BibleQuote{青年人如何才能改过自新？岂不是在于谨守天主的话吗？} 所以，首先要静默，并聆听，这样你的口舌才不会失足。\par

8. 一个人被自己的口舌定罪，实在是一桩大恶。的确，如果每个人都要为自己的闲话交账，\BibleRef{玛 12:36} 那么，对于污秽和不洁之言，岂不更要如此吗？因为仓促说出的话比闲话更加恶劣。因此，既然连一句闲话都要交代，那么，对于亵渎不敬的话语，岂不更要遭受严厉的惩罚吗？\par

% structure: p0020 source=h2 target=subsection reason=relative-heading-rank; base=h2

\subsection{第三章}

\Emph{静默不应始终不被打断，也不应出于懒散。如何防范内心和口舌免受过度偏情的侵害。}\par

9. 那么，又当如何？难道我们应该变成哑巴吗？当然不是。因为：\BibleQuote{静默有时，说话有时。} \BibleRef{训 3:7} 那么，如果我们连一句闲话都要交账的话，我们就该留意，免得起也因无益的静默而交账。因为还有一种积极的静默，就像\Person{苏撒纳}的静默那样：她藉着静默所做的，比她若开口说话所做的还要多。因为她在人前保持静默，却向天主说话，并发现没有比静默更能证明她的贞洁的了。她不用一言一语，她的良心却在说话；她不在人前为自己寻求判断，因为她有天主作证。所以她渴望由那位她深知绝不会受欺骗的天主宣判她无罪。是的，主自己在福音中也是默默地完成了人类的救恩工作。\BibleRef{玛 26:63} 因此，\Person{达味}恰当地给自己定下的，不是持久的静默，而是警惕。\par

10. 所以，我们要守护我们的心，守护我们的口。这两样在经上都有记载。在此处我们被告知要留意我们的口；在另一处又对你说：\BibleQuote{你要以全副心灵守护自己的心。} \BibleRef{箴 4:23} 如果\Person{达味}都留意了，难道你还不留意吗？如果\Person{依撒意亚}也曾嘴唇不洁——他曾说：\BibleQuote{我有祸了，我完了，因为我是个唇舌不洁的人} \BibleRef{依 6:5} ——如果一位主的先知尚且嘴唇不洁，我们又怎能洁净呢？\par

11. 可是，这是为谁写的呢？难道不是为我们每个人写的吗：\BibleQuote{要用荆棘围住你的产业，捆好你的金银，并为你的口装上门和门闩，为你的言语装上轭和天平？} \BibleRef{德 28:24} 你的产业就是你的心灵，你的金就是你的心，你的银就是你的言语：\BibleQuote{主的言语是纯净的言语，如同在火中炼过的白银。} 良好的心灵也是良好的产业。而且，纯洁的内修生活更是宝贵的产业。那么，你要把这产业围上篱笆，用思想把它围起来，用荆棘，也就是用虔敬的关怀守护它，免得肉身的强烈情欲冲进来，把它掳去；免得强猛的情绪攻击它，越过界限，夺去它的果实。要守护你的内在自我。不要轻忽或鄙视它，好像它毫无价值，因为它是一件贵重的产业；确实非常珍贵，因为它的果实不是短暂易朽、仅限于今世的，而是永恒的，用于永远的救恩。所以，要耕耘你的产业，使它成为你的耕田。\par

12. 要约束你的言语，免得它们放荡不羁，变得放肆，并在多言中为自己积聚罪过。宁可让它们受到限制，守在自身的堤岸之内。泛滥的河水很快便会聚拢污泥。也要约束你的心意；不可任其松懈放纵，免得有人指着你说：\Quote{没有治伤的香膏，没有油，也没有绷带可敷。} 心灵的自制有其辔头，藉以驾驭和引导。\par

13. 你的口当有一扇门，必要时可以关闭，并且要小心闩好，免得有人激起你的怒气，使你以辱骂还辱骂。今天你们听到的读经是：\BibleQuote{你们纵然动怒，却不可犯罪；}\BibleRef{弗 4:26} 因此，我们虽然动怒（这源于本性的冲动，而非意志的抉择），但不可口出恶言，以免我们陷于罪中；你们的言语当有轭和天平，即谦逊和节制，使你们的舌头服从于心志。要牢牢勒住舌头，让它有自制之方，能及时回心转意，趋于节制；让它说出经正义之天平衡量过的话语，使我们的心意庄重，言谈有分量，言语有分寸。\par

% structure: p0027 source=h2 target=subsection reason=relative-heading-rank; base=h2

\subsection{第四章}

\Emph{同样必须谨慎，使我们的言语不发自邪情，而出于善意；因为正是在这里，魔鬼特别伺机要捉住我们。}\par

14. 人若留意于此，便成为温和、良善、谦恭。因为他谨慎口舌，约束舌头，在未经过省察、思考和权衡之前——这话该不该说，那话该不该答，此时说这话是否合宜——便不轻易发言，他正是在操练谦恭、良善和忍耐。这样，他就不会因不快或愤怒而冲口发言，言语中不流露任何偏情，也不致使言辞间燃起欲火，或流露出愠怒的刺激。他应如此行事，以免那些本应为其内心生活增辉的话语，最终竟显露并证实其品行中存在某些恶习。\par

15. 因为敌人特别在我们内心生出偏情时策划计谋；那时他提供火种；那时他布置罗网。因此，正如我们今天听到的读经，先知并非无故地说：\BibleQuote{是他救我脱免猎户的罗网，脱免一切害人的恶言。}\BibleRef{咏 91:3} \Person{Symmachus}解释说，这是指\Quote{挑衅的言语}；另有人说是指\Quote{使人不安的言语}。敌人的罗网就是我们的言语——但这言语本身也同样是我们的大敌。我们常常说出一些话，被仇敌逮住，从而像用自己的剑一样刺伤自己。死于他人剑下，远胜于死在自己的剑下！\par

16. 因此，敌人试探我们的兵器，碰撞他的武器。他若见我心神不宁，便插入他的箭矢，以引发一连串争吵。我若说出不合宜的话，他便布下罗网。随后，他将报复的机会摆在我面前作为诱饵，好使我在渴望复仇时自投罗网，为自己系紧死亡的结索。人若察觉这敌人临近，就当格外谨慎口舌，以免给敌人留地步；但能看透他的人并不多。\par

% structure: p0032 source=h2 target=subsection reason=relative-heading-rank; base=h2

\subsection{第五章}

\Emph{我们也必须提防那看得见的敌人，当他以沉默激怒我们时；惟有藉着沉默之助，我们才能躲避那些比我们强大的人，并保持那我们必须向所有人表现的谦逊。}\par

17. 但我们还必须提防那看得见的，他激怒我们，驱策我们，使我们恼火，并提供那些会引发纵欲或邪情的事物。因此，若有人辱骂我们，激怒我们，煽动我们使用暴力，试图使我们争吵；我们就当保持沉默，不要以闭口不言为耻。因为那激怒我们、伤害我们的人，是在犯罪，并希望我们变得像他一样。\par

18. 确实，你若沉默并隐藏你的感受，他惯于说：\Quote{你为什么沉默？你若敢，就开口；但你不敢，你哑了，我使你哑口无言。} 你若沉默，他越发激动。他认为自己被打败了，被嘲笑，受轻视，被愚弄。你若回答，他认为自己成了得胜者，因为他找到了和自己一样的人。因为倘若你沉默，人们会说：\Quote{那人辱骂，但这人不理会他。} 你若以辱骂还辱骂，他们便会说：\Quote{两人都在辱骂。} 两人都受谴责，无一得免。因此，他的目的就是激怒我，好让我像他一样说话行事。但义人的本分是隐藏自己的感受而一言不发，保存无愧良心的果实，宁信托于善人的判断，也不信托于诽谤者的傲慢，并满足于自己品格的坚稳。因为这就是：\BibleQuote{连善言我也缄口不语。}\EditorialNote{拉丁通行本此处作\Quote{silui a bonis}。}\BibleRef{咏 39:2} 因为那拥有无愧良心的人，不该为虚妄的言语烦扰，也不该重视他人的辱骂过于自己内心的见证。\par

19. 所以，人也要保守自己的谦逊。然而，若他不愿显得过于卑微，便会这样想，并在心中说：\Quote{难道我要任凭这人轻看我，当着我的面说这些反对我的话，好像我在他面前不能开口似的？我为什么不说些什么，好让他难过呢？难道我要让他冤枉我，好像我不是一个人，好像我不能为自己报仇？难道他要控告我，好像我不能收集更恶劣的控告对付他？}\par

谁若如此说话，便不是温和谦卑的，也并非不受诱惑。那诱惑者搅动他，亲自将这样的思想放在他心里。恶神也常多次利用他人，使他向这人说出这样的话；但你要将你的脚立在磐石上。纵然奴隶辱骂，义人也当缄默；即使弱者出言侮辱，他也当缄默；即便穷苦人提出指控，他也无需回答。这些乃是义人的武器，好使他能以退让而得胜，如同那些精于投掷标枪者，惯于以退让取胜，并在逃跑中以更沉重的打击伤及追兵。\par

% structure: p0038 source=h2 target=subsection reason=relative-heading-rank; base=h2

\subsection{第六章}

\Emph{在此事上，我们必须效法\Person{达味}的沉默与谦卑，甚至不显露出应受伤害的样子。}\par

我们听见辱骂时，有何必要烦恼呢？为何我们不效法那说：\BibleQuote{我缄默不言，谦卑自下，甚至绝口不提善言。}的人呢？难道\Person{达味}只说了这话，却未实践吗？不，他也实践了。因为当\Person{革辣}的儿子\Person{史米}辱骂他时，达味默然不语；他虽被武装的人围绕，却没有还口，也不寻求报复；甚至当\Person{责鲁雅}的儿子因想报仇而对他说时，达味也不准许。他犹如哑巴和谦卑的人一样前行；他默默行走；虽被骂作杀人凶手，他亦不惊惶，因为他知道自己温和的性情。所以他不受辱骂的困扰，因为他深知自己的善行。\par

因此，谁若因受屈而立即发怒，即便他想表明自己不该受辱，却使自己看起来该受侮辱。轻看屈辱的人，比为此悲伤的人更好；因为轻看屈辱的人鄙视它们，仿佛不觉其存在；而为此悲伤的人却饱受折磨，正如真切感受到一般。\par

% structure: p0042 source=h2 target=subsection reason=relative-heading-rank; base=h2

\subsection{第七章}

\Emph{\BibleBook{圣咏}{三十九篇{[}三十八篇{]}}如何绝妙地充当了引言的角色。受此圣咏激励，圣人决定撰写\WorkTitle{论义务}。他这样做甚至比\Person{西塞罗}更有理由，西塞罗曾为其子撰写此主题。此外，为何如此。}\par

我并非不加思索地引用了这篇圣咏的开端，而我写信给你们，我的孩子们。因为这篇由先知\Person{达味}交给\Person{耶杜通}咏唱的圣咏，我敦促你们重视，我自己也喜乐其意义的深邃及格言的高超。因为在我们刚简要提及的那些话中，我们学到，这篇圣咏教导了缄默的忍耐和等待适当说话的时机，以及其后诗节中对财富的轻视，这些乃是诸德的主要基础。所以，当我默想这篇圣咏时，便想到要写\Quote{\WorkTitle{论义务}}。\par

虽然有些哲学家曾就此主题撰写著作——例如，希腊作家中的\Person{帕奈提乌斯}及其子，拉丁作家中的\Person{西塞罗}——我却不认为亲自撰写也有违我的职分。正如西塞罗为其子的教育而写，我也同样写信教导你们，我的孩子们。因为我爱你们，我在福音中所生的你们，不下于爱我的亲生儿子。因为本性并不比恩宠更能使我们热烈地爱。我们当然应该爱那些我们相信将永远与我们同在的人，胜过那些只在今世与我们同在的人。后者常生出不肖子孙，给父亲带来羞辱；但你们是我们预先拣选来爱的。他们出于本性，出于必然的爱，这并非足够贴切恒常的导师，以灌输持久的爱。而你们是基于我们刻意的选择而被爱，由此，一种深厚的情感结合着爱的力量：所以，人检验他所爱的，并爱他所选择的。\par

% structure: p0046 source=h2 target=subsection reason=relative-heading-rank; base=h2

\subsection{第八章}

\Emph{\Quote{义务}一词既常被哲学家们使用，亦见于圣经；其来源将在下文说明。}\par

既然，相关人物适合撰论义务，让我们看看主题本身是否同样成立，这词是否仅适用于哲学家的学派，抑或也见于圣经。恰巧，今天在宣读福音时，圣神美妙地为我们提出了一段经文，仿佛敦促我们执笔；由此我们坚定信念，认为\OriginalTerm{义务}{officium}一词也可在我们的语境中使用。因为当司祭\Person{匝加利亚}在圣殿中成了哑巴，不能说话时，经上说：\BibleQuote{他供职的日期一满，就回家去了。}（\OriginalTerm{供职}{officii}）因此，我们读到，\OriginalTerm{义务}{officium}一词可为我们所用。\par

这也并非不合乎理性，因为我们认为\OriginalTerm{义务}{officium}一词源自\OriginalTerm{施行}{efficere}，为求悦耳而变更了一个字母；或至少你应当做那些不伤害任何人，却造福众人的事（\OriginalTerm{伤害}{officiant}）。\par

% structure: p0050 source=h2 target=subsection reason=relative-heading-rank; base=h2

\subsection{第九章}

\Emph{义务应从有德之事、有益之事中拣选，亦须由两者的比较而择一；但在基督徒看来，凡无助于来世生命者，皆不被视为有德或有益。因此，这篇论义务的著作并非多余。}\par

哲学家们认为，义务源自美德与利益，且在这二者之间应择其更优者。继而，他们说，可能发生两种美德或两种利益彼此冲突的情况，问题在于，哪一更为有德，哪一更为有益？因此，首先，\Quote{义务}分为三类：有德的、有益的，以及二者中更优者。然后，这三者又再分为五等；即，二种有德的、二种有益的，以及最后，对于选择它们时正当的判断。他们说，第一种关乎生活的道德尊严与正直；第二种关乎生活的便利、财富、资源、机遇；而正当的判断则必须是选择其中任何一种的基础。哲学家们所言即如此。\par

28. 但我们所度量的，无非是合宜与有德之事，且是按照将来的事物而非现在的事物的准则；我们决不认为任何事物有益，除非它有助于我们获得永生的祝福；绝不是那些仅有助于我们享受现世的事物。我们也不承认机遇和世物的财富有何益处，反认为它们若不放弃便是害处；当我们拥有它们时，应视为负担，而非耗用时的损失。\par

29. 因此，我们的这项工作并非多余，因为我们和他们看待本分的方式完全不同。他们视现世利益为好事，我们却视之为恶事；因为在此世享福的人，如同比喻中的那个富人，在彼世受苦；而在此世忍受苦难的\Person{拉匝禄}，在彼世得到了安慰。\BibleRef{路 16:25} 最后，那些不读他们著作的人，如果愿意，可以读我们的——就是说，如果不需要华丽的辞藻或巧妙的论述，而只满足于主题本身的朴素魅力即可。\par

% structure: p0055 source=h2 target=subsection reason=relative-heading-rank; base=h2

\subsection{第十章}

\Emph{合宜的事在圣经中的出现，远早于哲学家们的著作。\Person{毕达哥拉斯}从\Person{达味}那里借用了沉默的规则。然而，\Person{达味}的规则是最好的，因为我们的首要本分是在说话上要有节制。}\par

30. 我们受教导，\Quote{合宜的事}在我们的圣经中被置于首位（在希腊文中称之为\OriginalTerm{合宜的事}{\GreekText{πρέπον}}）。因为我们读到：\BibleQuote{赞美诗在\Place{熙雍}合宜于祢，天主啊。}在希腊文中这是：\OriginalTerm{赞美诗在\Place{熙雍}合宜于祢，天主啊}{\GreekText{Σοί} \GreekText{πρέπει} \GreekText{ὕμνος} \GreekText{ὁ} \GreekText{θεὸς} \GreekText{ἐν} \GreekText{Σιών}}。而宗徒说：\BibleQuote{你要讲论那合乎健全道理的事。}\BibleRef{铎 2:1} 又在另一处说：\BibleQuote{原来那为万物所属，为万物所本的，既要领众子进入光荣，使救恩的领袖因受苦难而得以成全，本是合宜的。}\BibleRef{希 2:10}\par

31. 难道\Person{帕奈提奥斯}或\Person{亚里士多德}——他们也写过论本分的著作——比\Person{达味}更早吗？其实，生活在\Person{苏格拉底}之前的\Person{毕达哥拉斯}本人，就是追随先知\Person{达味}的足迹，给他的门徒定下了沉默的规则。他甚至要求门徒五年内不使用言语。而\Person{达味}定下他的规则，不是为了损害自然赋予的恩赐，而是为了教导我们谨慎自己的言语。\Person{毕达哥拉斯}制定他的规则，是为了通过不说话来教导人如何说话。但\Person{达味}制定他的规则，是为了通过说话让我们学习如何更好地说话。没有练习，何来教导？没有实践，何来进步？\par

32. 一个人若想接受战斗训练，就每日用武器操练自己。他好似随时准备行动，演练自己在战斗中的角色，挺身而立，仿佛敌人就在面前。或者，为了获得投掷标枪的技巧和力量，他试验自己的兵器，或闪避敌人的攻击，藉着警醒的注意躲开它们。一个想要在海上行船或划桨的人，先在河上练习。那些希望获得悦耳的歌唱风格和优美嗓音的人，先从歌唱中逐渐发出声音开始。而那些寻求藉由强壮的身体和在正规的角力比赛中赢得胜利桂冠的人，则通过每日在角力学校的操练来锤炼四肢，培养耐力，习惯艰苦的劳作。\par

33. 自然本身在婴儿身上教导我们这一点。因为他们先练习发音，然后才学会说话。因此，这些语音是一种练习，也是声音的训练场。所以，那些想要学会在言语上谨慎的人，不要拒绝合乎本性的事，而要凡事谨守；正如瞭望塔上的人通过警醒守望保持警觉，而不是通过睡觉。因为一切事物都因适当且适合自身的操练而变得更加完美和坚固。\par

34. 因此，\Person{达味}并非总是沉默，而只是有时沉默；他不是永远不对任何人说话，而是不回答那激怒他的敌人，那惹他恼怒的罪人。正如他在别处所说：\BibleQuote{他如同聋子，不听那些说虚幻和谋算欺骗的人；如同哑巴，对他们不开口。}又在另一处说：\BibleQuote{不要照愚人的愚昧回答他，免得你也像他一样。}\BibleRef{箴 26:4}\par

35. 所以，首要的本分就是在我们的言语上有节制。这样，赞颂之祭便献于天主；这样，在诵读圣经时便显出虔敬的敬畏；这样，父母便受到孝敬。我很清楚，许多人说话是因为他们不知道如何保持沉默。但通常，当说话于人无益时，保持沉默才是明智的。一个有智慧的人，在打算说话时，会先仔细思量他要说什么，向谁说，在哪里说，以及何时说。因此，在保持沉默和说话上都有一种节制；在我们的行为上也有一种节制。能持守正确的本分标准是荣耀的事。\par

% structure: p0063 source=h2 target=subsection reason=relative-heading-rank; base=h2

\subsection{第十一章}

\Emph{圣经的见证证明，一切本分要么是\Quote{普通的}，要么是\Quote{完全的}。此外还加上称赞慈悲的话，并劝勉人实行慈悲。}\par

36. 一切本分要么是\Quote{普通的}，要么是\Quote{完全的}，我们也可以凭圣经的权威证实这一点。因为我们在福音中读到主说：\BibleQuote{你若愿意进入生命，就该遵守诫命。他问：哪些？耶稣对他说：不可杀人，不可奸淫，不可偷盗，不可作假见证，应孝敬父母，应爱你的近人如你自己。}这些是普通的本分，尚有欠缺。\par

37. 那时，那少年人对祂说：\BibleQuote{这一切我自幼就都遵守了，还缺少什么？耶稣对他说：你若愿意是成全的，去！变卖你所有的，施舍给穷人，你必有宝藏在天上；然后来跟随我。}\BibleRef{玛 19:20} 更早的时候，经上同样记载，主说我们必须爱我们的仇敌，为那诬告和迫害我们的人祈祷，祝福那诅咒我们的人。\BibleRef{玛 5:44} 我们必须这样做，如果我们愿意成全，如同我们天上的父一样；祂使太阳光照恶人，也光照善人，降雨给义人，也给不义的人，不分彼此。\BibleRef{玛 5:45} 因此，这就是完美的职责（希腊人称之为\OriginalTerm{成全之功}{\GreekText{κατόρθωμα}}），借此，一切可能有缺陷的事物都被纠正。\par

38. 同样，慈悲也是一件美事，因为它使人成全，在于它效法成全的父。没有什么比慈悲更能为基督徒的灵魂增辉；慈悲尤其应向穷人显示，这样你便可将他们视为与你共享自然出产的人，因为大地出产的果实是为所有人使用的。因此，你可以自由地将你所有的给予穷人，以这种方式帮助你的弟兄和同伴。你施舍银钱，他却获得生命。你给出金钱，他视之为他的财富。你的钱币成了他的全部财产。\par

39. 再者，穷人给予你的，比你给他的更多，因为他是你救恩上的负债者。你若给裸者衣穿，便是给自己披上义德；你若收留旅客，扶助穷乏，他便为你赢得圣人们的友谊和永恒的居所。这绝非微小的报酬。你播种属地的，却收获属天的。你惊奇天主在圣约伯身上的判断吗？倒不如惊奇他的德行，因为他能说：\BibleQuote{我作了盲人的眼，跛者的脚。我成了穷人的父亲。我的羊羔的毛皮温暖了他们的肩头。客旅从不露宿我的门外，我的门向来为一切行人敞开。}\BibleRef{约 29:15} 显然，凡从未让穷人空手离开其家宅的人，是有福的。再者，没有谁比那体察穷人的需求、并体察弱小无助者的困苦的人更有福了。在审判之日，他将从主那里获得救恩，主因他所施的慈悲而成为他的负债者。\par

% structure: p0069 source=h2 target=subsection reason=relative-heading-rank; base=h2

\subsection{第十二章}

\Emph{为使任何人在施行慈悲时不受阻碍，他指出天主眷顾人的行为；并以约伯的实例证明，一切恶人于其财富的丰盛中仍是不幸的。}\par

40. 然而，许多人因下列缘故而滞于施行主动慈悲的职责，因为他们认为天主不顾念人的行为，或认为祂不知晓我们在暗中所做的、以及我们良心所注视的。有些人又认为祂的审判似乎绝不公正；因为他们看见罪人拥有丰富的财富，享受尊荣、健康与子女；反之，义人却生活在贫穷与卑贱中，他们没有子女，身体患病，且常处于忧伤之中。\par

41. 这并非小事。因为约伯的那三位身为王者的朋友，断言他是个罪人，因他们看见他先是富有，后来变得贫穷；先前儿女众多，后来尽失所有；如今满身毒疮，遍体鞭痕，从头到脚尽是伤口。但圣约伯向他们宣告说：\BibleQuote{如果我因我的罪孽受此苦难，恶人为何得以生存？他们享长寿，积财富，子孙满堂，随心所欲，儿女绕膝，家宅兴旺；但他们无所畏惧，主的鞭笞也不临于他们。}\BibleRef{约 21:7}\par

42. 心志薄弱的人，见此状况，心神不安，便将注意力转离。圣约伯将要以这种人的口吻发言时，开始这样说：\BibleQuote{请宽容我，我也要说话；然后你们可以嘲笑我。我若受指责，无非是以人的身份受指责。因此，请容忍我话语的重担。} 因为我要说的（他的意思是）并非我所赞同的；但我将说出错谬之言，为驳倒你们。或者换一种译法：\BibleQuote{这是怎么一回事？我竟被一个人指责吗？} 意思是：不能因我犯了罪而指责我，尽管我应该受指责；因为你们不是因明显的罪过指责我，而是凭我遭遇的不幸来衡量我应受的惩罚。因此，心志薄弱的人，看见恶人顺利兴旺，而自己却被不幸压垮，便向主说：\BibleQuote{离开我罢，我不愿知晓你的道路。}\BibleRef{约 21:14} \BibleQuote{我们事奉祂有什么益处？奔向祂有什么好处？一切美物都在恶人手中，祂却看不见他们的作为。}\par

43. \Person{柏拉图}备受称赞，因为在他的书\WorkTitle{论国家}中，他让那位扮演反对正义的反对者角色的人为自己的话请求原谅，那些话连他自己都不赞同；并且他说，那个角色只是为了发现真理和探究问题的需要而假定的。\Person{西塞罗}对此甚为赞同，以致他在自己所写的\WorkTitle{论共和国}一书中，也认为必须说些反对那观点的话。\par

44. \Person{约伯}比这些早许多年！他是第一个发现这事，并考虑应当为此提出什么理由，并不是为了修饰他的口才，而是为了寻求真理。他立刻说明，恶人的灯将熄灭，他们的毁灭将来到；\BibleRef{约 21:17} 那教导智慧和训诲的天主不受欺瞒，而是真理的判官。因此，个人的真福不能凭他们所知的财富多寡来估量，而应依据他们内在良心的声音。因为良心作为赏罚的真实而不受贿的判官，判定无辜者与有罪者的功过。无辜者离世时，具有自己纯朴的力量，充分拥有自己的意志；灵魂仿佛充满骨髓。\BibleRef{约 21:24} 但罪人虽在世享尽丰盛，生活奢靡，馨香环绕，却终其一生处于灵魂的痛苦中，结束自己的末日，毫不带走他曾经享受的那些美好事物---只带走他邪恶的代价。\par

45. 思想此事时，你能否认天主审判会予以报应吗？前者内心感到幸福，后者则悲惨；前者自判无罪，后者则是罪人；前者乐于离开世界，后者却为之悲伤。在自身良心中不自以为无辜的人，谁能被宣称为无罪呢？他说：\BibleQuote{请告诉我，他的帐幕有何遮盖？他的记号必不再见。}\BibleRef{约 21:28} 罪人的生活犹如梦幻。他睁开了眼睛。他的安息已逝，他的享乐已逃。岂止如此，恶人即使在世时的安息也只是假象，如今已在地狱中，因为他们活着时便已落入地狱。\par

46. 你看到罪人的享乐；但审问他的良心吧。他难道不比任何坟墓更污秽？你看见他的喜乐，你羡慕他子女身体康健、财富丰裕；然而往内观看他灵魂的疮痏、他内心的哀伤。对于他的财富，我还能说什么呢？当你读到：\BibleQuote{因为人的生命并不在于他拥有财物的丰裕}\BibleRef{路 12:15}，而你又知道，尽管他在你看来富有，他却自认贫穷，并以自身驳斥了你的判断。至于他子女众多且免于疾苦，我还能说什么呢---他满怀悲伤，决定不要继承人，也不希望效法他行事的人接续他。因为罪人确实没有留下继承人。因此，恶人是自己的惩罚，而义人是自己的恩宠---无论善恶，各自行为的报应都由他们自身承受。\par

% structure: p0078 source=h2 target=subsection reason=relative-heading-rank; base=h2

\subsection{第十三章}

\Emph{哲学家们那些否认天主眷顾整个世界或其任何部分的想法被驳斥了。}\par

47. 但让我们回到正题，免得我们看起来忽略了之前为回应某些人的想法而打断的部分。这些人看见有些恶人富有、喜乐、备受尊荣、有权有势，而许多义人却贫穷困顿---因此设想天主或者毫不关心我们（这是\OriginalTerm{伊壁鸠鲁学派}{Epicureans}所说的），或者如恶人所说，不知道人的行为---又或者，祂虽知晓万物，却是不公义的判官，竟允许善人匮乏而恶人富足。然而，为回应这类想法，并将它们与那些被他们视为幸福之人---他们却自认为悲惨---的感受相对照，作一番离题之论并非不合宜。我想他们宁愿相信自己，也不愿相信我们。\par

48. 在此离题之后，我认为驳斥其余各点并非难事---首先是那些以为天主丝毫不关心世界的人的说法。例如，\Person{亚里士多德}宣称，\OriginalTerm{天主眷顾}{providence}仅延伸到月亮。然而，哪个匠人会不关心自己的作品呢？谁愿意舍弃并抛弃那自己认为所造之物？若统治世界有损于祂的尊威，那创造岂不更有损吗？虽然不创造任何事物并无不义，但不关心自己所创造者，则是极端的残忍。\par

49. 然而，如果有人否认天主是造物主，从而将自己等同于野兽和无理性的受造物，那么对于那些甘陷如此不尊严境地的人，我们还有什么可说呢？他们自己宣称天主充满万有，万物都依赖祂的权能，祂的德能和尊威渗透一切元素---大地、上天和海洋；但他们却认为，进入人的灵---这是祂赐予我们最高贵之物---并带着天主尊威的全知居于其中，于祂乃是一种贬损。\par

50. 然而，那些被视为理性哲学家的，却嘲笑这学说的提出者，视其为愚昧放纵。但对于\Person{亚里士多德}的观点，我该说什么呢？他认为天主满足于自身狭窄的界限，生活在自己王国的限定范围内。然而，这也是诗人们传说所述；因为据他们说，世界被三位神明瓜分，一位统管苍穹，一位掌管海洋，一位治理下域；他们还必须留意，不要彼此挑起战争，切莫让对他人领域的思虑牵挂占据心头。同样，\Person{亚里士多德}也宣称，天主对于大地既不关心，对于海洋和下域也是如此。这些哲学家怎会将那些他们追随其足迹的诗人排除在他们的行列之外呢？\par

% structure: p0084 source=h2 target=subsection reason=relative-heading-rank; base=h2

\subsection{第十四章}

\Emph{没有什么能逃过天主的知晓。圣经的见证和太阳的比喻证明了这一点：太阳虽为受造物，却以它的光和热进入万物。}\par

51. 接下来的问题是：天主既对祂的工程显示了\OriginalTerm{天主眷顾}{providence}，是否现在对工程失去了知晓？经上这样写：\BibleQuote{栽植耳朵的，难道自己听不见？制造眼睛的，难道自己看不见？}\par

52. 这种错误的想法，圣先知们并非不知道。\Person{达味}自己就曾引介那些被骄傲充满并占为己有的人说话。因为还有什么比这更显骄傲呢：就是那些生活在罪恶中的人，竟认为其他罪人不该存活，并说：\Quote{主，恶人嚣张，要到几时？恶人嚣张，要到几时？}后又有人说：\Quote{他们说：\InnerQuote{主必看不见，\Person{雅各伯}的天主决不管。}}对此先知回答说：\Quote{民间的愚昧者！你们该当知悉，糊涂的人！你们几时才能明白？装置耳朵的，难道自己听不见？制造眼睛的，难道自己看不见？训戒万民的，难道不惩罚？——教导人知识的，难道不辨是非？主认透人的思念，无非虚幻。}那位洞察一切虚幻的，难道不晓得什么是圣善的，难道不认识祂自己所造的吗？工匠岂能不认识自己的作品？这个人是人，尚且能识透他作品中隐藏的事；天主——祂岂不认得自己的作品？那么，作品岂能比作者更有深度？难道祂造了什么比祂自己更高超的事物，以致作为造作者，却不认识它的价值，作为掌管者，却不晓得它的状况？对于这些人，就说这些吧。\par

53. 但我们以那位的见证为满足，祂说：\BibleQuote{我究察人心与肺腑。}\BibleRef{耶 17:10} 在福音中，主耶稣也说：\BibleQuote{你们为什么心里思念恶事呢？因为祂知道他们正在思念恶事。}\BibleRef{玛 9:4} 圣史也为此作证说：\BibleQuote{因为耶稣知道他们的心意。}\BibleRef{路 6:8}\par

54. 如果我们观看这些人的行为，这种想法就不会太困扰我们。他们不愿让那不受任何蒙蔽者作他们的审判者；他们不愿承认祂知晓隐秘的事，因为他们害怕自己的隐密事被揭发。然而主，\BibleQuote{知道他们的作为，把他们交付于黑暗；他夜间行动如盗贼；奸夫的眼注视黑暗，说：\InnerQuote{没有眼睛看见我。}他蒙住了自己的脸。}\BibleRef{约 24:14} 因为凡躲避光明的，都喜爱黑暗，设法隐藏自己，虽然他不能躲过天主，天主不但知道人所做的事，连人心中所想的，无论在宇宙深处还是在人心中，都知晓。同样，\BibleBook{德训}{篇}的作者也说：\BibleQuote{谁看见我？黑暗包围了我，墙壁遮蔽了我；我还怕谁？}\BibleRef{德 23:18} 然而，他即使躺在床上这样想，也会在他从未想到的地方被拿获。\BibleQuote{他必将蒙羞，}经上说，\BibleQuote{因为他不知道什么是敬畏主。}\BibleRef{德 23:31}\par

55. 还有什么比这更愚昧的呢：认为有什么事能逃过天主的眼目，而那供予光明的太阳，甚至能透入隐蔽之处，其热力的强度能直达房子的地基和内室？谁能否认，被冬寒冰封的大地深处，会被春天的温和所暖化？的确，树木的心核都能感受到冷热之力，以致其根部或被严寒冻伤，或因太阳的温暖而萌发。总之，凡上天的温和向大地展露笑颜之处，大地便盛产各种果实。\par

56. 那么，如果太阳的光线普照全地，射入隐秘之处；如果铁栏杆或沉重门户的阻碍都不能阻挡它们进入内室，那么，充满生命的天主的光荣，又怎能无法进入祂自己所造的人的思想和心灵呢？又怎能看不见祂自己所造的呢？难道祂造了自己的作品，竟然比祂自己更好、更有能力，（倘若如此）以至于这些作品可以随意躲避造物主的鉴察吗？难道祂在我们的心灵中赋予了如此完美和能力，以至于当祂愿意时，自己也无法明了吗？\par

% structure: p0092 source=h2 target=subsection reason=relative-heading-rank; base=h2

\subsection{第十五章}

\Emph{那些对善人遭遇恶事、恶人享受善事这一事实不满的人，可通过\Person{拉匝禄}的比喻，并凭\Person{保禄}的权威，知晓惩罚与赏报乃保留给来世。}\par

57. 我们已经充分讨论了两个问题；而且我们以为，这番讨论对我们来说并非毫无益处。还有第三个问题尚待解决：为什么罪人拥有大量的财富和资产，生活奢华，无忧无虑；而义人却贫穷匮乏，并因丧失妻儿而受罚呢？对于有这类疑问的人，福音中的那个比喻应当使他们满足：因为那富人穿著紫红袍和细麻布衣服，天天奢华地宴乐；而那乞丐，满身疮痍，却只能捡他桌上掉下的碎屑充饥。然而两人死后，那乞丐在\Person{亚巴郎}怀中得享安息；那富人却在苦刑中。由此岂不清楚显明，赏罚是按功过留待死后的吗？\par

58. 这无疑是十分公道的。因为竞赛需要付出大量劳苦——竞赛过后，胜利归于某些人，耻辱归于另一些人。赛程未完成之前，难道就会颁发棕榈枝或授予冠冕吗？\Person{保禄}写得好；他说：\BibleQuote{这场好仗，我已打完；这场赛跑，我已跑到终点；这信仰，我已保持了。从今以后，正义的冠冕已为我预备下了，就是主，正义的审判者，到那一日必要赏给我的；不但赏给我，也赏给一切爱慕他显现的人。}\BibleRef{弟后 4:7} \BibleQuote{在那一日，}他说，主将赏给——不在此世。他在这里奋斗，在劳苦、危险、翻船中，如一位优秀的角力士；因为他知道\BibleQuote{我们必须经过许多的艰难，才能进入天主的国。}\BibleRef{宗 14:22} 因此，除非按规矩奋斗，无人能领受赏报；若非竞赛艰苦，胜利也不荣耀。\par

% structure: p0096 source=h2 target=subsection reason=relative-heading-rank; base=h2

\subsection{第十六章}

\Emph{为证实上文关于赏罚的论述，他又说，若有些人在将来没有为他们保留的赏报，这并不奇怪；因为他们在此世既未劳苦，也未奋斗。他继续说，正是由于这个缘故，现世的财物才赐给这些人，好使他们毫无借口可寻。}\par

59. 难道在竞赛结束之前就给予赏报的人不是不义的吗？因此，主在福音中说：\BibleQuote{神贫的人是有福的，因为天国是他们的。}\BibleRef{玛 5:3} 祂没有说：\Quote{富足的人是有福的，}却是说：\Quote{贫穷的人。}按天主的判断，真福正始于人类痛苦所由来的地方。\BibleQuote{饥渴的人是有福的，因为他们要得饱饫；哀恸的人是有福的，因为他们要受安慰；慈悲的人是有福的，因为天主会怜悯他们；心里洁净的人是有福的，因为他们要看见天主；为义而受迫害的人是有福的，因为天国是他们的；当人因义而辱骂你们、迫害你们，并捏造种种恶言毁谤你们时，你们是有福的。你们欢喜踊跃吧！因为你们在天上的赏报是丰厚的。}祂应许要给予的赏报是未来的、而非现在的——在天上，而非在地上。你还期待什么？还该得什么？为什么在得胜之前就这样急着要求冠冕？为什么想要抖落灰尘、得享安息？为什么在赛程结束之前就渴望入席赴宴？众百姓还在观望，竞技者尚在场上，而你——难道已经寻求安逸了吗？\par

60. 或许你会说：为什么恶人欢乐？他们为何生活奢靡？为何不与我一同劳苦？这是因为，那些没有报名争夺冠冕的人，不必经受竞赛的劳苦。那些没有下场进入跑道的人，无须傅油，也不必沾满尘土。因为对那些等候荣耀的人，劳苦就在眼前。那些涂着香膏的观众，惯于旁观，而不参与角逐，也不忍受日晒、炎热、尘土和暴雨。让竞技者对观众说：来，与我们一同角逐吧。观众却要回答说：我们坐在这里，是为裁判你们；而你们若得胜，便会获得冠冕的荣耀，我们却不能。\par

61. 那些醉心于享乐、奢靡、劫掠、利益或尊荣的人，与其说是竞技者，不如说是观众。他们享有劳苦的收益，却没有德行的果实。他们贪图安逸；靠诡诈和邪恶聚敛财富，但终须偿付自己不义的惩罚，尽管来得迟晚。他们的安息将在地狱，你们的安息却在天堂；他们的居所在坟墓，你们的居所在乐园。因此，\Person{约伯}精妙地说，他们在坟墓里警醒，\BibleRef{约 21:32} 因为他们不能享有那将来复活者所享有的宁静安息。\par

62. 因此，不要像小孩子那样理解、说话或思考；也不要像小孩子那样，现在就要求那些属于将来的事物。冠冕属于成全的人。要等待，直到那成全的来到，那时你们就会认识——不是像在镜中猜谜，而是面对面地\BibleRef{格前 13:12}——真理的清晰面容。那时才会显明：为什么那个邪恶、劫掠他人财物的人竟是富足的，为什么另一个人执掌大权，为什么第三者儿女众多，为什么第四位满载尊荣。\par

63. 或许这一切的发生，正是为了请强盗回答：你本是富有的，为何还要攫取他人的财物？贫穷并未压迫你，匮乏并未驱使你。我岂不是使你富有，好让你没有借口吗？同样，对掌权者也可说：你为何不援助寡妇，以及那些遭受冤屈的孤儿？难道你没有能力？不能帮助吗？我造你正是为此，不是要你行不义，而是要你制止不义。难道经上没有为你写道：\BibleQuote{拯救那受冤屈的人吗？}\BibleRef{德 4:9} 难道经上没有为你写道：\BibleQuote{从罪人的手中拯救穷苦无告的人}吗？对那些丰衣足食的人也可说：我赐你儿女与尊荣，赏你身体健康；你为何不遵行我的诫命？我的仆人，我对你做了什么，或者怎样使你难过？难道不是我把儿女赐给你，把尊荣赐给你，把健康赐给你的吗？你为何否认我？为何以为你的行为不达于我？为何领受我的恩赐，却藐视我的诫命？\par

64. 我们从叛徒\Person{犹达斯}的例子中也可得出同样的结论。他曾在十二宗徒中被拣选，掌管钱囊，用以施舍穷人，\BibleRef{若 12:6} 这样他就不能看似因未受尊敬或穷困而背叛了主。因此，主赐给他这职务，好让主自己也能在他身上得以显明为义；他将负有更大的罪咎，不是像因受委屈而被迫犯罪的人，而是像滥用恩宠的人。\par

% structure: p0104 source=h2 target=subsection reason=relative-heading-rank; base=h2

\subsection{第十七章}

\Emph{接着提出青年时期的义务，以及适合该年龄的榜样。}\par

65. 既然已经充分阐明，恶行必受惩罚，德行必获赏报，那么让我们继续讨论那些自青年时期就应铭记在心的义务，好使这些义务随着年岁增长而日益加深。一个好青年应当敬畏天主，服从父母，尊敬长者，持守纯洁；他不应轻视谦卑，却要喜爱坚忍与端庄。所有这些，都是青年岁月的装饰。正如庄重是老年人的真正优雅，热忱是青年人的特征，同样，端庄——好似出于天性的一种恩赐——在青年身上也显得格外合宜。\par

66. \Person{依撒格}敬畏主，这在他——\Person{亚巴郎}的儿子——身上原是理所当然的；他顺服父亲，甚至为遵从父命而不惜丧失性命。\BibleRef{创 22:9} \Person{若瑟}也一样，虽然他梦见太阳、月亮和星辰向他下拜，却仍甘心服从父命。\BibleRef{创 37:9} 他极为贞洁，若非纯洁的话，连听都不愿听；谦卑至甘愿做奴隶的工作；端庄至宁可逃避；忍耐至甘愿坐监；宽恕过犯甚至以善报恶。他的端庄如此之大，以致被妇人抓住时，宁愿在逃跑中将外衣留在她手中，也不愿舍弃端庄。\BibleRef{创 39:12} 同样，蒙主拣选向民众宣讲天主圣言的\Person{梅瑟}\BibleRef{出 4:10}和\Person{耶肋米亚}\BibleRef{耶 1:6}，也出于端庄而推辞他们凭恩宠原可胜任的事。\par

% structure: p0108 source=h2 target=subsection reason=relative-heading-rank; base=h2

\subsection{第十八章}

\Emph{论端庄的不同功能。端庄应如何规范言谈与静默，陪伴贞洁，使祷告蒙天主悦纳，约束肢体动作；就最后一点，他以极为得体的言辞提及两位圣职人员。他进而指出，行走的姿态也须与这同一德行相称，并且务必谨慎，不让任何不端之言出口，也勿使任何不雅之状显于身体。所有论点均以十分贴切的例证加以说明。}\par

67. 端庄\OriginalTerm{端庄}{modesty}之德何等可爱，其恩宠何等甘饴！它不仅见于行动，更表现于言辞上，使我们言语不致失度，言词不致发出不得体的声音。心灵的明镜往往在言词中映照其形像。节制甚至衡量我们声音的高低，以免过高的声音冒犯他人之耳。即使在歌唱中，首要规则也是端庄，各类言语也是如此，好让人能在端庄原则美化其进程时，逐渐学会赞颂天主，吟唱诗歌，乃至一般言谈。\par

68. 再者，沉默——一切德行安息于其中的沉默——是端庄的首要表现。只不过，它若被视为幼稚或骄傲的标记，便被认为是一种耻辱；若被视为端庄的标记，则被当作一种赞美。\Person{苏撒纳}在危难中保持沉默，认为丧失端庄比丧失性命更糟。她不认为应以贞洁为代价来保全己身。她只向天主说话，因为在天主前她能以真正的端庄倾吐衷曲。她避免观看男人的面容。因为眼神中亦有端庄，使女人不愿注视男人，也不愿被他们看到。\par

69. 无人能认为这赞美仅归于贞洁。因为端庄是纯洁的伴侣，与纯洁相伴，贞洁本身便更为安全。再者，羞愧作为贞洁的同伴和指引亦为美善，因为它不容纯洁在接近危险边缘时受玷污。正是这端庄，在圣经读者面前，从圣母蒙主识认之初，便举荐了她，且如可信之见证，宣示她堪当承受如此重任。因为当她在室中独处，天使向她问安时，她默默不语，因天使的进入而惊惶，童贞女因异见男性形貌而面容不安。故此，她虽谦卑，却并非因此，而是出于端庄，才没有回覆他的问候，也没有给他任何答覆，除了在得知自己将怀孕主之后，询问此事如何成就。她确非仅为答话而发言。\par

70. 甚至在我们的祈祷中，端庄也最蒙悦纳，并为我们从天主那里赢得丰厚恩宠。难道不是这端庄举扬了那税吏，使他连举目望天也不敢，却因此蒙主称赞吗？\BibleRef{路 18:13} 主如此判定他比法利塞人更成义\OriginalTerm{成义}{justification}，法利塞人则因过分骄傲而显得丑恶。正如圣\Person{伯多禄}所说：\BibleQuote{因此，让我们以不朽的温柔、安静的精神祈祷，这在主前是宝贵的。}\BibleRef{伯前 3:4} 所以，端庄是尊贵的；它虽舍弃己权，一无所取，一无所求，且在某些方面退缩于自身能力范围内，但在天主眼中却是富有的——在天主眼中，人本无一富足。端庄是富有的，因为它是天主的分。\Person{保禄}也嘱咐，祈祷当以端庄和节制献上。\BibleRef{弟前 2:9} 他愿此为首要，宛如引领后续祈祷一般，好使罪人的祷词不致自夸，却宛如蒙上羞愧的红晕，因在己罪的追忆前向端庄让步，而堪得更丰厚的恩宠。\par

71. 甚至在我们的举止、姿态和步态上，也须谨守端庄。因为心灵的状态往往显于身体的姿态。为此缘故，我们心中隐藏的人（即我们的内在自我）或被判定为轻浮、夸耀或暴躁，反观则被看作沉稳、坚定、纯洁和可靠。因此，身体的运动可说是一种灵魂的声音。\par

72. 孩子们，你们还记得，我们的一位朋友，虽表面以恪尽职守见称，我却未接纳他进入圣职人员之列，因为他的举止过于失当。我还曾命一位已在圣职班中的人，绝不可走在我前面，因为他步态中的自大实在令我痛心。那是在他犯过一次错误后，重回岗位时我所说的话。我只不容许这一件事，而我的判断并未欺骗我。因为两人都离开了教会。他们的步态显露了他们的本质，而他们心中的背信也证实了这点。那一位在\Person{亚略}异端肆虐之时背弃了信仰；另一位则因贪婪金钱，否认属于我们，好避免身受教会的判决。他们的步态中，可分辨出轻浮的姿态，有若游荡的小丑一般。\par

73. 有些人走路时明显模仿演员的姿态，仿佛自己是游行中抬着神像的人，并且有着点头雕像般的动作，以至于每次换步似乎都踩着某种节拍。\par

74. 我也不认为急急忙忙地走路是得体的，除非有某种危险情况或真正的急需。因为我们常见那些匆匆而行的人，气喘吁吁，面容扭曲。但如果并没有这种匆忙的必要，就会招致正当的反感。然而，我并非在谈论那些偶尔因特殊原因而必须快走的人，而是指那些因长期习惯的羁绊，已将匆忙视为第二天性的人。对于前者，我不能认可他们那缓慢而故作庄严的步履，那会让人想起幽灵的形态。对于后者，我也不喜欢他们那横冲直撞的速度，因为那会使人联想到亡命之徒的败亡景象。\par

75. 相宜的步态，是那种外观上带有权威、分量和尊严，并具平静而沉着的气度。但它的特征必须是全无做作和自负，纯朴而自然。任何虚假的东西都不会令人愉悦。让本性来陶冶我们的举动。倘若本性中确实有什么缺点，就当努力改之。而为了不流于造作，改正本身也不可或缺。\par

76. 然而，如果我们对这些事都如此留意，那么更应当何等谨慎，不让任何可耻的话从我们口中说出，因为那会严重地玷污人。玷污人的不是食物，而是对他人的不义诋毁与污秽的言语。这些事是公然可耻的。在我们的职分中，确实连一句不得体的话，或冒犯端庄的话都不应出口。但我们不仅不该说任何不合宜的话，甚至不该竖起耳朵去听这类言语。因此，\Person{若瑟}逃走并丢下衣裳，为要不致听到任何有违自己端庄的话。\BibleRef{创 39:12} 因为喜欢听的人，会怂恿别人继续说下去。\par

77. 完全明了何为污秽，是极可耻的事。倘若偶然看见此类事物，那是何等的恐怖！在别人身上令我们不悦的，怎能在我们自己身上成为可喜的呢？大自然不就是我们的导师吗？她将我们身体的每个部分都完美地塑造，既供应必需之用，又使之形态匀称优美。然而，她把赏心悦目的部分，显眼而敞开地呈现在眼前；其中，首冠全身的头，悦目的体态轮廓，以及面容的外观，尤为突出，同时又随时便于劳作使用。至于那些顺应自然需求的部分，她则部分地将它们藏匿在身体深处，以免显出令人厌恶的样子；部分地，她也教导并促使我们将它们遮盖起来。\par

78. 那么，大自然不就是\OriginalTerm{端庄}{modesty}的导师吗？依循着她的榜样，人类的端庄——我想\Quote{端庄}之名源于辨识什么是合宜的方式——已经遮盖并隐藏了在身体构造中被置于暗处的部分；就像\Person{诺厄}奉命在方舟旁凿出的那扇门\BibleRef{创 6:16}，其中我们看到教会的预象，也看到人身体的预象，因为那扇门是用来排泄食物残渣的。我们本性的创造者就是这样顾念我们的端庄，这样守护我们身体中合宜与美善之处，将不雅的部分置于背后，挪出我们视线之外。关于这一点，宗徒说得好：\BibleQuote{身体上那些似乎比较软弱的肢体，原是必需的；我们认为不太尊贵的肢体，就越发加上尊贵的装饰，我们不雅的肢体，就越发显得端庄。}\BibleRef{格前 12:22} 确实，因着遵循自然的引导，精心的照料增添了身体的优雅。在另一处，我已较充分地探讨过这个主题，并指出：我们不单隐藏起那些原本为隐藏而赐予我们的部分，而且我们认为，说出这些肢体的名称及其用途，也是不雅的。\par

79. 若这些部位偶然暴露在外，端庄便受到了干犯；但若是故意为之，则被视为极端的无耻。因此，\Person{诺厄}的儿子\Person{含}使自己蒙羞；因为他看见父亲赤身露体而嗤笑，但那些为父亲遮盖的儿子却获得了祝福。\BibleRef{创 9:22} 也正因此，在\Place{罗马}以及许多其他城邦，古代习俗规定：成年的儿子不得与父母同浴，女婿也不得与岳父同浴，免得儿女对父母应有的极大敬意被削弱。然而，许多人在浴场中尽可能地将自己遮盖起来，这样，在全身赤裸的情况下，至少那一部分是被遮盖的。\par

80. 在旧法律下，司祭们——正如我们在\WorkTitle{出谷纪}中所读到的——也穿了裤子，这是主吩咐\Person{梅瑟}的：\BibleQuote{你要给他们做亚麻布的裤子，从腰部直到大腿，用以遮盖他们的裸体。亚郎和他的儿子们进入会幕，或走近祭台，在圣所内行礼时，要穿上，免得他们招致罪罚而死。}\BibleRef{出 28:42} 据说我们当中仍有些人遵行此礼，但多数人加以灵性的解释，认为这话是为谨守端庄、保持贞洁而发。\par

% structure: p0124 source=h2 target=subsection reason=relative-heading-rank; base=h2

\subsection{第十九章}

\Emph{说话者应如何表现出合宜的风范？美丽是否给\OriginalTerm{德行}{virtue}增添什么，若有，增添多少？最后，我们应该当心，切勿让人在我们身上看到任何矫饰或柔弱之气？}\par

81. 我在端庄的种种功用上稍作详尽论述，甚感欣悦；因为我是对你们这些，若非在自己身上认出它的好处，至少也晓得失去它的损失的人说话。端庄适宜于一切年龄、身份、时间和场合，而最相称于幼年与少年时期。\par

82. 但在每个年龄阶段，我们都必须留意，我们的一切行为都要合宜相称，使我们人生的过程形成一个和谐而完备的整体。因此，\Person{西塞罗}认为，在合宜之事上，应当遵守某种秩序。他说，这秩序在于美、在于条理，以及在于为行动所作的适当安排。诚如他所说，这难以用言语解释，但可以相当充分地加以理解。\par

83. \Person{西塞罗}为何要引入美，我并不十分理解；尽管他也确实称赞身体的能力。我们当然不把德性置于身体的美之中，尽管另一方面，我们确也承认某种优雅，比如，当\OriginalTerm{端庄}{modesty}习惯以羞赧的红晕覆盖面容，使之更为悦目。正如工匠的材料越合用，往往工作得越好，同样，端庄在身体的秀雅中更为彰明。只是，身体的秀雅不应是假装的；它应当是自然而不加雕琢的，质朴而非刻意的，不是靠昂贵闪亮的衣裳来衬托，而是仅以普通的衣物蔽体。必须留意，凡关乎诚信或必要所需之物，无所缺欠，而不可为着华丽而增添一物。\par

84. 声音也不应是软弱无力的，或微弱，或带着女人气的声调——许多人习惯于用这种声调，似乎是显得重要的样子。声音应保持某种品质、韵律，以及男子气概的活力。因为凡做到最适合于自己品格与性别之事的，那就是达成了生命之美。这是举止的最佳条理，这是适合每种行动的安排。然而，正如我不能赞同一种轻柔微弱的声音，或柔弱无力的身体姿态，同样，我也不能赞同粗野村俗的风格。让我们随从本性。效法本性，便为我们提供了训练的法则，并给予我们德性的范型。\par

% structure: p0130 source=h2 target=subsection reason=relative-heading-rank; base=h2

\subsection{第二十章}

\Emph{若我们要保守我们的端庄，就必须避免与放荡之人结交，也要避免外人的宴席，以及与妇女的交往；我们在家中的余暇，应当用于虔诚而有德性的操行。}\par

85. 端庄确实有其暗礁——不是她自身带来的，我指的是，她时常会撞上的暗礁，比如当我们与放荡之人交往时，他们在戏笑的外表下，向善人投以毒药。而后者，若他们经常赴宴游乐，频频参与戏谑，便会削弱他们那男子汉的庄重。那么，让我们留意，在想要放松心神时，我们不可摧毁一切的和谐，即仿佛一切善工的融合。因为习惯很快就会使本性转向另一方向。\par

86. 为此缘故，我以为你们明智之所行，正与\OriginalTerm{圣职人员}{clerics}的职责相称，尤其与\OriginalTerm{司铎职}{priesthood}的职责相称——就是说，你们回避外人的宴席，然而仍能款待旅人，并且因着你们在此事上的格外谨慎，不给人口实以加责难。外人的宴席占据人的心思，并且很快就滋生对宴乐的贪爱。世俗及其逸乐的闲谈也常常潜涌而入。人无法掩耳不闻；而禁止谈论却又被视为倨傲的表示。人的酒杯，即使非己所愿，也被一次又一次地斟满。的确，宁可在自己家中一次就致歉推辞，胜过常在他人家中频频辞谢。至少，当人清醒地起身时，自己的临在也无需因他人的无礼而受到谴责。\par

87. 年轻的圣职人员无需前往寡妇或童贞女的家中，除非为着一次明确的探访，且即便那时，亦须由年长的圣职人员陪同，即由主教，或者，若事情较为重要，则与司铎们同往。我们为何要给世人留下毁谤的余地呢？那些频繁的拜访有何必要，以致授人以流言的把柄呢？假如那些妇女中有人偶然失足跌倒，又该如何？你为何要因他人的跌倒而承受责难呢？有多少人，甚至是坚强的人，曾被自己的\OriginalTerm{偏情}{passions}引入歧途？又有多少人，虽确实未曾屈从于罪，却已给人留下了怀疑的口实？\par

88. 你为什么不把职务之外的空闲时间，用在圣堂里阅读呢？你为什么不回去看基督呢？你为什么不向祂发言，并听祂的声音呢？我们祈祷时是向祂发言，我们阅读天主的神圣神谕时是听祂的声音。我们与外人的房屋有何相干？有一间屋宇，容纳众人。需要我们的人可以到我们这里来。我们与闲谈和虚构的故事有何相干？我们领受的，是服务基督\OriginalTerm{祭台}{altar}的职务；并没有强制我们去做取悦于人的职责。\par

89. 我们应当谦卑、温和、柔善、庄重、忍耐。我们必须在一切事上遵守中道，以致平静的面容和安详的言语，都能显明我们的生活中毫无邪恶。\par

% structure: p0137 source=h2 target=subsection reason=relative-heading-rank; base=h2

\subsection{第二十一章}

\Emph{我们必须防范忿怒，在它兴起之前；如果它已经兴起，我们就必须遏制并平息它，而如果连这也做不到，至少我们应当保守口舌不出谩骂，好让我们的偏情如同男童的争吵一般。他引述了\Person{阿基塔斯}所说的话，并指出\Person{达味}在这事上，无论在其行为上，或在其著作中，都作了先导。}\par

90. 应防范\OriginalTerm{忿怒}{anger}。然而，若不能避免忿怒，便要约束它，使之不出限度。因为愤慨是引人犯罪的可怕动力。它扰乱心智，到了不给理性留下任何余地的地步。因此，如果可能，首要目标是藉着恒常的练习、对更好事物的向往、坚定不移的决心，使品格的宁静成为我们的本性倾向。然而，由于偏情在很大程度上植根于我们的本性与品格中，以致无法根除和回避，所以如果能够预见的话，就必须用理性加以约束。但若未及预见或以任何方法防备，而心已充满愤慨，我们就必须考虑，如何克服心神的偏情，如何抑制忿怒，使之不再如此充斥于心。可能的话，抗拒忿怒；若是不能，便让一步，因为经上记载说：\BibleQuote{给天主的忿怒留有余地。}\BibleRef{罗 12:19}\par

91. 雅各伯在兄长发怒时，顺从地退让；对黎贝加也是如此。这就是说，他受到忍耐的劝导后，宁愿离开，寄居异地，也不愿激起兄长的忿怒；直到他以为兄长平息了怒气，才返回。\BibleRef{创 27:42} 他因此在天主前获得了极大的\OriginalTerm{恩宠}{grace}。他何等甘愿服侍，何等厚赠礼物，使兄长与自己重修旧好，不再记挂往日的祝福，只记所献的补赎。\par

92. 因此，若是忿怒已经先发，占据你的心灵，升至心头，你不可离弃自己的阵地。你的阵地是忍耐、智慧、理性和平息的义愤。即便对手的固执惹动你，他的邪恶激你发怒，你若不能平静心神，至少应管束舌头。因为经上记载：\BibleQuote{你要管束你的口舌远离邪恶，使你的嘴唇不说诡诈；寻求和平，追随她。}请看圣雅各伯的平安何等伟大！所以，先要安静心神；若不能，就约束口舌；最后，不可不求和解。这些道理，世上的演说家借自我们，写进他们的著作。但最先说这话的人，才真正明白其意义。\par

93. 所以，我们应当避免忿怒，至少抑制它，不使应受的赞美落空，也不增添罪过。平息忿怒并非易事；与完全不动怒相比，其困难程度毫不逊色。后者出于意志的操练，前者源于本性的效果。因此，孩童的争吵并无恶意，愉快多于怨毒。即便孩子很快相争，也易于平息，转瞬之间，便会更亲密地和好。他们不懂欺骗和狡诈。不要轻视这些孩子，因为主说：\BibleQuote{你们若不变成如同小孩一样，你们决不能进天国。}\BibleRef{玛 18:3} 同样，主自己——\OriginalTerm{天主的德能}{Power of God}，身为孩童，受辱骂时不还口，被打击时不还手。\BibleRef{伯前 2:23} 所以，你要以此为念：如同孩童，不记恨、不怀怨，凡事无可指摘。不要计较别人的回报，守住你的阵地，保持内心的纯朴与洁净。不要以怒对怒，以愚对愚。一个过错必引发另一个过错，石头相击，岂能不生火？\par

94. 异教人（他们说话时惯于夸大其辞）常津津乐道塔兰托的哲学家阿尔基塔斯对管家说的话：\Quote{你这可怜的人，我若不发怒，定要重重责罚你。}但从前的达味在愤恨之际，早已收住他握刀的手。不还口岂不远比不报仇伟大？达味的战士准备向纳巴耳复仇，阿彼盖耳却以恳求劝阻。由此可见，我们不仅要接受适时的祈求，也当以此为乐。达味因此满心喜悦，祝福那出面调停的人，因为他被劝阻而未逞复仇之欲。\par

95. 他先前论及仇敌时曾说：\BibleQuote{因为他们将邪恶加在我身上，在忿怒中苦待我。}让我们听他在怒涛澎湃时所说的话：\BibleQuote{谁愿赐我像鸽子一般的翅膀，我就飞走安息。}他们不停地激怒他，他却寻求安宁。\par

96. 他还说过：\BibleQuote{应发怒，却不可犯罪。}这位伦理教师深知，本性应该受合理教导的引导，而非铲除；故他教导伦理，说：\Quote{在有过犯应当发怒的事上发怒。}因为人不可能不被许多卑劣之事激怒；否则，我们或许不被视为有德，反被当作麻木与疏忽。所以，发怒时切勿犯罪；换言之：你若发怒，不要犯罪，却要用理智胜过忿怒。或可说：你若发怒，就当对自己发怒，因为自己被激动，你便不会犯罪。对自己轻易被激动而发怒的人，就不再对别人发怒了。但是，想要证明自己发怒是正义的人，只会越发火上浇油，迅速陷入罪恶。诚如撒罗满所说：\BibleQuote{能克制忿怒的人，胜过攻城略地的勇士。}\BibleRef{箴 16:32} 因为忿怒连勇士也引人迷途。\par

97. 因此，我们应当在理智准备心灵之前，避免惊慌失措。因为忿怒、忧虑或死亡的恐惧，常使灵魂几乎失去生命，遭受突如其来的打击而跌倒。所以，预先思考并操练心灵，反复考虑事理，是件好事。这样，心灵就不会被突然的纷扰所动，而能在理智的轭和缰绳控制下，保持平静。\par

% structure: p0147 source=h2 target=subsection reason=relative-heading-rank; base=h2

\subsection{第22章}

\Emph{论思考与情感，并论普通谈话及讨论中保持言辞体统}\par

98. 心灵的运动有两种：一是\OriginalTerm{思考}{reflection}，一是\OriginalTerm{情感}{passion}。前者关乎思考，后者关乎情感。两者互不混淆，因其性质分明，殊为不同。思考寻求并仿佛研磨出真理；情感则催促并激发我们行动。因此，按本性而言，思考散发宁静与平和；情感则发出做事的冲动。所以，我们应准备让对美善事物的思考进入心中，并使情感服从理智（若我们确愿引导心灵，守护庄重之事），免得对任何事物的欲望排斥理性。宁可让理性检验并辨明什么合乎德行。\par

99. 我们既已说过应当致力于遵守\OriginalTerm{合宜}{seemly}之事，从而明了言行中的恰当尺度，又因为在行为以前，言语中的秩序更为首要；言语可分为两种：其一，用于友好交谈；其二，用于论述和讨论\OriginalTerm{信德}{faith}与正义之事。无论哪种情况，我们都必须小心，不要激动发怒。我们的言语应当温和宁静，充满和蔼与礼貌，毫无侮辱。在亲切的交谈中，不可有固执的争论，因为这样的争论往往只会引发无益的话题，而非提供任何有用的东西。当有不动怒的讨论，不刻薄的风雅，不严厉的劝诫，不得罪的规劝。正如在我们生活的每一行动中，我们都应谨慎，以免任何压倒性的心灵冲动将理性排除（我们总要为思量留有空间），同样，我们也应在言语中遵守这一规则，免得激起怨怒或仇恨，也绝不显露我们贪婪或懈怠的任何迹象。\par

100. 我们的言语应当如此，尤其当我们谈论\WorkTitle{圣经}时。因为除了这上好的谈话主题——它劝勉人警醒，关切良善的教导——还有什么更值得我们反复谈论呢？我们开始时应有一个理由，结束时应保持在适当的限度内。因为冗长乏味的言论只会激起愤怒。然而，当各种谈话通常都能带来额外愉悦时，唯独它却引起反感，这无疑是极其不合宜的！\par

101. 同样，对于诸如\OriginalTerm{信德}{faith}的教导、节制自律的训诲、正义的讨论、勤奋活动的劝勉等主题的处理，我们不可一时之间全盘托出、详尽阐述，而应按适当进程逐步进行，视我们力所能及以及经文段落的内容所允许的范围而定。我们的讲话不应过于冗长，也不可仓促中断，以免前者留下厌恶之感，后者产生草率疏忽。讲论应当朴素简明，清晰显豁，充满庄重与分量；既不应刻意雕琢、过于文雅，也不可落得令人不快、风格粗俗。\par

% structure: p0153 source=h2 target=subsection reason=relative-heading-rank; base=h2

\subsection{第二十三章}

\Emph{玩笑，虽然有时可能相当得体，也应在圣职人员中完全禁止。声音应质朴坦白。}\par

102. 世俗之人给出了许多关于说话方式的进一步规则，我认为我们可略过不提；例如，应如何进行戏谑。因为尽管戏谑有时可能得体而令人愉快，却不适合于圣职生活。我们怎能采纳那些未见于\WorkTitle{圣经}的事物呢？\par

103. 我们还应注意，在讲述故事时，不可改变我们摆在面前的更严格规则的严肃目的。主说：\BibleQuote{祸哉，你们欢笑的人！因为你们将要哭泣。}\BibleRef{路 6:25} 我们难道还在寻找可笑的材料，以致在此欢笑，到彼却要哭泣吗？我认为我们不仅应避免粗俗的笑话，而且应避免一切戏谑，除非或许在某些时候，使我们的谈话亲切愉快并非不合时宜。\par

104. 论到声音，我当然认为它应当质朴清晰。声音悦耳乃天资，非凭努力可得。发音应清晰，满有男子气概的活力，但应避免粗俗土气的音调。还要留意，声音不可带有戏剧腔调，而应忠实地表达它所说词句的内在含义。\par

% structure: p0158 source=h2 target=subsection reason=relative-heading-rank; base=h2

\subsection{第二十四章}

\Emph{在我们的生活行为中，有三件事需要注意。首先，我们的\OriginalTerm{情欲}{passions}应由\OriginalTerm{理性}{reason}控制；其次，我们应在欲望上保持适度的节制；最后，凡事都应在合适的时机和适当的秩序中完成。所有这些品质在旧约时代的圣人身上如此显著地闪耀，显然他们充分具备了人们所说的\OriginalTerm{枢德}{cardinal virtues}。}\par

105. 我想，关于说话的艺术我已说得够多了。现在让我们来思考什么是适合实践生活的。我们注意到，与此主题相关有三件事要留意。第一是\OriginalTerm{情欲}{passion}不应抗拒我们的\OriginalTerm{理性}{reason}。只有这样，我们的职责才能与\OriginalTerm{合宜}{seemly}之事相合。因为情欲若顺从理性，我们便容易在职责中保持合宜。其次，我们应小心，不要因表现出比所从事事务所需更大或更小的热忱，而看似是在大张旗鼓地处理小事，或以极小的用心处理大事。第三，关于我们努力和工作的节制，以及关于行事秩序和时机得当，我认为万事都应坦率直接。\par

106. 但首要的，我可称之为一切的基础，即我们的情欲应当服从理性。第二点和第三点其实是相同的——无论哪种情况都是节制。我们有时也要审视悦目的形态，即所谓的美，并考量尊严。然后便是考量事物的秩序与时机。因此，这便是三点，我们必须看看是否能在某一位圣人身上完美地展示它们。\par

107. 首先是我们的父亲\Person{亚巴郎}，他受塑造并被召叫，以教导未来的世代。当他奉命离开自己的故乡、亲族和父家，尽管被众多亲属关系所束缚牵累，他岂非证明了在他身上，情欲服从于理性吗？有谁不喜爱故乡、亲族和自己家的甜蜜魅力呢？那些甜蜜确曾使他欢喜。但是，对天上命令和永恒赏报的思念更深深地影响了他。他岂未思量过，他不能毫无极大危险地带着妻子同行，因她不惯于艰辛，柔弱难当侮辱，又貌美足以激起放荡男人的贪欲？然而他颇为沉着地决心忍受这一切，而非以推诿之辞逃避。最后，当他进入\Place{埃及}时，他劝妻子说自己是他的妹妹，而非妻子。\par

108. 你看，是何等的偏情在起作用！他害怕妻子的贞洁，为自己的安全担忧，他猜疑埃及人的淫欲，然而履行对天主本分的合理之道却在他内得胜了。因为他想，靠着天主的恩宠，他在任何地方都能平安，但如果冒犯了主，即使在家中也不能无恙。这样，理性征服了偏情，并使它服从自己。\par

109. 当他的侄儿被掳去时，\BibleRef{创 14:14} 他既不为众多君王的军旅所惊吓，也不惊惶，重新投入战斗。胜利后，他拒绝分享战利品，尽管那确是他自己赢得的。还有，当应许他一个儿子时，虽然他想及自己身体的活力已失，如同已死，妻子的不育，以及自己的年迈，他依然相信了天主，尽管这相反自然的规律。\par

110. 请注意，这里一切都汇集在一起了。偏情并非没有，而是受到了节制。这里有一颗在行动中持平的心灵，既不将大事当作小事，也不将小事当作大事。这里有对不同事务的节制，事情有秩序，时机恰当，言语有分寸。他在信德上是首屈一指的，在德行上显著的，在战斗中刚勇的，在胜利时不贪婪，在家时好客，对妻子关心。\par

111. 雅各伯，他的圣洁的孙子，也喜爱在家中度日，远离危险；但他的母亲却愿他居于外方，从而给兄弟的忿怒让位。健全的计议胜过了自然情感。离乡背井，远离父母，然而无论在哪里，在所行的一切事上，他都遵守恰当的尺度，行事合宜，并利用适当的时机。他在家时如此为父母所爱，以致父亲为他顺从之速所感，给了他祝福，母亲则怀着温柔的爱向他倾斜。在兄弟的眼中，他也被视为当受尊敬，当他以为应当将自己的食物让给兄弟时。因为尽管按自然倾向他想要食物，但当兄弟向他索要时，他出于兄弟情谊而让出。他为主人是个忠信的牧羊人，对岳父是个殷勤的女婿；他工作勤勉，饮食有节，赔偿时显著，报恩时慷慨。而且，他如此平息兄弟的忿怒，以致虽然他曾畏惧其仇恨，却获得了他的恩待。\BibleRef{创 33:4}\par

112. 论及若瑟，我该说什么？\EditorialNote{创世纪第39章} 他固然渴望自由，却忍受了奴役的束缚。在奴隶境地他是多么温良，在德行上是多么恒常，在监狱中是多么仁慈！在解梦时他明智，在运用权力时他自制！在丰年，他不谨慎吗？在荒年，他不公正吗？他岂非值得称赞地在一切事上按序而行，并依时利用机会；以他职务的约束性引导，向自己的人民施行正义吗？\par

113. 约伯在顺境与逆境中，都是无可指摘、忍耐、悦乐天主且蒙接纳的。他被痛苦所折磨，却能寻得安慰。\par

114. 达味在战争中也是勇敢的，在逆境时忍耐的，在耶路撒冷和平的，在胜利时仁慈的，犯罪时悔改的，老年时有远见的。他在行动中保持了恰当的尺度，并随而利用时机。他将这一切记在了后世岁月的诗歌中；因此在我看来，他不仅藉其圣咏的甘美，也以其生活，向天主倾泻了一首关于自己功绩的不朽之歌。\par

115. 在这些人物身上，缺少了什么与主要德行相关的本分呢？首先，他们显示了\OriginalTerm{明智}{prudence}，这明智用于寻求真理，并赋予对全备知识的渴望；其次，\OriginalTerm{正义}{justice}，它分配给每人其所应得的，不贪求他人的，并漠视自己的利益，以维护众人的权利；第三，\OriginalTerm{勇德}{fortitude}，在战争中与在家中都显于心志的伟大，并以体力见长；第四，\OriginalTerm{节制}{temperance}，它在我们认为应当言行的万事上，保持正确的方法与秩序。\par

% structure: p0171 source=h2 target=subsection reason=relative-heading-rank; base=h2

\subsection{第二十五章}

\Emph{本书何以不以讨论上述德行为开端的理由已陈明。亦简明指出，古代圣祖们同样具有这些德行。}\par

116. 或许，既然不同类别的本分源自这四种德行，有人可能会说，应当首先对它们加以描述。然而，若一开始就给出本分的定义，继而将其划分为各类，便显得做作。我们避免了做作，而是提出了古代圣祖们的榜样。这些榜样确实使我们在理解上毫无疑问，也让我们在讨论时没有巧辩的余地。因此，让圣祖们的生活，成为我们德行的明镜，而非仅仅是精明巧计的行为集锦。让我们在效法他们时表示敬畏，而不仅仅是讨论他们时的聪敏。\par

117. 明智在圣亚巴郎身上居于首位。因为圣经论及他说：\BibleQuote{亚巴郎信了天主，这算为他的正义。}\BibleRef{创 15:6} 因为不认识天主的人，不能是明智的。又说：\Quote{愚妄的人说：没有天主。} 因为明智的人不会这样说。那不求寻他的创造者，却对石头说：\BibleQuote{你是我的父亲}的，如何是明智的呢？\BibleRef{耶 2:27} 有谁像\OriginalTerm{摩尼教徒}{Manichæan}那样对魔鬼说：\Quote{你是我生命的创造者}呢？\OriginalTerm{亚略}{Arius}如何是明智的呢？他宁可选择一位不完美且低级的创造者，而非一位真实而完美的创造者。\OriginalTerm{玛尔西盎}{Marcion}或\OriginalTerm{欧诺弥}{Eunomius}如何能是明智的呢？他们宁愿有一位恶神，而非善神。不敬畏自己天主的人，又如何能是明智的呢？因为：\Quote{敬畏主是智慧的开端。} 在别处又说：\Quote{智者不离开主的口，而是颂扬他的伟大，亲近他。} 所以当圣经说：\Quote{这算为他的正义}时，这便为他带来了另一种德行的恩宠。\par

118. 我们的前人中最杰出者曾声称，\OriginalTerm{智德}{prudence}在于认识真理。但他们中谁能在这方面超越\Person{亚巴郎}、\Person{达味}或\Person{撒罗满}呢？然后他们接着说，\OriginalTerm{义德}{justice}顾及整个人类团体。因此\Person{达味}说：\BibleQuote{他散财赒济贫穷，他的仁义存留于永远。}义人怜悯，义人借贷。整个世界财富都在智慧和正义者的脚下。义人把他人的东西视如己物，而把自己的东西视为公共财产。正义的人责备自己而不责备他人。因为正义的人不宽待自己，也不容许自己的隐秘行为被隐藏。现在请看\Person{亚巴郎}多么正义！他在老年依预许得了一个儿子，当主要求他献子为祭献时，他竟不认为应当拒绝，尽管那是他唯一的儿子。见：\BibleRef{创 22:3}。\par

119. 请在此留意，这四枢德在同一行为中一齐呈现。相信天主，不将对儿子的爱置于造物主的命令之上，这是智慧的。将所受的归还，这是正义的。以理性约束自然情感，这是勇敢的。父亲牵着祭品，儿子问祭品在何处：父亲的情感备受考验，却未被战胜。儿子又说：\Quote{我的父亲，}这一声再次刺透父亲的心，却丝毫未减弱他对天主的虔敬。第四种德行，即\OriginalTerm{节德}{temperance}，也临在。他作为义人，在自己的虔诚中保持了应有的尺度，并在所须执行的一切中维持了秩序。因此，他按序带来祭献所需之物，点燃火焰，捆绑儿子，拔出刀，依序献祭；从而他堪受酬报，得以留住自己的儿子。\par

120. 还有比圣祖\Person{雅各伯}更大的智慧吗？他面对面看见了天主，并获得了祝福。见：\BibleRef{创 32:29}。 还有比他更高的正义吗？他将自己所得的分给兄弟，并作为礼物献上。见：\BibleRef{创 33:8}。 还有比他更大的\OriginalTerm{勇德}{fortitude}吗？他曾与天主搏斗。见：\BibleRef{创 32:24}。 还有比他更真实的节制吗？他在时间与地点上都表现出如此的节制，宁愿隐藏女儿的耻辱，而不为自己复仇。见：\BibleRef{创 34:5}。 因为他身处于仇敌之中，认为赢取他们的好感，比让他们将仇恨集中在自己身上更好。\par

121. \Person{诺厄}又是多么智慧，他建造了整艘方舟！见：\BibleRef{创 6:14}。 他又是多么正义！因为他独自一人，从全人类中被保存为人类的始祖，成为过去世代的幸存者，并成为未来世代的肇始者；他的出生，与其说是为他自己，不如说是为了世界和宇宙。他战胜洪水，是多么勇敢！他忍受洪水，是多么节制！当他进入方舟后，他度过多时光，是何等有节制！他放出乌鸦和鸽子，又迎接它们归来，他趁机离开方舟，他对这些时机的运用又是何等有节制！\par

% structure: p0179 source=h2 target=subsection reason=relative-heading-rank; base=h2

\subsection{第二十六章}

\Emph{在探究真理方面，哲学家们违背了自己的准则。然而，\Person{梅瑟}却显示出比他们更有智慧。智慧的尊严越大，我们就越必须热切地努力去获得它。自然本性本身催促我们所有人为此而努力。}\par

122. 因此，据说在探究真理时，我们必须遵守合宜之道。我们应当以最大的审慎寻求真实之事。我们不可将虚假当作真理提出，也不可将真理隐藏于黑暗之中，也不可用无益、繁琐或可疑之事充满心灵。还有什么比敬拜人手所造的木质之物更不合宜的呢？还有什么比讨论几何学和天文学（他们所赞同的学科）的问题、测量空间的深度、将天地限制在固定数字的界限内，而忽略救恩的根基、寻求谬误，更能显示这种黑暗呢？\par

123. \Person{梅瑟}虽然精通埃及人的一切智慧，见：\BibleRef{宗 7:22}。 却不赞同那些事，反而认为那种智慧既有害又愚昧。他由此转离，全心渴望寻求天主，从而看见了天主，向他询问，并在他发言时聆听了祂。见：\BibleRef{出 3:4}。 还有谁比他更智慧呢？他蒙天主教导，并以自己工程的威力使埃及人的一切智慧和一切法术的能力化为乌有。他不曾将未知之事当作已知，从而轻率接受。然而，这些哲学家，尽管不认为敬拜偶像并向毫无知晓的偶像求助是违背本性或可耻的，却教导我们刚才所提到的两件事——即合乎本性和德行的两件事——应该避免。\par

124. 智慧的德行越高超，我再说，我们就越应为获得它而努力，好使我们能达至它。并且，为了不让我们产生违背本性、可耻或不适当的观念，我们应当付出两样东西，即时间与关注，为了探究而对事物加以思考。因为人之所以超越其他一切生物，最突出之处莫过于他拥有理性，探求事物的本源，并认为应当寻求他存在之造主。因为我们的生命与死亡都在祂手中；祂以点头治理这个世界。我们知道我们必须向祂交代我们的行为。因为没有什么比相信祂将是我们的审判者更能有助于善度生活：隐秘之事无法隐瞒祂，不端之事触犯祂，善行则使祂喜悦。\par

125. 因此，按照人的本性，所有人心中都存在一种探求真理的渴望，这渴望引导我们向往知识与学问，并使我们对求知充满热情。在这方面超群出众，在人类看来是高尚的事；但只有少数人能达到。他们通过深入思考、仔细斟酌，付出不小的辛劳，才得以达至真福而有德的生活，并在行动中接近其样式。\BibleQuote{因为不是凡向我说\InnerQuote{主啊！主啊！}的人，就能进入天国；而是那承行我在天之父旨意的人。}见：\BibleRef{玛 7:21}。 唉！我不知道那还能带来什么更多的东西。\par

% structure: p0185 source=h2 target=subsection reason=relative-heading-rank; base=h2

\subsection{第二十七章}

\Emph{首要的责任源泉是智德，由此涌出其他三种德行；它们不能分离或被割裂，因为是彼此相连的。}\par

126. 那么，责任的首要源泉是智德。因为有比向造物主献上全心的虔诚与敬畏更大的责任吗？然而，这一源泉也流向他种德性。因为义德若无智德便不能存在，它要求相当程度的智德来判断一件事是否正义或非义。任何一方面的差错都是非常严重的。\Quote{称义人为恶，称恶人为义的，在天主面前是可咒骂的。为何恶人的正义反倒增多呢？}\Person{撒罗满}如此说。另一方面，智德也不能没有义德而存在，因为对天主的虔敬是智慧的开端。我们注意到，这观念在世间的智者中是借来的，而非原本就有，因为虔敬是一切德性的根基。\par

127. 然而，义德的虔敬首先指向天主；其次指向祖国；再次指向父母；最后指向众人。这也是符合本性的指引的。从生命的起始，当悟性初次注入我们时，我们便爱生命为天主的恩赐，我们爱祖国与父母；最后我们爱与我们有来往的同伴。由此而生出真爱，这爱以他人为先而不求己益，其中蕴藏着义德的卓越。\par

128. 所有生物的首要本能是保存自身的安全，防范有害之事，追求有利之物。它们寻求食物和隐蔽之所，以保护自己免受危险、暴风和烈日的侵害——这一切都是智德的标记。其次，我们发现各类不同的生物本性上习惯于群居，起初与同类同种的伙伴，然后也与其他种类。我们看到牛乐于成群，马乐于结队，尤其同类相喜，鹿也与鹿相伴，且常与人相处。对于它们渴望繁衍后代，关于它们的子嗣，甚至它们的爱欲，我又何必多言呢？在这些事上，义德的影子是显而易见的。\par

129. 那么，很明显这些以及其余的德性是彼此关联的。因为勇德，在战争中保卫祖国免受蛮族侵犯，或在家庭中保护弱者，或保护同伴免遭强盗之害，便充满了义德；而知道以何种策略进行防御和给予援助，如何善用时地之机，则是智德与节制所司之事，而节德本身若无智德便不能遵守应有的尺度。认识恰当的时机，并按正义施行回报，这属于义德。在所有这些事上，宽宏大量也是必要的，还需要心灵的坚忍，且往往需要身体的毅力，好使我们能完成所愿之事。\par

% structure: p0191 source=h2 target=subsection reason=relative-heading-rank; base=h2

\subsection{第二十八章}

\Emph{一个团体建立在义德与善意之上。前者的两个部分，复仇和私人占有，不被基督徒所承认。斯多葛派关于公共财产和互相帮助的说法是从圣经中借来的。义德荣耀的伟大，以及什么阻碍了通往义德之路。}\par

130. 那么，义德与人类社会的联系，以及与整个团体的关系。因为维系社会之事可分为两个部分——义德与善意，后者也称为慷慨与慈惠。在我看来，两者之中义德更高，慷慨更令人愉悦。一个给出判断，另一个展示慈惠。\par

131. 然而，哲学家们视为义德本分的事，在我们这里是被排除的。因为他们说，义德的首要表现是：不伤害任何人，除非因受委屈而被迫如此。福音的权威将此撇开。因为圣经愿人子的灵在我们内，祂来是为赐予恩宠，而非带来伤害。\BibleRef{路 9:56}\par

132. 其次，他们认为义德适宜的做法是：将公共财产，即公产视为公共的，将私产视为私人的。但这甚至不符合本性，因为本性已为所有众人倾流万物，以供共同使用。天主使万物生长，为使人人有共同的食物，且使大地成为人人的共有产业。因此，本性为所有人产生了共同的权利，而贪婪却使之变为少数人的权利。这里，我们也得知斯多葛派教导说，地上所生的一切为人类的使用而被造，而人类彼此为对方而生，为能互相裨益。\par

133. 但若非从圣经中，他们能从哪里得到这些观念呢？因为\Person{梅瑟}写道，天主说：\BibleQuote{让我们照我们的肖像，按我们的模样造人，叫他管理海中的鱼、天空的飞鸟、牲畜、各种在地上爬行的生物。} \BibleRef{创 1:26} \Person{达味}也说：\BibleQuote{你将一切置于他的脚下：所有的羊和牛，以及田野的走兽、天空的飞鸟和海中的鱼类。} 因此，这些哲学家从我们的著作中学到万物被置于人下，于是他们认为万物也是为人而受造的。\par

134. 我们在\Person{梅瑟}书中也读到，人是为了人而受造，当上主说：\BibleQuote{人单独不好，我要给他造个与他相称的助手。} \BibleRef{创 2:18} 于是女人被赐予男人，为帮助他。她应为他生育子女，使一个人常能帮助另一个人。再者，在女人受造之前，论到\Person{亚当}曾说：\BibleQuote{没有找着个与他相称的助手。} \BibleRef{创 2:20} 因为一个人除非从另一个人那里，不能获得适当的帮助。所以，在所有生物中，没有与他相称的，或者坦率地说，没有能做他助手的。因此，需要一个女人来帮助他。\par

这样，按照天主的旨意和本性的联系，我们理应彼此相助，在履行职责上相互竞争，将我们的一切利益仿佛摆在众人面前，（用圣经的话说）出于虔诚或本分之心互相帮助，或捐钱，或作事，无论如何总以某种方式相助；好使人类友谊的温馨在我们中间日益甘甜，没有人因畏惧危险而退缩其职责，反倒是无论顺逆，都视为自己的事。如此，圣\Person{梅瑟}为了他的百姓，不畏惧从事可怖的战争，不怕最强悍君王的武力，也不为野蛮民族的凶残所吓倒。他将自身的安危置之度外，为要给予百姓自由。\par

如此，正义的荣耀何其大；因为正义与其说是为自己，不如说是为他人而存在，帮助我们中间团结和友谊的纽带。她拥有如此高的地位，以致万物都服在她的权下，更能帮助他人，供给金钱；她不拒绝为人服务，甚至为他人承担危险。\par

若非贪婪削弱并减损这德能的力量，谁不乐意攀登并持守这德行的巅峰呢？因为只要我们想要增加财产，聚敛金钱，占据新的土地，并要成为众人中最富有的，我们便已将正义的形式抛弃，丧失了善待众人的福分。那企图从别人夺取自己所欲之物的人，又怎能称得上正义呢？\par

图谋权势的欲望也削弱了正义的完美力量和美丽。因为一个企图使他人臣服于自己权下的人，怎能替他人出面呢？一个以牺牲自由为代价而渴望大权的人，又怎能帮助弱者反抗强者呢？\par

% structure: p0202 source=h2 target=subsection reason=relative-heading-rank; base=h2

\subsection{第29章}

\Emph{正义即使在战争中和对待敌人时也当遵守。\Person{梅瑟}和\Person{厄里叟}的榜样证明了这一点。古代作家从希伯来人那里学会了用更温和的词语称呼他们的敌人。最后，正义的基础在于信德，其和谐之美在教会中臻于完美。}\par

正义何等伟大，可由以下事实推知：没有地方，没有人物，没有时间，与它毫无关联。即使在与敌人交往的一切事务中，亦当保守正义。例如，若与敌人商定了交战的日期或地点，则预先占据该地，或提前发动攻击，便被视为违背正义。因为在战役中被击败，或由于激烈交锋，或由于对手技高一筹，或仅为偶然所致，其间确有区别。但对于更为凶残的敌人，以及诡诈和对我们犯下更大恶行者，则施以更严厉的报复，正如对待米德杨人那样。\EditorialNote{户 31} 因为他们曾藉其女子引诱众多犹太百姓犯罪；为此，主的烈怒倾泻在我们祖先的百姓身上。因此，\Person{梅瑟}获胜后，未曾留下一个活口。另一方面，\Person{若苏厄}并未攻击那些以诡计而非战争试探我们祖先百姓的基贝红人，却以奴役之轭惩罚了他们。\EditorialNote{苏 9} 再者，当以色列王想要击杀\Place{叙利亚}人时，\Person{厄里叟}也不容许。敌人围困他时，他先以瞬间的盲目击打他们，使他们看不见所走的路，遂将他们引入城中。因为他说：\BibleQuote{你不要击杀那些你未用枪用剑掳来的人；要给他们摆上饼和水，让他们吃喝，然后回到自己主人那里去。}因受此恩待，他们应将所受的恩惠昭示于世。\BibleQuote{于是}（我们读到）\BibleQuote{\Place{叙利亚}的部队不再侵犯\Place{以色列}地域。}\par

那么，如果正义即使在战争中也具有约束力，我们在和平时期岂不更当遵守？先知对那些前来捉拿他的人显示了何等的恩惠。我们读到，\Place{叙利亚}王派兵设伏，因他得知自己的计划和谋略都是\Person{厄里叟}泄露的。先知仆人\Person{革哈齐}看见军队，开始为自己的性命担忧。但先知对他说：\BibleQuote{不要害怕！因为在我们这里的比他们在那里的更多。}当先知祈求打开他仆人的眼睛后，他的眼睛就开了。于是\Person{革哈齐}看见整个山上满布火马火车围绕着\Person{厄里叟}。当他们下到先知那里时，先知说：\BibleQuote{天主啊，求你降罚\Place{叙利亚}的军队，使他们失明。}这祈求蒙允后，他对\Place{叙利亚}人说：\BibleQuote{跟我来，我把你们领到你们所寻找的人那里去。}于是他们看见了他们想要捉拿的\Person{厄里叟}，及至看见他，却不能拘禁他。由此可见，即使在战争中亦当持守信德和正义；若信德遭背弃，便不得不为耻辱之事。\par

同样，古人习惯以较温和的名字称呼他们的敌人，称之为外方人。因为在古代习俗中，敌人常被称为外方人。我们也可以说，他们采纳了我们的典籍；因为希伯来人称呼他们的敌人为\OriginalTerm{外族人}{allophyllos}，翻成拉丁文即\OriginalTerm{外族人}{alienigenas}（意为异族）。因为我们在\BibleBook{列王纪}{上}中读到：\BibleQuote{那时，那些异族列阵攻击\Place{以色列}。}\par

142. 因此，正义的基础在于信德，因为义人的心思念信德，而那控告自己的义人，正是将正义建立在信德上；当他宣认真理时，他的正义便显明了。主藉\Person{依撒意亚}说：\BibleQuote{看，我要在\Place{熙雍}放置一块石头作为基础。}\BibleRef{依 28:16}这指的是\Person{基督}作为教会的基础。因为\Person{基督}是所有人信仰的对象；而教会仿佛是正义的外在形式，她是所有人的共同权利。因为教会为所有人共同祈祷，为所有人共同工作，在所有人的诱惑中经受考验。因此，那弃绝自己的人确实是义人，确实配得上\Person{基督}。为此，\Person{保禄}将\Person{基督}立为基础，好使我们在祂上面建造正义的工程\BibleRef{格前 3:11}，而信德则是地基。那么，在我们的行为中，如果是邪恶的，便显出不义；如果是善的，便显出正义。\par

% structure: p0208 source=h2 target=subsection reason=relative-heading-rank; base=h2

\subsection{第三十章}

\Emph{论慈惠及其各部分，即善意与慷慨。如何将二者结合。为使人以可称赞的方式展现慷慨，尚需什么其他条件。}\par

143. 现在我们可以接着谈论慈惠，它分为两部分：善意与慷慨。完美的慈惠必须兼具这两种品质。仅仅心存善念是不够的；我们还得付诸行动。反过来说，仅有行动也不够，除非这行动出自良好的源头，即出自善意。\BibleQuote{因为天主喜爱乐善好施的人。}\BibleRef{格后 9:7}如果我们行事出于勉强，那还有什么赏报呢？因此宗徒概括地说：\BibleQuote{我若甘心做这事，就有赏报；若不甘心，那只不过是将责任托付于我了。}\BibleRef{格前 9:17}在福音中，我们也领受了关于正义之慷慨的许多规诫。\par

144. 因此，善意待人、慷慨施予，唯一愿望是行善而非作恶，这实为光荣之事。因为如果我们认为有责任提供财物给挥霍者让他继续挥霍，或给奸淫者为他的奸淫付钱，那便不是慈惠之举，因为其中缺乏善意。如果我们给一个人钱，协助他密谋反对自己的国家，或试图用我们的钱财纠集一些堕落之徒攻击教会，那么我们是在害人而非助人。再者，帮助那严酷压迫寡妇孤儿的人，或者帮助那试图以任何暴力手段夺取他们财产的人，看起来也非慷慨之举。\par

145. 从一个人那里勒索东西来送给另一个人，或不义地赚钱，然后认为可以好好花用，这并非慷慨之心的表现；除非我们效法\Person{匝凯}\BibleRef{路 19:8}的榜样，将我们所抢夺的以四倍偿还，并以信德的热忱和真正的基督徒努力来弥补这类外教人的罪行。我们的慷慨必须建立在坚实的基础上。\par

146. 施行慈惠的当务之急是要怀着善意行事，奉献时不可弄虚作假。绝不可我们实际做得少，却自称做得多。何必非要言说呢？承诺中隐藏着欺骗。我们有能力给予我们所愿意的。欺骗会毁坏根基，从而破坏工程。\Person{伯多禄}的愤怒难道仅至于希望\Person{阿纳尼雅}和他的妻子被杀吗？\BibleRef{宗 5:11}当然不是。他愿其他人因知悉他们的榜样而不致丧亡。\par

147. 若你施予是为了炫耀，而非出于怜悯，这也不是真正的慷慨之举。你内心的情感为你的行为命名。行为如何从你发出，别人就会如何看待它。看，你有一位多么真实的判官！祂与你商议如何着手你的工作，首先便审问你的心意。\BibleQuote{祂说：\InnerQuote{不要叫左手知道你右手所行的。}}\BibleRef{玛 6:3}这并不是指我们实际的身体，而是说：不要让与你同心志的人，连你的弟兄，也不要知道你所作的，以免你在此世借炫耀寻求酬报，从而失去将来赏报的果实。然而，当一个人默默隐藏他的作为，暗中帮助他人的需要，而赞颂他的是穷人的口，而非他自己的唇舌，这才是真正的慷慨。\par

148. 完美的慷慨由其善意、它所帮助的事件、施予的时机与场合而得到证实。然而，我们首先务要援助那些\OriginalTerm{信仰之家}{household of faith}的人。\BibleRef{迦 6:10}如果一位信徒处于匮乏，而你知道此事，或知道他一无所有、忍饥受饿、遭受困苦，尤其当他耻于自己的穷乏，而你却不伸出援手，这便是重罪。如果他被监禁或因家人的诬告而备受压迫，你不去援助，这也是大罪。如果他身陷囹圄，虽然为人正直，却不得不因债务而受苦受刑（因为虽然我们应对所有人施以怜悯，但更当这样对待正直的人）；如果在患难之时，他从你那里得不到任何帮助；如果在危急关头，当他被拉去处死的时候，你的钱财在你眼中比垂死者的生命更为宝贵，这对你是何等的罪过！因此，\Person{约伯}优美地说道：\BibleQuote{愿那临亡者的祝福临到我身上。}\BibleRef{约 29:13}\par

149. 天主确实不偏待人，因为祂洞悉一切。我们也本当向众人施予怜悯。然而许多人以虚假的理由寻求帮助，并声称自己处境悲惨；因此，在情况清楚、人物为人熟知、且刻不容缓时，更应迅速施予怜悯。因为主并不苛求至极。诚然，舍弃一切而跟随祂的人是蒙福的，但同样蒙福的，是以自己所有、尽己所能施予的人。\BibleQuote{主看重穷寡妇的两个小钱，胜过富人的一切献仪，因为她把一切所有的都投上了，而他们只从丰盈中取了一小部分。}\BibleRef{路 21:3}因此，是那份心意使施予有价值或无价值，并使事物各得其值。主不希望我们将全部财产一次分尽，而是要我们逐步分施；除非我们如同\Person{厄里叟}那样，宰杀了牛，用所有的一切款待了百姓，这样，家室的忧虑才不会牵绊他，使他能舍弃一切，全心投入先知的训导。\par

150. 真正的慷慨还当经受这样的考验：如果我们知道自己的近亲有需要，不可轻视他们。因为与其让他们羞于向他人求助，或在需要时到别人那里乞求，还不如你去帮助他们。然而，这并不表示他们应藉你本可施予穷人之物而致富。我们要考虑的是事实，而非人情。你将自己奉献于主，并非为了使你的家族致富，而是为了藉善工的果实赢得永生，并藉施予怜悯来补赎你的罪过。他们或许以为自己所求甚少，但他们索取的是你为赎罪而应付的代价。他们企图夺去你生命的果实，却自以为行事正当。有人指责你没有让他变得富有，而他始终试图骗取你永生的赏报。\par

151. 至此，我们已提出劝告，现在让我们寻求权威依据。首先，无人当以从富有变为贫穷为耻，若这是因他为穷人慷慨施舍所致；\BibleQuote{因为基督本是富有的，却成了贫穷的，为使你们因着祂的贫穷而成为富有的。}\BibleRef{格后 8:9}祂赐给了我们效法的准则，好使我们能善用所减少的产业来交账；凡使饥饿者得饱饫的人，便是减轻了自己的困苦。\BibleQuote{论到这事，我给出我的意见，}宗徒说，\BibleQuote{因为这对你们是有益的，你们当效法基督。}\BibleRef{格后 8:10}劝告是对善人的指引，但警告则约束作恶者。他又像是对善人说：\BibleQuote{因为你们不但已经开始实行，也在一年前就有了乐意的心。}\BibleRef{格后 8:10}这两者，而非仅此其一，是成全的标志。因此他教导，没有善意的慷慨，或有善意而无慷慨，两者都不成全。所以他也敦促我们迈向成全，说：\BibleQuote{现今，你们当完成这件事；好使你们既有乐意的渴望，也按你们所有的去完成。因为人若有乐意的心，必蒙悦纳，乃是照他所有的，并不是照他所没有的。}\BibleRef{格后 8:11}这并不是要使别人富足，你们自己却匮乏，而是要出于均平：你们的富余现在要补他们的不足，好使他们的富余将来也补你们的不足；这样就均平了，如经上记载：\BibleQuote{多收的没有剩余，少收的也没有不足。}\BibleRef{出 16:18}\par

152. 我们注意宗徒如何将善意与慷慨并论，连同方式、正当施予的果实，以及有关的人。方式确实，因为他劝告了那些未成全的人：因为只有未成全的人才会感到焦虑。但若有任何司铎或其他圣职人员，不愿加重教会的负担，并不将所拥有的一切全部施舍，而是按自己职务的要求行事得体，他在我看来并非不成全。我也认为宗徒在这里所说的焦虑，并非指心灵的忧虑，而是指家室的烦扰。\par

153. 并且我认为他论及有关的人时说：\Quote{使你们的富余补他们的不足，好使他们的富余也补你们的不足。}这句话的意思是，民众在物质上的富余能激励他们行善，以补他人食物之缺；而后者在精神上的富余，则能补民众自身灵性功绩之不足，并因此为他们赢得祝福。\par

154. 为此他给出了一个极好的例子：\Quote{多收的没有剩余，少收的也没有不足。}这个例子极大地鼓励了众人去施予怜悯。因为那拥有许多黄金的人并没有剩余，因为世间的一切都如同虚无；而那少有财物的人也没有不足，因为他所失去的本就是虚无。整个事情毫无损失，因为一切本已失落。\par

155. 我们也可以这样正确地理解它。那拥有很多的人，即使不施舍，也没有剩余。因为无论他获得多少，他总是感到缺乏，因为他渴望更多。而那仅有少许的人却没有不足，因为周济穷人无需太多花费。同样，那以精神祝福回报金钱的穷人，虽然拥有许多恩宠，却没有剩余。因为恩宠不会使心灵负担沉重，反而使其轻省。\par

156. 还可以这样理解：人啊，你没有剩余！因为无论你看来拥有多少，你真正领受的又有多少呢？\BibleQuote{在妇人所生的人中，没有一位比\Person{若翰}更大，然而他在天国里最小的，比他还大。}\BibleRef{玛 11:11}\par

157. 或者再次说明。天主的恩宠，以人的说法，永不会过多，因为它是属灵的。谁能度量它的伟大或广阔？那是人眼不能见的。信德若如一粒芥子，就能移山——而所赐给你的，不会超过一粒。若恩宠完全寓居于你内，你岂不须恐惧，你的心意因如此大恩而高举？因为许多人从属灵的高处跌落，比从未从主领受恩宠者更为惨重。而那拥有少的并不缺乏，因它并非可触摸以作分配；在拥有者看来为少的，在缺乏者看来却是多的。\par

158. 施与时，我们也须考虑年龄和衰弱；有时也要顾念那出于本性的羞惭之心，这羞惭表明良好的出身。对于因年老不能再靠劳动自给食物的人，应多予施舍。同样，身体的孱弱也须帮助，且要爽快地帮助。再者，若有人由富足陷入穷困，我们必须援助，特别是若他并非因己罪丧失所有，而是由于抢劫、流放或诬告。\par

159. 也许有人会说：一个盲人坐在一处，人们经过却漠视他，而一个健壮青年却常获施舍。这话有理；因为那青年以他的纠缠强求感动了人。这并非因人们判定他当得施舍，而是因他的乞求让人厌烦。因为主在福音中论到那个已关了门的人：当有人猛烈敲门时，他虽已关门，却起来给予所需，正是因那人的纠缠强求。\BibleRef{路 11:8}\par

% structure: p0227 source=h2 target=subsection reason=relative-heading-rank; base=h2

\subsection{第三十一章}

\Emph{领受的恩惠应以更慷慨的心报答。大地提供了这方面的例证。引用撒罗满关于宴饮的经文以证明此点，后文将以灵性的意义加以阐释。}\par

160. 同样，那曾向你施恩或赠礼的人，若日后陷入匮乏，你更应多加顾念。因为有什么比受了恩而不报答更与本分相悖呢？我也不认为报答应只求等价，反而应更多。人当因曾从他人那里享受恩情，而如此补偿，以致能援助那人，甚至消解他的急需。因为若在报答中不胜于施恩时，便是位次更低的；那先施与的人，不仅在时间上为先，在表达仁善心意上也是占先。\par

161. 因此，在这方面我们当效法大地的本性：大地惯于将所领受的种子，以千倍的增多归还。所以经上写着：\BibleQuote{愚昧人如田地，无知人如葡萄园。你若舍之，他必荒芜。} 智慧人也如田地，归还所受的种子时更为丰盛，仿佛借出而取利。大地或是自发生出果实，或是将自己受托之物丰饶地偿还并归还。这两种方式都要求你在承受父业之用时有所回报，以免你被弃置如同不结果实的田地。人或许能找借口不施与什么，但岂能推卸不报答所领受的呢？完全不施与已不算正当；不报答所受之恩，则更绝对不当。\par

162. 所以撒罗满说得好：\BibleQuote{当你坐在统治者席前吃喝时，该细心观察摆在面前的，然后伸手取用，并当知你必须作好预备。但你若贪婪无厌，不可贪恋他的珍馐美味，因为那些不过虚妄的生命。} 我写下这些话，愿我们众人都能遵行。服务他人固然是善事，但若不知回报，则甚为冷酷。大地本身便提供了仁爱的榜样：她自发生出你未播种的果实；她也以多倍偿还所领受的。你不该否认已收受的金钱，又岂能忽视所受的服侍而不为所动呢？在\BibleBook{}{箴言}书中也有话说：报答恩情在天主前有如此大能，即便在毁灭之日，人也能借此蒙受恩宠，纵使他的罪恶比其他一切都重。我何需再列举其他例证呢？因为主亲自在\BibleBook{}{福音}中许给圣徒们更丰满的赏报，并劝勉我们行善工，说：\BibleQuote{你们要宽恕，你们也将被宽恕；你们要施与，你们也将得到施与；人将以好的升斗，连按带摇，满溢外流，倒入你们怀中。}\BibleRef{路 6:37}\par

163. 但撒罗满所说的宴饮不仅关乎普通的食物，而是应理解为关乎善工。灵魂岂有比以善工为宴更善的方式？有什么能像知道自己行了一件善工那样，轻易充满义人的心灵呢？有何食物比承行天主的旨意更为甘美？主管告诉我们，祂单单以此食物为丰足，正如\BibleBook{}{福音}中所载说：\BibleQuote{我的食物就是承行我在天之父的旨意。}\BibleRef{若 4:34}\par

164. 我们当以此食物为乐，先知说：\BibleQuote{你当以主为乐。} 他们以此食物为乐，是那以奇妙知识学会了领略更高乐趣的人；谁能明白那纯净且能为理智所领悟的乐趣是什么呢？因此，我们当吃智慧的食粮，以天主的话饱饫己身。因为按天主肖像所造的人，生活并不单靠饼，也靠由天主口中所发的一切言语。\BibleRef{玛 4:4} 论及杯爵，圣约伯也清楚地说：\BibleQuote{犹如大地等候甘雨，他们同样等候我的言语。}\BibleRef{约 29:23}\par

% structure: p0234 source=h2 target=subsection reason=relative-heading-rank; base=h2

\subsection{第三十二章}

\Emph{在说明了应为上述宴筵的服侍作出何种回报之后，列举了报恩的种种理由。接着，他赞扬善意，论及善意的效果及其秩序。}\par

165. 因此，为我们来说，被圣经的劝谕所滋润，并使天主的话如同甘露降在我们身上，是一件好事。所以，当你坐在那位伟大者的餐桌前时，要明白那位伟大者是谁。你被安置在喜乐的天堂，出席智慧的盛筵，要想想摆在你面前的是什么！圣经就是智慧的盛筵，各卷书就是不同的菜肴。首先，要知道这筵席提供了哪些菜肴，然后伸出手去，好使你阅读的，或从上主你的天主领受的，能在行动中实行出来，这样，通过你的职责，能彰显所赋予你的恩宠。\Person{伯多禄}和\Person{保禄}正是如此，他们在宣讲福音时，对那位白白赐给他们一切的主作出了一些回报。因此，他们都能说：\BibleQuote{因天主的恩宠，我成为今日的我；祂赐给我的恩宠没有落空，我比他们众人更劳碌。}\BibleRef{格前 15:10}\par

166. 一个人报答所受的恩惠，会以金还金，以银还银。另一人付出劳力。还有一人——我不知道他是否做得更为充分——仅献上内心的祝福。然而，如果当前没有机会回报呢？如果我们希望回报恩情，更重要的在于我们所怀的精神，而非我们财产的多少，因为人们更看重我们的善意，而非我们充分回报的能力。因为所受的恩惠，是按人所拥有的来衡量的。所以，善意是一件伟大的事。即使它没有什么可给予，却奉献得更多；虽然它自己一无所有，却慷慨地施予许多人，而且这样做，自己毫无损失，却使许多人受益。因此，善意比慷慨本身更美好。它在品德上比后者在礼物上更为富有；因为需要恩惠的人，比富足的人更多。\par

167. 善意也与慷慨相伴，因为慷慨实际上源于善意，因为施予的习惯产生于施予的愿望。然而，它也是独立存在的。因为在缺乏慷慨之处，善意仍存——它仿佛是众人共同的父母，联结并维系着友谊。它在忠告上忠信，在顺境中喜乐，在逆境中忧伤。因此，人们更愿信靠有善意之人的忠告，而非智慧之人的忠告，正如\Person{达味}所做的那样。因为达味虽然更有远见，却听从了较为年轻的\Person{约纳堂}的忠告。若将善意从人世间除去，就如同从世界上移走了太阳。因为没有它，人们就不会再去关心为陌生人指路、引领迷途者归来、款待客人（这不是小德，因为\Person{约伯}曾以此自夸，他说：\BibleQuote{我从来没有拒绝过客旅，我的门户常为行人敞开。}\BibleRef{约 31:32}），甚至不会从家门前的流水给人水喝，或用自己的灯点亮他人的蜡烛。因此，善意存在于这一切之中，如同泉水滋润干渴者，如同光，照亮他人，却不会使那些从自己的光中给人光明的人有所减损。\par

168. 还有一种源于善意的慷慨，它使人撕毁自己持有的债务人的借据，而不要求他偿还任何债务。\Person{圣约伯}以自己的榜样吩咐我们如此行动。因为那富有的，不借贷；但那贫穷的，却不能终结约定。那么，如果你并不需要，为何要为自己贪婪的继承人积蓄那些你可以立即归还的东西，从而因善意而获得称赞，并且毫无金钱损失呢？\par

169. 追根溯源——善意首先始于家人，即子女、父母、兄弟，然后一步步扩展到整个世界。它从乐园起始，充满了世界。因为天主将善意之情放在男人和女人内，说：\BibleQuote{他们应成为一体，}\BibleRef{创 2:24} （可以加上）一心一意。因此，\Person{厄娃}也相信了蛇；因为那领受了善意恩赐的她，不认为会有恶意。\par

% structure: p0241 source=h2 target=subsection reason=relative-heading-rank; base=h2

\subsection{第三十三章}

\Emph{善意尤其存在于教会内，并滋养类似的德行。}\par

170. \OriginalTerm{善意}{good-will}在教会的身体内扩展，通过信仰的共融，圣洗的纽带，藉所领受的恩宠而来的亲缘，以及奥迹的共融。因为所有这些联系都自称为亲密关系、子女的孝爱、父母的权威与宗教关怀、兄弟的关系。因此，恩宠的联系清楚地指出善意的增长。\par

171. 对达到相似德行的渴望也大有裨益；同样，善意也会产生性格上的相似。因为君王之子\Person{约纳堂}因爱\Person{圣达味}，效法了他的温柔。所以，那些话：\Quote{你对圣者，必为圣，}看来不仅涉及我们日常的交往，也与善意有关。\Person{诺厄}的儿子们固然同住一起，但他们的性格却毫不相似。\Person{厄撒乌}和\Person{雅各伯}也在父家同住，但彼此非常不同。然而，他们之间没有善意，使一方看重对方胜过自己，反而争夺谁能先获得祝福。因为一方如此刚硬，另一方温柔，在如此不同的性格和冲突的欲望之间，善意无法存在。此外，\Person{圣雅各伯}不可能看重父家中不配为子的，胜过德行。\par

172. 但没有什么比正义和公平更为和谐。因为这正义，作为善意的同伴和盟友，使我们爱那些我们认为与自己相似的人。再者，善意本身也包含刚毅。因为当友谊源于善意的泉源时，它便会毫不犹豫地为朋友承受生命的巨大危险。它说：\BibleQuote{若因他而有灾祸临到我，我必承当。}\BibleRef{德 23:31}\par

% structure: p0246 source=h2 target=subsection reason=relative-heading-rank; base=h2

\subsection{第三十四章}

\Emph{这里列举了善意的其他一些益处。}\par

173. \Emph{善意}也常能除去愤怒之剑。善意也使朋友的伤痕胜过仇敌的甘吻。\BibleRef{箴 27:6} 善意又使多人合而为一。因为若多人成为朋友，他们就合而为一；在他们内只有一个心灵和一个意见。我们也注意到，在友谊中，规劝是令人愉快的。它们有刺，却不引起痛苦。我们被责备的言辞刺痛，但乐于见到善意所显示的关怀。\par

174. 最后结论是，并非人人应尽同样的责任。也不可偏待人，尽管通常要考虑事件的时机和情况，以至于人有时必须帮助邻人，而非自己的兄弟。因为\Person{撒罗满}也说：\Quote{近处的邻人胜过远方的弟兄。}\BibleRef{箴 27:10} 因此，人通常将自己交托给朋友的善意，而非与弟兄的血缘关系。善意如此普遍，以至于它常常超越了天性的担保。\par

% structure: p0250 source=h2 target=subsection reason=relative-heading-rank; base=h2

\subsection{第三十五章}

\Emph{论刚毅。这分为两部分：关于战争事务和关于家庭事务。第一部分除非与正义和明智结合，否则不能成为德行。另一部分主要依赖于忍耐。}\par

175. 我们已经从正义的角度充分讨论了德行的性质和力量。现在让我们讨论刚毅，它（作为一种比其他德行更高超的德行）分为两部分，因为它涉及战争事务和家庭事务。但战争事务的思想似乎与我们职务的责任无关，因为我们更多关注灵魂的责任，而非肉体的责任；我们也不该关注武器，而应关注和平的事务。然而，我们的先祖，如\Person{农}的儿子\Person{若苏厄}、\Person{耶鲁巴耳}、\Person{三松}和\Person{达味}，在战争中也赢得了极大的荣耀。\par

176. 因此，刚毅是一种比其他德行更高超的德行，但它也绝不孤立存在。因为它从不只依赖自己。而且，没有正义的刚毅是邪恶的根源。因为它越强大，就越容易压服弱小者，而在战争事务中，人们应该看战争是正义的还是不义的。\par

177. \Person{达味}除非被迫，从不发动战争。因此在战斗中，明智伴随刚毅与他同在。因为即使他即将与巨人\Person{哥肋雅}单打独斗，他也丢弃了身上所负的盔甲。他的力量更依赖于自己的手臂，而非他人的武器。然后，在远处，为了更有力的一掷，他仅一石就击杀了敌人。此后，他从未在不寻求上主的劝告下投入战争。因此他在所有的战争中都获得了胜利，甚至直到晚年仍准备好战斗。当与\Place{培肋舍特}人的战争爆发时，他加入战斗，与他们凶猛的军队交锋，渴望赢得荣耀，而不顾自己的安全。\par

178. 但这并非唯一值得注意的刚毅类型。我们认为那些以伟大的心怀，\BibleQuote{因着信德，堵住了狮子的口，熄灭了烈火的威力，脱开刀剑的锋刃，从软弱中变为刚强。}\BibleRef{希 11:33} 的人的刚毅是光荣的。他们并非与许多人一同获胜，既非被同伴环绕，也非由军团援助，而是仅凭自己灵魂的勇气，独自赢得对奸诈仇敌的胜利。\Person{达尼尔}是多么不可战胜，他不怕周围咆哮的狮子。野兽咆哮，而他却在进食。\par

% structure: p0256 source=h2 target=subsection reason=relative-heading-rank; base=h2

\subsection{第三十六章}

\Emph{刚毅的职责之一是保护弱者免受伤害；另一则是抑制我们灵魂中的错误冲动；第三，既要轻视羞辱，又要以平静的心灵做正确的事。所有这些显然应为所有基督徒，尤其是圣职人员所履行。}\par

179. 因此，刚毅的荣耀不只在于身体或臂膀的力量，更在于心灵的勇气。刚毅的法则不在于造成伤害，而在于驱除一切伤害。那能保护朋友免遭伤害却不去做的人，与那造成伤害的人同样有罪。因此，圣\Person{梅瑟}以此作为他在战争中刚毅的首次证明。当他看到一个希伯来人受到一个埃及人的苛待时，他保护了他，击倒了那埃及人，并将他埋在沙中。\BibleRef{出 2:11} \Person{撒罗满}也说：\Quote{解救那被带往死亡的人。}\BibleRef{箴 24:11}\par

180. 那么，\Person{西塞罗}和\Person{帕奈提乌}，甚至\Person{亚里士多德}，从哪里得到这些观念，就十分清楚了。因为虽活在这两人之前，\Person{约伯}已说过：\Quote{我搭救了贫穷者脱离强暴者的手，扶助了无依无靠的孤儿。那源于亡者之口唇边的祝福都临于我身。}\BibleRef{约 29:12} 他如此高贵地承受了魔鬼的攻击，并以心灵的力量战胜了它，岂不是最勇敢的吗？我们也没有理由怀疑他的刚毅，上主曾对他说：\Quote{束上你的腰，像一个勇士；穿上尊荣和权能；贬抑一切作恶之人。} 宗徒也说：\Quote{我们有坚固的安慰。}\BibleRef{希 6:18} 那么，在一切忧伤中找到安慰的人，才是勇敢的。\par

181. 而且的确，当一个人战胜自己，克制愤怒，不屈从于任何诱惑，不因厄运而沮丧，也不因顺境而自傲，又不随风飘荡般被各种变化左右时，此便恰当地称为刚毅。然而，还有什么比训练心灵、克制肉身、使之顺从，从而能听从命令、服从理智、并在经受劳苦时欣然完成心灵意旨和愿望的事，更为高贵和光彩的呢？\par

因此，这就是刚毅的第一个概念。因为心灵的刚毅可以从两方面来看。第一，它将所有身外之物视为极其无足轻重，把这些看作与其说是应追求的，不如说是多余的该受轻视的。第二，它尽力追求那些最崇高的事物，以及一切可以从中看出某种道德之美（或如希腊人所称的\OriginalTerm{适宜}{\GreekText{πρέπον}}）的事物，且是用全部的心力。因为有什么比训练自己的心灵，使它不看重财富、享乐和荣誉，也不把所有心思浪费在这些事上更为高贵呢？当你的心灵如此调适，你必须考虑一切有德且合宜的事如何必须置于万事之先；且必须如此坚定你的心意，以致如果发生任何能扰乱你心神的事，无论是财产损失、获得较少荣誉，还是不信者的贬损，你都毫不为所动，仿佛你超越这类事物之上；甚至那威胁你安全的危险，如果是为正义而承担，也不能搅扰你。\par

这是基督的精兵所拥有的真刚毅，他若不按规矩比武，便不得受赏。\BibleRef{弟后 2:5} 或者，那对刚毅的呼唤在你看来是否卑微：\BibleQuote{磨难生忍耐，忍耐生老练，老练生望德}？\BibleRef{罗 5:3} 请看，竞赛虽多，冠冕却只有一个！那种呼唤无人给予，惟有在基督耶稣内得到坚固的，他的肉身从未有过安宁。处处有患难，外有争战，内有恐惧。\BibleRef{格后 7:5} 虽然处于危险、无数的劳苦、监禁、死亡之中——他却不气馁，反而作战，使他的软弱成了他的力量。\par

那么，请想一想，他如何教导那些开始在教会中尽职的人，应轻视一切世俗之物：\BibleQuote{所以你们若是与基督同死，脱离了世俗的虚妄，为什么仍像活在世俗中一样，服从那\InnerQuote{不可拿，不可尝，不可摸}的规条呢？这一切都是为使用而终必败坏的。} 又说：\BibleQuote{所以你们若与基督一同复活了，就该寻求天上的事，不要思念地上的事。}\BibleRef{哥 3:1} 又说：\BibleQuote{所以要致死你们在地上的肢体。}\BibleRef{哥 3:5} 这诚然是对所有信众说的。但他特别催促你，我的儿子，要轻视财富，摈弃渎神的和如老妇闲扯的荒诞故事——只许这样：\BibleQuote{要操练自己于虔敬，因为身体的操练益处不多，唯独虔敬在各方面都有益处。}\BibleRef{弟前 4:8}\par

所以，要让虔敬操练你走向正义、节制、温和，使你避免幼稚的行为，并在恩宠中扎根建基，好能打这场有关信德的好仗。\BibleRef{弟前 6:12} 你不可被今生的俗务缠身，因为你是为天主作战。\BibleRef{弟后 2:4} 因为如果那为皇帝作战的士兵，被人间的法律禁止牵涉诉讼、处理法律事务或经营买卖，那么投身信仰之战的人，岂不更当避开一切俗务，只满足于自己一小块土地（如果有的话）的出产？如果没有，也当满足于自己服役所得的酬劳。关于此事有一个好见证，他说：\BibleQuote{我从前年幼，现在年老，却从未见过义人被弃，也未见他的后裔讨饭。} 这才是心灵的真正安息和节制，它既不因贪图利益而激动，也不因害怕匮乏而苦恼。\par

% structure: p0265 source=h2 target=subsection reason=relative-heading-rank; base=h2

\subsection{第三十七章}

\Emph{在逆境和顺境中都应保持平静的心。但恶事必须躲避。}\par

这也是心灵不受烦扰的真正自由，它使我们既不在忧伤中过度消沉，也不在顺境中过分得意。如果那些敦促人们从事国家事务的人尚且给出这些规条，那么我们这些蒙召在教会中尽职的人，岂不更应当这样行事，做那些悦乐天主的事，好使基督的德能在我们身上彰显出来吗？我们也必须向我们的元帅证明自己，使我们的肢体成为正义的武器；不是罪恶寄身的属血肉的武器，而是为天主坚固的武器，好藉以除去罪恶。让我们的肉身死亡，以便在它之内一切罪恶皆死。并且，如同死后复生，让我们起来从事新的工作和度新的生命。\par

这些就是刚毅的服务；并充满了有德且合宜的职责。然而，在我们所做的一切事上，必须留意不单看它是否合乎道德，还要看它是否可能，免得我们从事那不能完成的事。因此，主，用祂自己的话说，愿意我们在受迫害时从一城逃到另一城；\BibleRef{玛 10:23} 这样，没有人会因贪求殉道的荣冠，而冒昧自置于危险之中，可能软弱的肉体或放纵的心灵无法承受和忍耐这样的危险。\par

% structure: p0269 source=h2 target=subsection reason=relative-heading-rank; base=h2

\subsection{第三十八章}

\Emph{我们必须坚强心灵以对抗未来的忧患，并藉着提前预见来坚固它。这样做有何困难。}\par

但是，没有人应当出于胆怯而退避，或因害怕危险而放弃信德。为了坚定不移，从不被任何恐惧扰乱，不被任何忧患摧垮，也不向任何折磨屈服，灵魂该当以何等大的恩宠装备起来，心灵该当如何加以训练和教导啊！这些确实难以承当！但正如在更大痛苦的恐惧中，所有较小的痛苦似乎都减轻了，同样，如果你以宁静的谋略修养你的灵魂，并决定不偏离你的道路，又将天主审判的畏惧和永罚的痛苦摆在面前，你必能获得心灵的坚忍。\par

如果一个人如此预备自己，他便显出了极大的勤奋。另一方面，如果一个人能凭其心灵的预见力，预知未来，并仿佛将可能发生的事摆在眼前，决断倘使它发生时应当如何做，这便是天赋的才能。也可能他会同时思考两或三件可能单独或一起发生的事，并决定如何处理它们，无论单独或一起发生，他都认为能处理至最为有利。\par

200. 所以，勇敢者的本分是，当任何事威胁时，不闭上眼睛，而是把它摆在面前，在心灵的镜子中审视它，并以预见的思想迎接未来，以免事后不得不说：\Quote{这临到我身上，是因为我原以为它不会发生。}若不预先思量不幸，它们就会迅速支配我们。在战争中，不期而至的敌人难以抵挡，若发现对方未作准备，就轻易战胜他们；同样，未经思虑的邪恶容易摧垮灵魂。\par

200. 那么，灵魂的卓越在于这两点：你的灵魂，受过善念的训练，怀有纯洁的心，第一，能看清何为真实与德行（因为\BibleQuote{心里洁净的人是有福的，因为他们要看见天主}\BibleRef{玛 5:8}），并能断定唯有德行才是善；其次，永不被任何事务扰乱，也不被任何欲望摇荡。\par

201. 但这并非对任何人都是易事。因为，有什么比从守望台上分辨智慧的资源和那些多数人视为伟大高贵的事物更困难呢？再者，有什么比将决断置于稳固的根基，并蔑视已被断定为无价值、毫无益处之事更困难呢？还有，当某种不幸发生，被视为严重而令人哀伤的事时，以这样的方式忍受它，即视它无非是自然之事，如同经上所说：\BibleQuote{我赤身脱离母胎，也要赤身归去。主赐的，主收回。}\BibleRef{约 1:21}（说这话的人丧失了子女和财产），并在一切事上保持智者和正直人的品行，如那位所言：\BibleQuote{主看着好，就这样成就了。愿主的名受赞美。}\BibleRef{约 1:21} 又说：\BibleQuote{你说话真像一个糊涂女人！难道我们只由天主接受恩惠，而不接受灾祸吗？}\BibleRef{约 2:10}\par

% structure: p0276 source=h2 target=subsection reason=relative-heading-rank; base=h2

\subsection{第三十九章}

\Emph{必须以刚毅对抗一切恶习，尤其对抗贪婪。圣约伯教导了这一课。}\par

202. 那么，灵魂的\Emph{\OriginalTerm{刚毅}{Fortitude}}并非无足轻重，也不与其他德行分离，因为它结合众德作战，单独捍卫一切德行的优美，守护其分辨能力，并以不共戴天的仇恨对抗一切恶习。就劳苦而言，它不可征服；为忍受危险，它勇敢无畏；面对享乐，它严厉无情；对诱惑则刚硬不阿，不知如何听信，甚至可以说，连招呼也不打。它不关心金钱，逃避贪婪如同逃避毁灭一切德行的瘟疫。因为，没有任何事比让自己为利益所胜更与刚毅相悖。常有仇敌被击退、敌军溃逃之时，战士忙于收拾阵亡者的战利品，倒在自己所杀的人中凄惨死去；军队忙于掳物，反将被逐之敌召回，以致胜利被夺走。\par

203. 因此，刚毅必须击退这污秽的瘟疫，并将它粉碎。它不可让自己受欲望诱惑，也不可为恐惧所动摇。德行忠于自身，勇敢地追究一切恶习，视之为德行的毒药。它必须如同以武器击退愤怒，因为愤怒使谋略远离。它必须躲避愤怒，如同躲避重病。此外，它必须提防对荣耀的贪求，因为过分急切地追求，常造成危害，而一旦获得，也总无裨益。\par

204. 这一切在圣约伯身上，或在他的德行中，有何缺失？又有什么恶习临到了他？他如何忍受疾病、寒冷或饥饿的苦楚？他如何面对威胁他安全的危险？他那大量施舍穷人的财富，是藉掠夺积聚的吗？他何时允许对财富的贪婪、对享乐的渴求，或肉欲在他心中升起？那三位首领不友善的争论，或奴隶们的辱骂，可曾激起他的愤怒？当他呼唤加祸于己，若他曾隐藏哪怕无心的过失，或因惧怕众人而不敢在众人面前承认时，荣耀可曾像左右反复之人那样把他卷走？他的德行，与任何恶习毫无接触，却坚立于自己的根基之上。那么，谁像圣约伯那样勇敢？谁能屈居任何人之次？因为几乎无人能与他相提并论。\par

% structure: p0281 source=h2 target=subsection reason=relative-heading-rank; base=h2

\subsection{第四十章}

\Emph{我们的祖先不乏战场上的勇气，这由古人的表率，特别是\Person{厄肋阿匝尔}的辉煌事迹所证明。}\par

205. 但也许战场上的声誉如此紧紧地束缚某些人，使他们以为刚毅仅见于战斗，因而我此前转而言及这些事，是因为我们在战场上缺乏刚毅。然而，\Person{农}的儿子\Person{若苏厄}多么勇敢，他曾在一场战斗中一举击杀了五位国王及其人民！\EditorialNote{\BibleBook{若苏厄}{书}10章} 再者，当他与基贝红人作战，担心黑夜会阻止他获胜时，他以深厚的信德和高昂的精神喊道：\BibleQuote{太阳，停住吧！}\BibleRef{苏 10:12} 太阳便停住了，直到胜利完成。\Person{基德红}率领三百人战胜了一个大民族和残忍的敌人。\EditorialNote{\BibleBook{民长}{纪}7章} \Person{约纳堂}年轻时在战斗中表现了极大的勇气，关于玛加伯们，我又该说什么呢？\par

206. 首先，我要论及我们祖先的子民。他们随时准备为天主的圣殿和自己的权利而战，当在安息日被敌人的诡计袭击时，他们宁愿让未加防护的身体遭受创伤，也不卷入战斗，以免玷污安息日。他们全都欣然赴死。但是玛加伯家族认为，如果那样，整个民族都将灭亡，于是当他们在安息日被挑战作战时，也为无辜弟兄们的死报仇雪恨。此后，\Person{安提约古}王受此刺激而加倍努力，在他的将领\Person{里息雅}、\Person{尼加诺尔}和\Person{哥尔基雅}的率领下重新开战，结果他和他的东方及亚述军队被彻底击溃，在战场上遗尸四万八千具，而击败他们的军队只有三千人。\par

207. 请注意领袖\Person{犹大玛加伯}的英勇，正如在他的一名士兵的品格中所展示的那样。\Person{厄肋阿匝尔}\BibleRef{加上 6:43} 遇见一头比其余所有战象都更高大、全身披挂王室装饰的大象，以为国王骑在上面，便急速奔跑，冲入阵营当中；他抛掉盾牌，用双手击杀阻挡他的人，直到抵达那头巨兽跟前。然后他钻到象身之下，刺入自己的剑，将它杀死。但那头巨兽倒下时压住了厄肋阿匝尔，使他丧命。那时他心灵的勇气何等伟大：首先，他不畏惧死亡；其次，当他被敌人包围时，他被这种勇气带入敌阵最密集之处，直插中心！随后，他藐视死亡，抛掉盾牌，奔到巨兽身下，用双手将它刺伤，并任它压在自己身上。他钻到兽身之下，是为了给予更致命的一击。他被兽身倒下所包围，与其说是被压碎，不如说是埋葬在自己的凯旋之中。\par

208. 尽管他被王室的装饰所蒙蔽，但他的意图并未受骗。因为敌人被如此英勇的表现所震惊，竟不敢向这个手无寸铁、已被困住的人发起冲击。在头兽被杀死的意外之后，他们如此恐惧，以至于认为自己完全无法与一个人的勇猛相匹敌。不但如此，\Person{安提约古}王，即\Person{里息雅}的儿子，被一人的英勇所震慑，就求和了。他率领十二万武装士兵和三十二头战象前来作战；这些战象在日光照耀下，闪烁着兵器的光芒，犹如一排燃烧的灯烛，一头接一头地行进，大小如同群山。就这样，\Person{厄肋阿匝尔}将和平留作他的勇敢的继承人。这些就是凯旋的标记。\par

% structure: p0287 source=h2 target=subsection reason=relative-heading-rank; base=h2

\subsection{第四十一章}

\Emph{在称赞了\Person{犹大}和\Person{约纳堂}的高远心志之后，接下来要展示的是殉道者在忍受酷刑时所表现的坚忍，这也是\OriginalTerm{刚毅}{fortitude}之德的一个重要部分。}\par

209. 但是，既然刚毅不仅通过顺境，也通过逆境得到证明，现在让我们来思考\Person{犹大玛加伯}的死亡。他，在\Person{德默特琉}王的将军\Person{尼加诺尔}被击败之后，率领九百人勇敢地与国王的两万大军交战；这九百人因害怕被如此庞大的军队压倒而想要撤退，但他说服他们宁可光荣战死，也不可羞辱地逃跑。他说：\Quote{不要让我们，}他说，\Quote{在我们的光荣上留下任何污点。}于是，这样投入战斗，从日出战至黄昏，他攻击并迅速击退了右翼，那里他见到敌人最强劲的部队。但当他从后方追击逃敌时，给了敌人可乘之机对他造成创伤。\BibleRef{加上 9:8} 就这样，他为自己找到了一个比任何凯旋都更加光荣的死亡之地。\par

210. 我何必再提他的兄弟\Person{约纳堂}呢？他仅以一小队人马对抗国王的军队。\BibleRef{加上 11:68} 虽然他被自己的部下离弃，只剩下两个人，他却挽回战局，击退敌人，并召回那四散奔逃的部下，分享他的胜利。\par

211. 那么，这就是战争中的刚毅，它深深铭刻着德行与合宜的印记，因为它宁要死亡，不要为奴和羞辱。但是，对于殉道者们的苦难，我又能说什么呢？不必远求，玛加伯的儿子们不是对骄傲的\Person{安提约古}王赢得了与他们祖先一样伟大的凯旋吗？后者确实有武装，而他们却没有武器就获得了胜利。那七兄弟的队伍，虽被国王的军团包围，却依然坚不可摧——酷刑失效，施刑者停手；而殉道者们却毫不屈服。有一人，头皮被剥去，虽然容貌改变，勇气却增长了。另一人，被命令伸出舌头以便割掉，回答说：\Quote{天主不但倾听说话的人，因为当\Person{梅瑟}沉默时，他也倾听了他。他更垂听自己子民无言的思念，胜过其他人的一切声音。你惧怕我舌头的鞭笞——难道你不惧怕鲜血洒在地上所发出的鞭笞吗？鲜血也有声音，它向天主大声呼喊——如同\Person{亚伯尔}的血那样。}\par

212. 关于那位母亲\BibleRef{加下 7:20} 她欢欣地看着自己孩子们的尸首，如同看着许多战利品，并从垂死儿子们的声音中找到喜乐，仿佛是在歌手的歌曲中，她从孩子们身上觉察出自己心中那光荣竖琴的旋律，以及一种比任何弦琴所能发出的更甜美的爱的和谐。\par

213. 关于那些\Place{白冷}的两岁儿童\BibleRef{玛 2:16} 他们在尚未感受到自己内在的自然生命之前，就获得了胜利的棕榈枝。至于\Person{圣依搦斯}，当她面临贞洁与生命这两件大事的危险时，她保护了自己的贞洁，并以生命换取了不朽。\par

214. 也不可忽略\Person{圣老楞佐}，他眼见自己的主教\Person{西斯笃}被带去殉道，便痛哭起来，并非为主教所受之苦，而是因自己不得不留下。他这样开口对主教说：\Quote{父亲，你没有儿子跟随，往哪里去？圣善的司铎，你没有执事随行，往哪里急奔？你素常从未不在侍从陪同下举行祭献。我有什么地方令你不悦，父亲？你发现我不配吗？那么，请试试你是否拣选了一个合适的仆人。你曾将救主圣血的祝圣委托于我，曾恩准我共享圣事，如今却拒绝让我与你一同赴死吗？提防你良好的判断受到危害，而你的刚毅却只赢得称赞。抛弃弟子便是师长的损失；或者说，高贵的杰出人物岂不更常因门徒而非自己的比武获胜？\Person{亚巴郎}献上自己的儿子，\Person{伯多禄}派遣\Person{斯德望}先行！父亲，在儿子身上彰显你的勇气吧。献上你所训导的我，好让你深信对我的选择，在相称的同伴中获得荣冠。}\par

215. 于是\Person{西斯笃}说道：\Quote{我绝不离开你，也不舍弃你。更重大的斗争还在等待你。我们老年人必须接受较轻松的战斗；你这年轻人则等待着一场对暴君更光荣的凯旋。你很快就要到来。停止哭泣；三天后你将随我来。这个间隔必须放在司铎和他的肋未人之间。你不该在师傅的眼皮底下获胜，好像需要一位帮手似的。你为何寻求分担我的死亡？我将这死亡的完整遗产留给你。你何需我的临在？让软弱的门徒走在师傅之前，而让勇敢的跟从他，好叫他们在没有他的情况下获胜。因为他们不再需要他的引导。\Person{厄里亚}就这样留下了\Person{厄里叟}。我把我自己勇气的完整继承权托付于你。}\par

216. 这便是他们的争论，确实是一场相称的争论，司铎与侍从在其中争辩谁该首先为基督的名字受难。据说，当那段悲剧在剧院演出时，\Person{皮拉德斯}自称是\Person{俄瑞斯忒斯}，而\Person{俄瑞斯忒斯}却声明他才是自己本人，此时剧场里响起了热烈的掌声。前者这般作为，是为\Person{俄瑞斯忒斯}而死，而\Person{俄瑞斯忒斯}则不愿\Person{皮拉德斯}代己受戮。但他们本不该活下来，因为两人都犯了弑亲之罪：一个亲自犯下罪行，另一个协助完成。然而，此处却没有任何事促使\Person{圣老楞佐}如此行事，只出于他的爱德与虔敬。不过，三日后，他被自己讥嘲的暴君放置在铁烤架上，活活烧死。他说：\Quote{肉已烤好，翻过来吃吧。} 这样，藉由心灵的勇气，他胜过了烈火的力量。\par

% structure: p0297 source=h2 target=subsection reason=relative-heading-rank; base=h2

\subsection{第42章}

\Emph{不可无故激怒掌权者。勿听谄媚之言。}\par

217. 我想我们必须谨慎，以免因过分贪求荣耀，而滥用当权者，并激起那些反对我们的外教人的心思，使其渴望迫害，引发他们的愤怒。有些人使多少人丧亡，而自己却坚持到底，克胜了酷刑！\par

218. 我们也须留意，不要向谄媚者开启我们的耳朵。让人用谄媚来软化自己，似乎不仅是缺乏刚毅，更是真正的怯懦的表现。\par

% structure: p0301 source=h2 target=subsection reason=relative-heading-rank; base=h2

\subsection{第43章}

\Emph{论\OriginalTerm{节制}{temperance}及其主要部分，尤其\OriginalTerm{心灵的宁静}{tranquillity of mind}与\OriginalTerm{中和}{moderation}，对\OriginalTerm{德行}{virtue}之关爱，及对\OriginalTerm{合宜之事}{seemly}的反思。}\par

219. 我们既已论述了四枢德中的三项，就只剩下第四项需要讨论。这称为节制和中和；其中首先寻求并期待的是：心灵的宁静，达成温和，中和的优雅，对德行的关注，以及对合宜之事的反思。\par

220. 我们在生活中必须遵守某种秩序，好使我们能以最初的端庄感受打下基础，因为端庄是心灵平静的友伴与盟助。避免过分自信，厌恶一切过度，它喜爱节制，守护尊贵之事，只寻求合宜之事。\par

221. 其次便是交往的选择。我们应当与那些年长的、品行经过考验的人结交。因为与同辈人的相处虽然更为愉快，但与长者的交游却更为安全。藉由他们的引导和行事为人，他们给年轻人的品格涂上了色彩，仿佛以正直的深紫色浸染了他们一般。如果对某地陌生的人非常乐意在经验丰富的向导陪伴下旅行，那么年轻人更应当在长者的陪伴下踏上对他们而言全新的生命之路；这样他们才不会走错，或偏离德行的正路。因为再没有比有同一个人既指导我们的生活，又见证我们如何生活更好的事了。\par

222. 而且在一切行为中，我们也须考虑什么适合不同的人、不同的时机和年龄，以及什么与各人能力相符。因为常常适合一人的，并不适合另一人；一件适合青年人，另一件适合老年人；一件适合危难时，另一件适合顺境时。\par

223. \Person{达味}在主的约柜前跳舞。\Person{撒慕尔}却没有跳舞；然而\Person{达味}并未受责备，另一位却受到称扬。\Person{达味}在名叫\Person{阿基士}的国王面前改变了面容。\BibleRef{撒上 21:13} 如果他这样做时丝毫没有害怕被人认出，那他肯定无法逃脱轻浮的指责。\Person{撒乌耳}也曾在先知群中受感说话。但唯独论到他，仿佛他不配似的，经上说：\BibleQuote{撒乌耳也在先知中吗？}\BibleRef{撒上 19:24}\par

% structure: p0308 source=h2 target=subsection reason=relative-heading-rank; base=h2

\subsection{第44章}

\Emph{每人应致力于适合其性格的职责。然而许多人因追随父辈的职业而受阻。圣职人员则不同。}\par

224. 各人都认识自己的才能。因此，每人应致力于那他认为适合自己的事。但他必须首先考虑其后果。他或许知晓自己的优点，但也须知晓自己的缺点。他还须公正地判断自己，以便追求善而避免恶。\par

225. 有人更适合作读经员，另一人更善于歌唱，第三人更热心于驱逐附魔者，还有人被认为更适合管理圣物。司铎应留意这一切。他应分派给每人那最适合他的职务。因为每人随其心意所向，或任何适合他的职务，那职位或职务便更能充满恩宠。\par

226. 然而，在一切生活境况中这皆属难事，于我们而言则最为艰难。因为各人惯于随从父母的生活选择。因此，父辈从军者，通常亦入伍从军。其他各行业亦然。\par

227. 然而，在圣职界中，子承父业者最为罕见，或因艰难的职务使其却步，或因在动荡的青年时期持守贞洁过于困难，或因这生活对活跃的青年显得过于平静。于是他们转向那些被认为更为显赫的职业。的确，多数人重今生而轻来世。他们为今生而奋斗，我们则为来世。因此，我们所从事的事业愈伟大，便愈须专心致力于它。\par

% structure: p0314 source=h2 target=subsection reason=relative-heading-rank; base=h2

\subsection{第45章}

\Emph{论高尚与德行，及其间之区别，如世俗作者与神圣作者所述。}\par

228. 因此，我们要坚守端庄，和那为整个生命增添华美的节制。因为在凡事上保持适当的尺度并建立秩序，绝非小事，而其中显著昭然者，即我们称之为\Quote{\OriginalTerm{端庄}{decorum}}，或得体者。这与德行如此紧密相连，以致二者不可分离。因为得体者亦有德，而有德者亦得体。故此，其区别在言语而非在事物本身。二者之间确有分别，我们可以理解，却无法解释。\par

229. 为尝试对二者略作区分，我们可以说，德行可比作身体的良好健康和健全，而得体则如其秀雅与美丽。正如美丽似乎高于健全与健康，却非有它们便不能存在，也不能以任何方式与之分离——因为除非一人享有健康，便不能有美丽与秀雅——同样，德行本身亦包含着得体，以致二者似乎一同开始，且不能无彼此而存在。因此，德行如同我们一切工作与事业中的健全；得体则好似外在表现，它与德行结合时，只能在我们的思想中被分别认识。因为尽管在某些情况下它似乎显眼突出，其根却在德行之中，尽管花朵属于其自身。根于此，则繁荣茂盛；否则便衰败枯萎。何谓德行，若非避免可耻之事如同避免死亡？而德行的反面，若非那带来荒芜与死亡的事？因此，若德行的本质强壮有力，得体亦会迅速萌发如花，因其根是健全的。但若其意愿之根腐败，便无物能由之生长。\par

230. 在我们的经卷中，此理说得更为明白。因为达味说：\BibleQuote{上主为王，以尊威作衣冠。} 宗徒也说：\BibleQuote{行动要端庄，好像在白天一样。} \BibleRef{罗 13:13} 希腊原文作\OriginalTerm{\GreekText{ευσχημόνως}}{\GreekText{ευσχημόνως}}——意即：穿着得体，容貌端正。当天主造第一个人时，祂以美好的体型、匀称的四肢造了他，并赐予他极为高贵的容貌。祂并未赐予他罪的赦免。但后来，那以仆人的形体、以人的模样来临者，用祂的圣神更新了他，将祂的恩宠倾注于他心中，并将人类救赎的尊荣加于己身。因此，先知说：\BibleQuote{上主为王，以尊威作衣冠。} 他又说：\BibleQuote{天主，在熙雍，赞美你本是合宜。} 这就是说：敬畏你、爱慕你、祈求你、尊敬你，乃是适宜和美好的，因为经上记载：\BibleQuote{一切都该照规矩按次序而行。} \BibleRef{格前 14:40} 但我们也能敬畏、爱慕、祈求、尊敬人；然而赞美歌尤其向天主呈献。我们向天主呈献的这种得体，我们可信其远胜其它事物。女人以端庄的服装祈祷也是合宜的，\BibleRef{弟前 2:9} 但尤为合宜的是，她应蒙头祈祷，并在祈祷时，以良好的品行承诺纯洁。\par

% structure: p0319 source=h2 target=subsection reason=relative-heading-rank; base=h2

\subsection{第46章}

\Emph{论得体的双重划分。继而阐明，合乎本性者即是有德的，反之则必被视为可耻。并举例解释这一划分。}\par

231. 因此，显赫的得体有双重划分。其一可谓之为普遍的得体，它遍及一切德行，且可说见于整个身体。亦有特殊的得体，它清晰显于某个特定部分。前者具有一致的形式，以及德行在每一行动中和谐完备。因为其整个生命与自身一致，凡事毫无矛盾。后者则关乎在有德行的生活过程中所行的任何特殊行动。\par

同时我们要注意，合于本性地生活，并按本性地度过时光，是合宜的；相反本性之事则是可耻的。因为宗徒问道：\BibleQuote{女人不蒙头向天主祈祷，这相宜吗？本性自身不是也教导你们：男人若蓄长发，为他是羞辱吗？因为这是反本性的。}他又说：\BibleQuote{女人若蓄长发，为她是光荣。}\BibleRef{格前 11:13} 这是合于本性的，因为她的头发是给她当作帕巾的，因为这是本性的帕巾。因此本性为我们安排品格和仪表，我们应当遵守她的引导。但求我们能保守她的纯洁，不要因我们的邪恶而改变我们所领受的！\par

我们具备那种普遍的端庄；因为天主创造了这个世界的美。在其部分中也具备它；因为当天主造了光，分别昼夜，造了穹苍，分开陆地与海洋，设置日月星辰照耀大地时，祂逐一地赞许它们。因此，这光辉焕发于世界每一部分的端庄，闪耀于整体，正如\BibleBook{}{智慧篇}所表明的，说：\BibleQuote{我在祂那里，祂在完成世界的工程时，曾欢欣踊跃。} 同样，在人体的构造中，每个肢体都令人愉悦，但所有肢体搭配得宜则更令我们欣喜不已。因为它们似乎联结并配合成为一个和谐的整全。\par

% structure: p0324 source=h2 target=subsection reason=relative-heading-rank; base=h2

\subsection{第四十七章}

\Emph{我们应当在生活中处处彰显出端庄。那么，我们应当让哪些私欲偏情发展到极限，又应当克制哪些呢？}\par

若有人在整个人生中保持平稳的进程，在一切所行的事上有条不紊，察见自己的言语有秩序、前后一致，行为有节制，那么端庄便在这样的人的生活中格外显著，如同在镜子里闪耀。\par

此外，还应当有令人愉悦的说话方式，以便赢得听者的好感，并尽可能使自己与所有朋友和同邦居民和睦相处。不要让人显得喜好谄媚，也不要让人渴望从任何人那里得到谄媚。前者是狡诈的标记，后者是虚荣的标记。\par

不要轻视他人——尤其是善人——对自己的看法，因为这样他就学会了重视善人的意见。蔑视善人的评判，要么是自负的标志，要么是软弱的标志：前者出于骄傲，后者出于懈怠。\par

我们也要提防灵魂的冲动。灵魂总该警醒自省，防范自己。因为有一些冲动，带有某种偏情，仿佛在一阵奔突中爆发出来。所以在希腊文中称之为\OriginalTerm{冲动}{\GreekText{ὁρμή}}，因为它带着某种力量突然冲出来。在这些冲动中，蕴涵着灵魂或本性不小的力量。然而这力量有两种：一方面建立在偏情之上，另一方面建立在理性之上，理性抑制偏情，使它服从自己，随己意引导它；又借细心的教导，训练它知道何事当做，何事当避，从而使它顺服于温和的管束。\par

因为我们应该小心，绝不要鲁莽或疏忽地行事，也绝不要做任何我们无法提供合理理由的事。虽然我们行事的原因无需向每个人明言，人人却都会察究它。事实上，我们也没有任何可借以推脱的东西。我们每一种偏情固然都有某种本性之力，但按照本性之律，偏情本身却服从于理性，并受理性支配。因此，细心警醒的人有责任如此守望，使偏情不超越理性，也不彻底离弃它，免得因超越造成混乱，理性被排斥，被这样的遗弃化为乌有。纷扰破坏恒常。退缩暴露怯懦，且意味着怠惰。因为心神烦乱时，偏情向四方扩散，在暴烈的发作中忍受不了理性的约束，也感受不到驾驭者的指挥，无法回心转意。因此，通常不仅灵魂被扰乱，理性失去，而且面容或因愤怒、或因邪情而涨红，或因恐惧而苍白，在逸乐中不能自制，在喜乐中不能把持。\par

当这样的事发生时，那本性的判断力和品格的尊严便被抛弃，那种唯有在行为与思想中才能保持自身权威与端庄的恒常，就不再能持守了。\par

然而，更激烈的偏情由过度的忿怒而生，那是我们所受的某种冤屈引起的痛苦点燃的。开篇的圣咏训词教导了我们这一点。因此，我们在论述职责时，采用那段开场白，真是恰到好处，因为那段话本身就关乎职责的引导。\par

不过，由于我们当时（这是合宜的）仅就此点略加触及——即如何小心避免在受冤屈时受到扰乱——以免开场论述过于冗长，我想现在略作更详细的讨论。此刻正是良机，因为我们正在谈论节德的不同部分，来看忿怒当如何克制。\par

% structure: p0334 source=h2 target=subsection reason=relative-heading-rank; base=h2

\subsection{第四十八章}

\Emph{再次提出克制忿怒的论证。随后论及忍受冤屈的三等人；据说宗徒和达味达到了最完美的一等。趁机阐明今生与来世的区别。}\par

我们愿尽可能指出，圣经中记载了三种受冤屈的人。其中一类是那些遭受罪人辱骂、诽谤、肆意践踏的人。正因正义不彰，羞耻便加深，痛苦亦加剧。我同阶层的许多弟兄、我同侪中的许多人都与此相似。倘若有人对软弱的我施以伤害，也许，我虽软弱，仍可宽恕他对我的不义。若他指控我犯有某罪，我并不满足于自己良心的见证，尽管我知道自己对他所控之事清白无辜；然而，正因我软弱，我渴望洗去我天性中羞耻的痕迹。所以我要求以眼还眼、以牙还牙，以辱骂还辱骂。\par

然而，倘若我是那虽未臻完美却在精进的人，我便不以辱骂还辱骂；若他大发谩骂，以种种侮辱充斥我耳，我则默然不语，不予回答。\par

但是若我是完人（我这样讲仅为举例，因为在事实上我是软弱的），那么，若我是完人，我便祝福那咒骂我的人，正如\Person{保禄}也曾祝福，因为他说：\BibleQuote{受辱骂，我们就祝福；}\BibleRef{格前 4:12}他曾听见那一位说：\BibleQuote{应爱你们的仇人，为那迫害你们的人祈祷。}\BibleRef{玛 5:44}这样，\Person{保禄}遭受迫害，并且坚忍到底，因为他克制并平复了自己属人的情感，为了那摆在他面前的赏报，即：他若能爱仇人，就将成为\OriginalTerm{天主的儿女}{son of God}。\par

我们也能表明，圣\Person{达味}在这同样的德性上，肖似\Person{保禄}。当\Person{史米}的儿子咒骂他，并指控他犯有重罪时，起初他默然不语，谦卑自下，甚至对自己所行的善工、即他所知的善行也噤口不言。后来，他甚至请求那人咒骂他；因为当他被咒骂时，他希望能获得主的怜悯。\par

但请看，他是如何积聚谦逊、正义和明智，以获享由主而来的恩宠！起初他说：\Quote{因此他咒骂我，因为主曾吩咐他说：\InnerQuote{你应咒骂\Person{达味}。}}这是谦逊；因为他认为凡\OriginalTerm{天主所安排的}{divinely ordered}事，都应以平和之心忍受，仿佛自己不过是个奴仆小童。接着他说：\Quote{看，我所生的儿子，竟要索我的命。}这是正义。因为我们若从自己家人手中遭受苦难，又何必对外人所加之于我者动怒呢？最后他说：\Quote{由他咒骂吧，因为那是主吩咐他的。也许主会垂视我的困苦，因他今天的咒骂而恩待我。}这样，他不仅忍受了辱骂，而且在那人向他扔石、尾随追击时，他也不加以惩罚。不但如此，在他获得胜利之后，当那人前来求恕时，他便宽仁地赦免了他。\par

我写这些是为了表明，圣\Person{达味}怀着真正的福音精神，不仅没有见怪，反而感激那辱骂他的人；他对自己的冤屈，甚至欢喜多于恼怒，因他以为如此便能为他换得某种赏报。但他虽已臻完美，仍寻求更为完美的境界。作为凡人，他因受屈的痛苦而心中愤激，但如同善战的勇士，他克胜了这一切，又像勇敢的摔跤者，他坚忍不拔。他忍耐的终向与目标在于期待恩许的完成，因此他说：\Quote{主啊，求你使我认识我的终期，和我的寿数几何，好叫我明了我是多么短促。}那么，他所寻求的乃是天上恩许的终向，那时各人将依照自己的次序复活：\BibleQuote{首先是为初果的基督，然后是在基督再来时属祂的人，随后是结局。}\BibleRef{格前 15:23}因为，就如\OriginalTerm{宗徒}{Apostle}所说，当国度的权柄被交付于\OriginalTerm{天主父}{God, even the Father}，一切权势都被制伏时，那时圆满才开始。所以，在此世有阻碍，在此世成全之人仍有软弱；而在那里才有全然的圆满。因此，他所祈求的，是那真实存在的永生之日，而非流逝的时光，好使他知道自己所欠缺的是什么，什么是那结常生果实的应许之地，哪一处在父家中的住所是第一座，哪是第二座，哪是第三座，在那里各人将按自己的功绩安息。\par

因此我们必须努力追求那有圆满和真理之处。此世是影，此世是像；\BibleRef{希 10:1}彼处才是真理。影在法律中，像在福音中，真理则在天上。古时，人们献上羊、献上牛犊；如今，献上的却是\OriginalTerm{基督}{Christ}。但祂是以人并承受苦难的身份被献上的。祂又作为\OriginalTerm{司铎}{priest}，亲献己身以除去我们的罪，在此世是影像中如此，而在彼处则是真理中如此，在那里祂与圣父同在，作我们的\OriginalTerm{中保}{Advocate}，为我们转求。所以，我们在此世是行于影像之中，见的是影像；而在彼处则是面对面相见，那里有全然的圆满。因为一切圆满全在真理之中。\par

% structure: p0343 source=h2 target=subsection reason=relative-heading-rank; base=h2

\subsection{第四十九章}

\Emph{我们必须在自己内保守诸德之像。必须祛除魔鬼与恶习之像，尤其是贪婪之像；因为贪婪剥夺我们的自由，并使那些沉溺于虚幻中的人，丧失\OriginalTerm{天主的肖像}{image of God}。}\par

那么，当我们尚在此世时，让我们保守此像，好使我们在彼处能得见真理。让正义之像存于我们内，同样也让智慧之像存于我们内，因为我们将迎接那日子，并依照我们内所存之像获得赏报。\par

不要让仇敌在你身上寻见他的肖像，莫让他寻见暴怒和狂怒；因为这些便是邪恶的肖像。\BibleQuote{你们的仇敌魔鬼，如同咆哮的狮子巡行，寻找可吞噬的人。}\BibleRef{伯前 5:8}不要让他寻见对黄金的贪欲，也不要让他寻见堆聚的银钱，以及诸般恶习的表象，免得他从你口中夺去自由的声音。因为，当你能说出这句话时，真正自由的声音便被听见了：\BibleQuote{这世界的元首将要来到，他在我内一无所有。}\BibleRef{若 14:30}因此，若你确信，当他前来搜察你时，在你内找不到什么，你便会如圣祖\Person{雅各伯}对\Person{拉班}所说的：\BibleQuote{现在请你察看，在你这里有什么属于你的东西。}\BibleRef{创 31:32}我们实在应称\Person{雅各伯}为有福的，因为\Person{拉班}在他那里竟找不到任何属于自己的东西。原来\Person{辣黑耳}已将他神像的金银像藏了起来。\par

251. 如果智慧、\OriginalTerm{信德}{faith}、轻看世俗和属神的\OriginalTerm{恩宠}{grace}排除了所有无信德，你们便是有福的；因为你们不顾念虚妄、愚昧和谎言。夺去对手说话的机会，使他无从控告你们，难道是小事吗？那不顾念虚妄的便不受扰乱；但顾念虚妄的就被扰乱，而且毫无益处。积蓄财富岂不是虚妄吗？追求转眼即逝的事物，确实足够虚妄。即使积蓄了，又怎能知道它们必归你所有呢？\par

252. 商人日夜奔波，为要积蓄财宝，难道不虚妄吗？他聚集货物，又为其价格惶惶不安，唯恐卖价低于买价，难道不虚妄吗？他到处谋求高价，却因夸耀生意而招来盗贼嫉妒，以致突遭抢劫；或不待风平浪静，不耐迟延，在求利途中遭遇海难，这岂不也是虚妄吗？\par

253. 那辛苦劳碌积蓄财富，却不知留给哪个继承人的，岂不也是徒然受扰吗？贪婪之人费尽心思所积攒的一切，常被挥霍无度的继承人以疾速的奢侈挥霍殆尽。那厚颜无耻的浪子，无视眼前，不顾将来，将长久积聚的财富如吞入深渊般耗尽。往往，那期望中的继承人所分得的遗产只招来嫉妒，又突然死去，将尚未到手的全部遗产转归外人。\par

254. 那么，你们为何徒然纺织无益无果的网罗？又为何堆砌无用的财宝，如同蜘蛛的网？因为即使满溢，也毫无益处；反而剥去你们天主的肖像，给你们穿上属土的肖像。若有人具有暴君的肖像，岂不应当受惩罚吗？你们弃绝了永恒君王的肖像，在自己内树立死亡的形象。不如从你们灵魂的国度中驱逐魔鬼的肖像，建立起基督的肖像。让这肖像在你们内照耀；让它在你们的国度，即你们的灵魂中，明亮地燃烧，因为它摧毁一切恶习的肖像。达味论到这事说：\BibleQuote{主啊，祢在祢的国里，使他们的形像归于虚无。}因为当主按照自己的肖像装饰耶路撒冷时，仇敌的一切肖像便被毁灭了。\par

% structure: p0351 source=h2 target=subsection reason=relative-heading-rank; base=h2

\subsection{第五十章}

\Emph{肋未人应当完全摆脱一切属世的欲望。按宗徒的教导，他们的德行应当如何，纯洁又当如何伟大。还有他们的尊荣和职责是什么，为此必须具备首要的德行。他指出，这些德行并非哲学家们所不知，只是他们弄错了次序。有些德行本性符合义务，但因所附带的情形而变得相反义务。由此，他推论出肋未人的职务要求何种恩赐。最后，他解释了梅瑟祝福肋未支派时所说的话。}\par

255. 如果在主的福音中，民众受教并被引导去轻看财富\BibleRef{谷 10:23}，那么你们肋未人岂不更应当不再受属地欲望的束缚吗？因为天主是你们的产业。当梅瑟将地上的产业分给我们祖先的民众时，主不许肋未人在那地上有份\BibleRef{户 18:23}，因为祂自己要作他们产业的根基。因此达味说：\BibleQuote{主是我的产业，是我杯中的份。}由此我们得到\Quote{肋未}一名，意思是：\Quote{祂是我的}，或\Quote{祂为我。}所以，他的尊荣甚大，天主竟指着他说：祂是我的。或者说，正如对伯多禄论及从鱼口中所得的钱币：\BibleQuote{你拿去，为他们和我，给你和我吧。}\BibleRef{玛 17:27}因此，\OriginalTerm{宗徒}{Apostle}说：\BibleQuote{\OriginalTerm{主教}{bishop}应当是自持、庄重、善行、好客、善于教导、不贪财、不争吵、善于管理自己的家庭，}又加上：\BibleQuote{\OriginalTerm{执事}{deacons}同样应当庄重，不一口两舌，不饮酒过度，不贪不义之财，以纯洁的\OriginalTerm{良心}{conscience}持守\OriginalTerm{信德}{faith}的奥秘。也要先受试验，若无过犯，方可任执事职。}\BibleRef{弟前 3:2}\par

256. 我们当留意对我们的要求何等重大。主的仆役应当戒酒，好使他不但获得信徒，也获得教外人士的好评。因为我们行为和工作的见证，应当是公众的意见，以免这职务受到玷辱。这样，人看见祭台的司役具备相称的德行，便会赞美他们的创造者，并敬畏拥有如此仆人的主。哪里家中有纯洁的产业和无辜的管理，主的赞颂便在那里传扬。\par

257. 关于贞洁，只允许一次婚姻而不允许第二次结合，我还能说什么呢？就婚姻而言，法律禁止再婚，也不得寻求与其他妻子结合。许多人感到奇怪，为何在领洗前缔结的第二次婚姻竟会造成障碍，妨碍人获选任圣职和领受\OriginalTerm{圣秩}{ordination}的恩宠；因为即使罪行，若在\OriginalTerm{圣洗圣事}{sacrament of baptism}中已被清除，通常也不构成阻碍。但我们必须明白，\OriginalTerm{圣洗}{baptism}中罪可得赦免，法律却不能废除。在婚姻的事上，并没有罪，却有法律。任何罪过都可除去，但在婚姻中，任何法律都不可搁置。自己多次结婚的人，怎能劝导别人守寡呢？\par

258. 但你们知道，\OriginalTerm{圣职}{ministerial office}必须保持纯洁无瑕，不可被\OriginalTerm{夫妻结合}{conjugal intercourse}所玷污；我告诉你们，你们领受了\OriginalTerm{圣职}{sacred ministry}的恩赐，常保持着纯洁的身体与无玷的贞洁，从不知夫妻结合为何物。我之所以提及此事，是因为在某些偏远地方，当人进入圣职，甚或成为\OriginalTerm{司铎}{priests}时，他们已有子女。他们以古老习俗为此辩护，说当时祭献的间隔很长。然而，据我们在\OriginalTerm{旧约}{Old Testament}中所读的，就连百姓也必须提前两三天洁净自己，以便洁净地前来参与祭献。\BibleRef{出 19:10} 他们甚至要洗净自己的衣服。如果在预像中尚且如此敬谨，那么在实体中又该如何加倍呢！那么，司铎和\OriginalTerm{肋未人}{Levites}，你们要明白，洗净你们的衣服意味着什么。你们必须有清洁的身体，才能奉献圣事。百姓尚且不得在未洗净衣服时接近祭品，你们内心和身体污秽，怎敢为他人祈求呢？你们怎敢为他们献祭呢？\par

259. 肋未人的职责并非轻松，因为上主论到他们说：\BibleQuote{我要从以色列子民中选取肋未人，代替以色列子民中一切头胎的长子；肋未人应归我所有。因为凡在埃及地头胎生的，我都祝圣了他们，使之归我。} \BibleRef{户 3:12} 我们知道肋未人并不被算在其餘人当中，而是被置于众人之上，因为他们是从众人中被拣选出来的，且如属于上主的初果和首生一样被祝圣，因为誓愿的酬报和赎罪之祭献都是由他们呈上的。\BibleQuote{你们不可将他们}，\Emph{他说}，\BibleQuote{算在以色列子民中，却要派肋未人管理约柜的帐幕，以及其中的一切器具和所有物件。他们要搬运帐幕和其中的一切器具，并在帐幕中服务，在帐幕四周扎营。帐幕往前行时，肋未人要拆卸；扎营时，肋未人要支搭。凡外人走近的，应当处死。} \BibleRef{户 1:49}\par

260. 因此，你，从以色列全体子民中被拣选出来，被视为神圣祭品中的初果，被设立在圣幕上，以在圣洁和信德的营地中守望；凡外人走近的，必定处死。你被安置在那里，看守\OriginalTerm{约柜}{ark of the covenant}。并非所有人都能窥见奥迹的深处，因为这些奥迹对肋未人隐藏起来，以免那不该看见的人看见，那不能服务的人将其举起。的确，\Person{梅瑟}看见了\OriginalTerm{圣神的割损}{circumcision of the Spirit}，却将其遮蔽，只将割损表现为外在的标记。他看见了真诚与真理的\OriginalTerm{无酵饼}{unleavened bread}；他看见了主的苦难，却将真理的无酵饼隐藏于物质的无酵饼内，将主的苦难隐藏于羔羊或牛犊的祭献中。优秀的肋未人始终以自己信德的保障守护着托付给他们的奥迹，而你却轻视所托付给你的？首先，你应窥见天主的深奥事物，这需要智慧。其次，你必须为百姓守望；这需要义德。你必须保卫营地、守护圣幕，这需要勇德。你必须显出节制和清醒，这需要节德。\par

261. 这些主要德行，外界人士虽已认识到，但他们认为建立于社会上的秩序，高于建立于智慧上的秩序；然而智慧是根基，义德是建筑，若没有根基，建筑便不能直立。根基就是基督。\BibleRef{格前 3:11}\par

262. 首先有信德，它是智慧的标记，正如\Person{撒罗满}继承其父所说：\BibleQuote{敬畏上主是智慧的开端。} 法律也说：\BibleQuote{你应当爱你的主天主，并爱你的近人。} \BibleRef{申 6:5} 向整个人类施予仁慈并尽义务，是件高尚的事。然而，最适宜的莫过于你将你所拥有的最宝贵之物，即你的心灵，献给天主，因为你没有比这更好的了。当你向你的造物主偿还了债务，你便可以为人劳苦，向他们施仁慈，提供帮助；然后你可以用金钱，或藉某项义务，或你所承担的某项圣职中的服务，来帮助穷乏的人：以金钱支持他；替他还债，以释放被捆缚者；承担一项义务，为他人保管他存放并担心失去的寄托之物。\par

263. 因此，保管并归还原主托付给我们的物件，是一项义务。然而，时而时移世易，以致归还所收之物不再是义务。例如，当一个人以公敌的身份索回金钱，意图用来对抗自己的祖国，并将财富献给蛮族。或者，若你当归还时，旁边有人要用暴力从他手中强夺。如果你将金钱归还给一个不能保管的狂怒疯子；如果你将曾存放在你这里的一把剑交还给一个疯子，他很可能用来自杀，偿还这笔债岂不是违背义务？明知是盗贼偷来的财物却接收，致使失去财物的人遭受欺骗，这岂不是违背义务？\par

264. 有时，履行诺言或持守誓言亦属违背义务。就如\Person{黑落德}，他曾发誓，无论\Person{黑落狄雅}的女儿求什么，凡她所求的，他都要给，为不愿食言，竟致容许\Person{若翰}丧命。更遑论\Person{依弗大}，他献上自己的女儿作为祭献，因为她是他胜利返家时第一个迎接他的人；这样他实践了他所许的愿：他要把第一个迎接他的人献给天主。他宁可根本不发此愿，也胜过以女儿的死亡来履行。\par

265. 你们并非不知道关注此事的重要性。因此，一位\OriginalTerm{肋未人}{Levite}被拣选以守卫\OriginalTerm{圣所}{sanctuary}，他绝不会在劝言上有所缺失，不会背弃信仰，不会惧怕死亡，也不会做出任何逾矩之事，以至在他的整个外表上都能显出其诚切之心。因为他不仅应约束自己的灵魂，甚至应约束自己的双眼，免得任何意外不幸使他的面颊生出羞红。因为\BibleQuote{凡注视妇女而贪恋她的，已在心里奸淫了她。}\BibleRef{玛 5:28}由此可知，奸淫不仅由实际做出丑恶之事而犯下，甚至也由炽热目光的欲望而犯下。\par

266. 这似乎崇高且略显严苛，但处于崇高的职务中，这并非不合宜。因为\OriginalTerm{肋未人}{Levites}的\OriginalTerm{恩宠}{grace}是如此之盛，以至于\Person{梅瑟}在以下祝福中论及他们说：\BibleQuote{将他的男子赐给肋未，将他的忠信者赐给肋未，将他的产业之分赐给肋未，将他的真理赐给那些在试探中试探了他的圣者，并在争辩的水旁咒骂了他。他对自己的父母说：我不认识你；不认识自己的弟兄，也不承认自己的子女。他守护了祢的话，持守了祢的证言。}\BibleRef{申 33:8}\par

267. 因此，他们就是祂的男子，祂的忠信者，他们心里没有诡诈，不隐藏任何奸诈，反而谨守祂的话，并在心里默思这些话语，正如玛利亚所默思的那样；\BibleRef{路 2:19}他们不认识自己的父母，免得将父母置于责任之先；他们憎恨那侵犯贞洁的人，并为纯洁所受的伤害施行报复；他们知道履行自己责任的时间，也知道哪个责任更大，哪个较小，以及每一项职责适宜于何种场合。在这一切事上，他们只遵循那有德行的事。并且，当存在两个有德行的责任时，他们以为那更具德行的应置于优先。这些人确实是适当地蒙受祝福的。\par

268. 若有人宣扬主正义的工程，并向祂献香，那么：\BibleQuote{主啊，求祢祝福他的力量；悦纳他手中的工作，}\BibleRef{申 33:11}好使他能在那位生活并统治，直到永远的那一位面前，获得先知性祝福的恩宠。阿们。\par

\SourceRecord{Translated by H. de Romestin, E. de Romestin and H.T.F. Duckworth.；Nicene and Post-Nicene Fathers, Second Series；Vol. 10.；Edited by Philip Schaff and Henry Wace.；Buffalo, NY: Christian Literature Publishing Co.,；1896.；<http://www.newadvent.org/fathers/34011.htm>.}

% structure: p0001 source=h1 target=section reason=page-title

\section{论圣职人员的职责（第二卷）}

% structure: p0002 source=h2 target=subsection reason=relative-heading-rank; base=h2

\subsection{第一章}

\Emph{基督徒若轻视世间的荣耀与众人的赞赏，在所作所为上惟愿讨天主喜爱，便能借着修德行善获得人生的幸福。}\par

1. 在第一卷中，我们论述了我们认为适合德行生活的职责；的确，没有人曾怀疑，那\WorkTitle{圣经}称之为永生的真福生活，正是赖此而获得。德行生活的光辉如此浩大，以致平安的良心和宁静的无罪便成就了幸福的生活。正如旭日东升，遮掩了月轮和星光，同样，当德行生活的光辉在真正纯洁的荣耀中闪耀时，便将那一切按肉欲视为美好、或在世人眼中当作伟大尊贵的事物，尽都掩蔽了。\par

2. 那生活显然是有福的，它不以外在的评价为准，而是以自身的内在感受作为自身的审断，自我认知。它绝不需要世人的赞誉作为赏报，亦无惧惩罚。因此，它越不追求荣耀，便越超越荣耀。因为对那些寻求荣耀的人而言，现世事物的酬报不过是未来之事的影子，且是永生的阻碍，正如\WorkTitle{圣经}上所记载的：\BibleQuote{我实在告诉你们，他们已获得了他们的赏报。}\BibleRef{玛 6:2}这是指那等人，他们仿佛以号筒之声，愿向举世显扬自己对穷人的慷慨施舍。同样，那仅为了外表的炫耀而行的禁食也是如此。祂说：\BibleQuote{他们已获得了他们的赏报。}\par

3. 因此，施行仁慈与暗中禁食，乃属德行生活之事；这样一来，你便似乎只寻求你的天主的赏报，而非来自人的赏报。因为凡向人寻求赏报的，已获得了他的赏报；但向天主寻求的，却得永生，这永生唯永生之主能赐予，正如\WorkTitle{圣经}上说：\BibleQuote{我实在告诉你，今天你就要与我一同在乐园里。}\BibleRef{路 23:43}因此，\WorkTitle{圣经}明白地称那有福之生命为永生。它并没有被交付给人的观念去评断，而是被托付于天主的审断。\par

% structure: p0007 source=h2 target=subsection reason=relative-heading-rank; base=h2

\subsection{第二章}

\Emph{哲学家们关于幸福的不同见解。他首先从\WorkTitle{福音}证明，幸福基于对天主的认识及对善工的追求；接着，为避免人们认为这观念取自哲学家，他更引用了先知们的见证作为证据。}\par

4. 哲学家们使幸福的生活或依赖于无痛苦（如\Person{希罗尼穆斯}），或依赖于知识（如\Person{赫里路斯}）。因为\Person{赫里路斯}听到\Person{亚里士多德}与\Person{泰奥弗拉斯托斯}极力称赞知识，便唯独以知识为至善，尽管他们实际上称赞知识为一种善，而非唯一的善；另有人，如\Person{伊壁鸠鲁}，则称快乐为至善；还有人，如\Person{卡利福}，及随后的\Person{狄奥多罗斯}，则以这样的方式理解：使德行生活与某物结合，一个与快乐结合，另一个与无痛苦结合，因为幸福生活离了它便不能存在。斯多葛派\Person{芝诺}则认为，唯有德行生活才是至高而独一的善。但\Person{亚里士多德}、\Person{泰奥弗拉斯托斯}及其他逍遥派主张，幸福生活在于德行，即在德行生活，但其幸福却由身体的优越及其他外在善物所成全。\par

5. 然而，\WorkTitle{圣经}却说，永生在于认识天主的事理和善工的果实。\WorkTitle{福音}为这两项陈述作证。因为主耶稣论及知识时这样说：\BibleQuote{永生就是认识你，唯一的真天主，和你所派遣来的耶稣基督。}\BibleRef{若 17:3}论及善工，祂这样回答：\BibleQuote{凡为我的名，舍弃了房屋，或兄弟，或姊妹，或父亲，或母亲，或妻子，或儿女，或田地的，必要领取百倍的赏报，并承受永生。}\BibleRef{玛 19:29}\par

6. 谁也不可认为这是晚近才说的话，并以为在\WorkTitle{福音}提及之前，哲学家们早已讲过了。因为哲学家们，意即\Person{亚里士多德}、\Person{泰奥弗拉斯托斯}，以及\Person{芝诺}和\Person{希罗尼穆斯}，固然生活在\WorkTitle{福音}时代之前；但他们却迟于先知们。他们更应思量，早在哲学家的名字被听闻之前多久，这两者似乎已借着圣\Person{达味}的口公开宣告出来了；因为\WorkTitle{圣经}上记载说：\BibleQuote{上主，你所教导，并以你的法律训诲的人，是有福的。}\BibleRef{咏 94:12}我们在另一处也看到：\BibleQuote{敬畏上主的人，是有福的，他必极其喜爱祂的诫命。}\BibleRef{咏 112:1}关于知识，我们已经证明了我们的论点；先知指出，知识的赏报就是永恒的果实，并补充说，在敬畏上主、或受祂法律的教导、并极其喜爱上主诫命的人家中，\BibleQuote{有光荣与财富；他的仁义永存不替。}\BibleRef{咏 112:3}他在同一篇圣咏中，更进一步论及善工，说善工为正直的人赢得永生的赏赐。他这样说：\BibleQuote{乐善好施的人，是有福的；他必妥善处理自己的事务，他永不动摇，义人必永受怀念。}\BibleRef{咏 112:5-6}又说：\BibleQuote{他广施钱财，周济穷人；他的仁义永存。}\BibleRef{咏 112:9}\par

7. 信德，则持有[应许的]永生，因为这是良好的基础。善工同样如此，因为一个正直的人由他的言语和行为受检验。因为若一个人总是忙于说话却行事迟缓，他的行为就显明他的知识是何等无用：此外，知道当行之事却不去行所学应行的，更为糟糕。另一方面，忙于善工而心中不忠信，就如同想要在坏的基础上建造一座美丽高大的圆顶。建得越高，倒塌就越烈；因为没有信德的保障，善工无法站立。港口里不可靠的锚地会凿穿船只，沙质底部很快塌陷，无法承受建在其上的建筑物的重量。因此，在德行圆满，言语与行为合理一致的地方，才会找到圆满的赏报。\par

% structure: p0013 source=h2 target=subsection reason=relative-heading-rank; base=h2

\subsection{第3章}

\Emph{真福的定义根据圣经被考虑和证明。它不能被外在的好运增强，也不能被不幸削弱。}\par

8. 那么，正如知识孤立无援时，根据哲学家们多余的讨论，不是被当作毫无价值而弃置，就是仅被视为不完美的理念，我们就来注意神圣的圣经多么清晰地解释了一件事，关于此事我们见到哲学家们持有如此多复杂纷乱的想法。因为圣经声明，除了有德行的，一无所善；并宣布德行在任何境遇中都是有福的，它从不因身体的或其他外在的好运而增强，也不因逆境而削弱。没有哪种状态能如此有福，如同人免于罪，充满纯洁无罪，并满被天主的恩宠。因为经上记载：\BibleQuote{凡不随从恶人的计谋，不插足于罪人的道路，不坐于讥讽者的席位，而专心爱好主法律的，是有福的。} 又说：\BibleQuote{行为无玷，遵行主法律的，是有福的。}\par

9. 那么，纯洁无罪与知识使人有福。我们也已经注意过，永生的真福是善工的赏报。接下来要说明，当轻看了逸乐的庇护或痛苦的恐惧（前者被视为可怜而柔弱，后者被视为不丈夫而软弱），那么有福的生命就能在痛苦中兴起。当我们读到以下经文时，这很容易说明：\BibleQuote{几时人为了我而辱骂迫害你们，捏造一切坏话毁谤你们，你们是有福的。你们欢喜踊跃罢！因为你们在天上的赏报是丰厚的，因为在你们以前的先知，人也曾这样迫害过他们。}\BibleRef{玛 5:11} 又说：\BibleQuote{谁若愿意跟随我，该弃绝自己，背着自己的十字架来跟随我。}\BibleRef{玛 16:24}\par

% structure: p0017 source=h2 target=subsection reason=relative-heading-rank; base=h2

\subsection{第4章}

\Emph{这同一论证，即真福不因外在事物而减少或增加，由古人的例子得以说明。}\par

10. 因此，即使在痛苦与悲伤之中，也存在真福。这一切德行以其甘饴予以制止和约束，因为它在天然的助佑中充裕，既能安抚良心，又能增添恩宠。因为梅瑟并非小有福气，那时他被埃及人包围，被海困住，却因自己的功绩为自己和百姓找到了穿行水中的道路。\EditorialNote{出 14} 他何时曾比那一刻更勇敢，那时他身处极大危险，却不放弃安全之望，反而祈求凯旋？\par

11. 亚郎又如何呢？他何时曾认为自己比那一刻更有福，那时他站在生者与死者之间，以自己的临在阻止死亡从死者身体传至生者行列？\BibleRef{户 16:48} 对于少年达尼尔，我又该说什么呢？他如此明智，以致当他身处饥肠辘辘的猛狮之中时，竟丝毫不因野兽的凶猛而恐惧。他如此无惧，竟能进食，并不担心自己的榜样会激起野兽吞食自己。\par

12. 因此，在痛苦中存有一种德行，它能彰显善良心的甘饴，因而证明痛苦并不减损德行的愉悦。那么，就像痛苦不会导致德行的真福损失一样，肉身的逸乐以及利益所带来的享受也不能为其增添什么。关于此事，宗徒说得好：\BibleQuote{凡以前对我有利益的事，我如今为了基督，都看作是损失，} 并且他补充说：\BibleQuote{不但如此，而且我将一切都看作损失，因为我只以认识我主基督耶稣为至宝；为了他，我自愿损失一切，拿这一切当废物，为赚得基督。}\BibleRef{斐 3:7}\par

13. 梅瑟也同样，他认为埃及的珍宝是自己的损失，并因此在生活中彰显了主十字架的屈辱。当他拥有丰厚金钱时，并不富有；后来他缺乏食物时，也并不贫穷，除非也许有人认为，当日用的食物在旷野中缺乏，他和他的百姓就不那么幸福。然而玛纳，即天使的食粮——确实没有人敢否认那是极大美善与真福的标记——却从天上赐给了他；每日降下的肉食也足够喂养全体众人。\BibleRef{出 16:13}\par

14. 为那位圣善的厄里亚，如果他去寻找食粮，食粮也会缺乏；但似乎它并未缺乏，因为他没有去寻找。于是，靠着乌鸦每日的服侍，早晨给他带来饼，晚上带来肉。他是否因为对自己贫穷而真福稍减呢？绝非如此。相反，他更为有福，因为他对天主富足。为他人富足胜过为自己富足。他正是如此，因为在饥荒时期，他向一位寡妇求饼，打算偿还，以致面缸里的面三年零六个月没有短缺，油瓶里的油也足以供那贫困的寡妇在此期间每日使用。伯多禄愿意在那里看见他们，这是合理的。他们与基督在光荣中显于山上，也是合理的，\BibleRef{玛 17:3} 因为当他富有时，却成为贫穷的。\par

15. 因此，财富无助于度一种真福的生活，这一事实，主在\BibleBook{路加}{福音}中清楚地指出，说：\BibleQuote{你们贫穷的是有福的，因为天主的国是你们的。你们现今饥饿的是有福的，因为你们将得饱饫。你们现今哭泣的是有福的，因为你们将要欢笑。}\BibleRef{路 6:20}这样，极清楚地指出，贫穷、饥饿和痛苦，这些被认为是不幸的事，不但不妨碍真福的生活，反而实在有助于它。\par

% structure: p0025 source=h2 target=subsection reason=relative-heading-rank; base=h2

\subsection{第五章}

\Emph{那些通常被人视作善的事物，大多是对真福生活的障碍；而被人视作恶的事物，却是德行生长的材料。有关真福之事，由其他例子可见。}\par

16. 然而，那些看似美好的事物，如财富、丰足、没有痛苦的欢乐，实则是真福果实的障碍，正如主亲口所说的话清楚指出的：\BibleQuote{你们富有的是有祸的，因为你们已经获得了你们的安慰！你们现今饱饫的是有祸的，因为你们将要饥饿。你们现今欢笑的是有祸的，因为你们将要哀恸哭泣！}\BibleRef{路 6:24}因此，肉体或外在的善，不仅无助于达到真福的生活，反而阻碍着它。\par

17. 因此，\Person{纳波特}是有福的，尽管他被富人用石头砸死；他软弱贫穷，与王室的资源相比，却在志向和信仰上是富足的；他如此富有，竟不肯以祖传的葡萄园来换取国王的金钱；正因如此，他是全德的，因为他用自己的血捍卫了祖先的权利。同样，\Person{阿哈布}在自己看来也是悲惨的，因为他将穷人处死，以便自己占有他的葡萄园。\par

18. 完全可以确定的是，德行是唯一的、至高的善；唯有德行丰盈结果，带来真福生活；真福生活，通过它赢得永生，并不依靠外在或肉体的好处，而只依靠德行。真福生活是当下的果实，而永生是未来的希望。\par

19. 然而，有些人认为，在这具软弱而脆弱的身体内，不可能有真福的生活。因为在这身体内，人必须忍受痛苦和悲伤，必须哭泣，必须生病。因此我也可以说，真福的生活在于肉体的欢乐，而不在于智慧的高峰、良心的甘饴或德行的高尚。处于痛苦之中并不是有福的，但战胜它，不因短暂的痛苦而丧胆，才是有福的。\par

20. 试想，来了一些被认为因其导致的悲痛而可怕的事，比如失明、流亡、饥饿、女儿受辱、子女死亡。谁会否认\Person{依撒格}是有福的呢？他年老时看不见，却以他的祝福颁赐了祝福\BibleRef{创 27:28}。\Person{雅各伯}岂非有福？他离开父家，作为牧人受雇流亡\BibleRef{创 31:41}，为女儿受辱的贞洁而哀恸\BibleRef{创 34:5}，遭受饥饿\BibleRef{创 42:2}。那些以其诚信而蒙\Person{天主}作证的人，岂非有福？正如经上记载：\BibleQuote{\Person{亚巴郎}的天主，\Person{依撒格}的天主，\Person{雅各伯}的天主}\BibleRef{出 3:6}。奴役是可悲的事，但\Person{若瑟}并不悲惨；不，显然他是有福的，因为他在为奴时克制了主母的淫欲\BibleRef{创 39:7}。关于圣\Person{达味}，他哀悼三个儿子的死亡，更有甚者，哀悼他女儿的乱伦，我又能说什么呢？他怎能是无福的呢？真福的创立者正是从他而出，这位创立者使许多人成为有福的。因为：\BibleQuote{那些没有看见而相信的，才是有福的}\BibleRef{若 20:29}。所有这些人感受到自己的软弱，却勇敢地战胜了它。我们能想象比圣\Person{约伯}更悲惨的吗？无论是他的房屋被焚，十个儿子瞬间死亡，还是他身体的苦痛。难道他比没有经历这一切，从而真正表明自己经得起考验的人更少福份吗？\par

21. 的确，在这些痛苦中有某种苦涩，意志的力量也无法隐藏这痛苦。我不会因为近岸处水浅就否认海洋的深邃，也不会因为时而乌云密布就否认天空的清澈，不会因为某些地方贫瘠就否认大地的丰产，也不会因为庄稼有时夹杂着野燕麦就否认庄稼的丰盈。同样，也要确信：一个快乐良心的收获可能混有一些悲痛的苦涩。在整个真福生活的禾捆中，倘若偶尔不幸或苦涩潜入，岂不像野燕麦被隐藏，或毒麦的苦味被谷物的香甜所掩盖吗？不过，让我们现在继续我们的主题。\par

% structure: p0033 source=h2 target=subsection reason=relative-heading-rank; base=h2

\subsection{第六章}

\Emph{论有用之事：并非指有利可图的事，而是指正义而有德的事。它可于损失中发现，分为对身体有用的，和对虔敬有用的。}\par

22. 在第一卷书中，我们曾那样分类，将美善的与合宜的置于首位；因为一切责任皆源于此。在其次，我们放置有用之事。正如我们一开始所说，美善的与合宜的之间存在着区别——人们领会它比解释它更容易——同样，当我们思考有用之事时，也必须认真思考什么更有用。\par

23. 但我们不以金钱收益的价值来衡量用处，而以获得虔敬来衡量，正如宗徒所说：\BibleQuote{虔敬在各方面都有益，因为有今生与来生的应许。}\BibleRef{弟前 4:8}因此，在圣经中，如果我们仔细考察，就常会发现：凡有德之事即被称为有用之事：\BibleQuote{凡事我都可行，但不全有益} [有用的]\BibleRef{格前 6:12}在此之前他谈论的是恶习，因此他的意思是：犯罪是允许的，但并不合宜。罪在自己权下，却非有德之事。奢华度日是容易的，但并非正理。因为食物不是为服事天主，而是为肚腹。\par

24. 因此，既然有益的就是正义的，那么服事那救赎了我们的基督，便是正义的。那些为祂的名的缘故舍身赴死的人，也是正义的；逃避死亡的人，便是不义的。圣经论到这些人说：\Quote{我流鲜血，有何利益？}意即：我的正义有何进步？因此他们也说：\Quote{我们要束缚义人，因为他对我们无用，}意即：他是不义的，因为他指责我们，谴责并叱责我们。这也可指不虔敬者的贪婪，那贪婪与背叛极其相似；正如我们在叛徒\Person{犹达斯}的例子中所读到的，他因贪图利益和渴求金钱，将头伸进了背叛的圈套，就跌倒了。\par

25. 因此，我们要谈论那充满德行的有益之事，正如宗徒本人已明言的：\BibleQuote{我说这话，是为你们的益处，不是要设下圈套陷害你们，而是为指给你们一件合宜的事。}\BibleRef{格前 7:35} 那么显然，凡有德行的便是有益的，凡有益的便是有德行的；同样，凡有益的便是正义的，凡正义的便是有益的。我可以说这话，因为我不是对因贪图利益而贪婪的商人说话，而是对我的孩子们说话。我正在谈论我切望铭刻并传授给你们的那些本分，你们是我为服事主所拣选的；好使那些已借习惯与训练深植并固定在你们心意和品格中的事，如今能借着讲解和教导，更充分地展现在你们面前。\par

26. 因此，我既将谈论有益之事，便愿采用先知的这句话：\Quote{使我的心倾向祢的法度，不倾向贪财，}免得\Quote{有益}一词的发音，激起我们对金钱的贪欲。的确，有人这样读：\Quote{使我的心倾向祢的法度，不倾向有益的事，}那是指那种时刻伺机从生意中谋利、因人的习俗而偏向并扭曲为追求金钱的所谓有益。因为通常大多数人只以能赚钱的为有益，但我们所谈论的有益，却是在世上的损失中寻求，\BibleQuote{为获得基督，}\BibleRef{斐 3:8} 祂的利益就是\BibleQuote{虔敬与知足。}\BibleRef{弟前 6:6} 我们借以达至虔敬的利益也何其大，这虔敬在天主前是富足的，不在于转瞬即逝的财富，而在于永恒的恩赐，其中没有不可靠的考验，只有恒久不息的恩宠。\par

27. 因此，按照宗徒的划分，有一种与身体相关的有益，另有一种与虔敬相关的有益：\BibleQuote{身体的操练益处有限，但虔敬在各方面都有益。}\BibleRef{弟前 4:8} 有什么像正直那样有德行？有什么像保持身体无瑕无玷、纯洁不受沾染那样合宜？再者，有什么像寡妇对已故丈夫信守誓约那样合宜？有什么比这更有益，藉此可达至天国？因为\BibleQuote{有些人为天国的缘故自阉了。}\BibleRef{玛 19:12}\par

% structure: p0041 source=h2 target=subsection reason=relative-heading-rank; base=h2

\subsection{第七章。}

\Emph{有益与有德是同一回事；没有什么比爱更有益，爱由温和、礼貌、善良、正义及其他德行获得，正如我们从\Person{梅瑟}和\Person{达味}的历史中得知的。最后，信赖出自爱，而爱又出自信赖。}\par

28. 因此，不仅德行与有益之间有着密切的关联，而且同一件事既是有益的又是有德的。因此，那位愿意向所有人敞开天国的那一位，不求自己的益处，而求众人的益处。所以我们必须有一定的次序，从习惯或普通的行为逐步进到更卓越的行为，以便借许多范例显示有益之处的进步。\par

29. 首先我们要知道，没有比受人爱戴更有益的事，也没有比不为人所爱更无益的事；因为在我看来，受人憎恨简直是致命的，彻底是死亡的。我们谈论这事，是为了留心使人们对我们有良好的评价和看法，并借心灵的平静与灵魂的善良，努力在他人心中获得一席之地。因为良善为众人所喜悦和悦纳，没有什么比这更容易触动人的情感。如果再加上性情的温柔与乐意，以及发号施令时的适度、言语的彬彬有礼、言谈中的敬重、对话的随机应变和谦逊的优雅，这一切对增加爱的助益，实在难以尽述。\par

30. 我们读到的，不仅是在个人身上，就是在君王身上也如此：欣然且乐意的礼貌能产生什么效果，骄傲与夸大其辞又带来了什么损害，甚至使王国动摇，权势毁灭。若有人借劝告或服务，借履行其职务或职责的种种责任而赢得人民的爱戴，或者为了整个民族而身陷危险，毫无疑问，人民将要向他显示如此的爱，甚至将他的安全与福祉置于自己的之上。\par

31. \Person{梅瑟}从他的人民那里承受了多少指责！但当主因那些辱骂他的人要为他报仇时，他却常为人民献上自己，为从神圣的义怒中拯救他们。\BibleRef{出 32:32} 即使在遭受冤屈之后，他仍以何等温和的言语向人民说话！他安慰他们的劳苦，以他对未来的先知性宣告抚慰他们，以他的工作鼓励他们。虽然他时常与天主交谈，但他总是温和愉快地对人说话。他当之无愧被视为超乎众人之上。因为他们甚至不能仰视他的面容，\BibleRef{出 34:30} 并且不肯相信他的坟墓能被找到。\BibleRef{申 34:6} 他如此吸引了全体民众的心，以致他们因他的温和而爱他，更甚于因他的事迹而钦佩他。\par

32. 也有\Person{达味}追随他的脚步，他从众人中被拣选治理百姓。他是多么温柔和蔼，心灵谦卑，多么勤勉并乐于表达亲情。在他登基之前，就替众人献上自己。作为君王，他在战争中与众人平等，分担他们的劳苦。他作战勇敢，治理温和，受辱时忍耐，更愿承受委屈而不报复。他深得众人喜爱，以致虽然尚是青年，便被迫违背自己的意愿被选作他们的统治者，虽曾抗拒，仍被推上这一职责。年老时，百姓请求他不要亲临战场，因为他们全都宁愿为他冒险，也不愿他为他们经历危险。\par

33. 他在履行职责时，自由地将百姓与自己联结在一起；首先，在百姓分裂期间，他宁愿像流亡者一样住在\Place{赫贝龙}，而不在\Place{耶路撒冷}为王；其次，当他在仇敌身上表现出对勇武的喜爱时。他也认为，对那些曾与他为敌的人，应像对待自己的人一样施予正义。此外，他钦佩\Person{阿贝乃尔}——敌方最勇敢的战士，那时阿贝乃尔仍是他们的领袖，且正与他交战。当阿贝乃尔前来求和时，他并未轻视他，反而设宴款待。当阿贝乃尔被奸计杀害时，他为阿贝乃尔哀哭。他为他送葬，尊敬地举行葬礼，并表明自己渴望为谋杀复仇的诚意；因为他将这一责任嘱咐给儿子，交付给他，他更关心的是无辜人的死不可不报仇，而非任何人为自己的死哀哭。\par

34. 尤其对于一位君王而言，履行谦卑的职责以至于使自己与最卑微的人相似，这绝非小事。不殃及他人而求食，拒绝饮水，承认罪过，并为自己的人民舍命，这是高尚之举。达味正是这样做的，为使天主的义怒转向他自己，当他向毁灭天使献上自己，说：\BibleQuote{看，是我犯了罪，我为牧者行了恶，但这羊群作了什么？愿你的手打击我。}\par

35. 我还有什么可说呢？他对着图谋诡诈的人不开口，如同听不见，他认为不应还口，也不回答他们的辱骂。遭人恶言时，他祈祷；受人诅咒时，他祝福。他心地单纯而行，远离骄傲的人。他效法那些在世上无玷的人，哀哭己罪时将灰掺入食物，以泪水调和饮品。因此，众百姓都前来召请他。\Place{以色列}各支派来对他说：\BibleQuote{看，我们是你的亲骨肉。以前，连昨天和前天，当\Person{撒乌耳}为王治理时，是你率领\Place{以色列}出入；上主曾对你说：你应牧养我的百姓。}我还要说什么呢？关于他，上主已发出话说：\BibleQuote{我找到了合我心意的\Person{达味}。}谁曾像他那样心地圣善、行事正义，以承行天主的旨意？因他的缘故，他的子孙犯罪时蒙受宽恕，他的后嗣仍保有权利。\par

36. 看到他对朋友多么情深，谁不喜爱他呢？因为他真心爱他的朋友，所以他也以为自己的朋友同样地爱他。不仅如此，父母甚至将他置于自己子女之前，子女爱他甚于爱父母。因此，\Person{撒乌耳}非常愤怒，企图用矛刺死儿子\Person{约纳堂}，因为他认为，在约纳堂心目中，\Person{达味}的友谊比孝爱之情或父亲的权威更为重要。\par

37. 如果一个人对爱我们的人也报以爱，并证明自己爱他们不亚于自己被爱，特别是通过忠实友谊所给予的证明来表达，这将极大地增进彼此的爱。有什么比感恩更能赢得好感呢？有什么比爱那爱我们的人更自然呢？有什么比希望让我们想要他爱的人知道我们爱他这种情感，更深植于人心呢？智者说得对：\BibleQuote{为你的兄弟和朋友，可以失去金钱。}\BibleRef{德 29:10} 又说：\BibleQuote{我不以护卫朋友为耻，也不躲避他。}\BibleRef{德 22:31} 的确，\BibleBook{德训}{篇}中的话证实：朋友是生命与不死之药；\BibleRef{德 6:16} 然而，从来无人怀疑，我们最好的保障在于爱。正如宗徒所说：\BibleQuote{凡事包容，凡事相信，凡事盼望，凡事忍耐；爱永不止息。}\BibleRef{格前 13:7}\par

38. 因此，\Person{达味}未曾失败，因为他为众人所爱，并且他愿意臣民爱戴他而非畏惧他。恐惧只能提供暂时的保护，却不懂得如何持久守卫。所以恐惧一旦离去，大胆便常会潜入；因为恐惧不能强制出信任，只有爱才能唤起信任。\par

39. 所以，爱是给予我们赞誉的首要因素。因此，我们的见证在于众人的爱，这实为美善。随后信心便会产生，以致连外方人看到你如此深受众人爱戴，也不怕将自身交托于你的仁慈。同样，人经由信心走向爱，因此那对一两个人表现出诚信的人，便仿佛感召了众人的心灵，赢得了众人的善意。\par

% structure: p0055 source=h2 target=subsection reason=relative-heading-rank; base=h2

\subsection{第八章。}

\Emph{赢取善意莫过于给予劝告，其效果无与伦比；但除非劝告基于正义与明智，无人能够信赖。这两样德行在\Person{撒罗满}身上何其彰显，他那众所周知的判决可资证明。}\par

40. 因此，两件事——爱与信心——在举荐我们于他人时最为有效；如果你拥有这第三种品质，即许多人认为你身上值得钦佩、并认定为确实配受尊敬的事（事实上，即给予良策的能力），那么这也会有所助益。\par

41. 既然给予良善的建议是赢得他人好感的重要手段，那么在一切情况下都需要\OriginalTerm{智德}{prudence}与\OriginalTerm{义德}{justice}。大多数人正是寻求这些，因此人们立即信赖那些具备这些德行的人，因为这样的人能给予任何需要者有益且可信的建议。谁会将自己交托给一个他认为并不比自己更明智的人来求取建议呢？因此，被求问者必须比求问者更胜一筹。因为我们若认为某人无法使事情比我们自己所见更清楚，又为何要去请教他呢？\par

42. 然而，如果我们寻得一人，因其品格之刚健、心智之力量与感召力而超越众人，而且比他人更善于以身作则、富有经验；能消除当前的危难，预知将来的灾祸，指明近在咫尺的危险，能阐明事理，及时给予救助，不仅乐于给予建议，更乐于提供帮助——这样的人便值得信赖，以至于寻求建议者会说：\BibleQuote{即便因他而遭遇凶恶，我也情愿忍受。}\BibleRef{德 22:31}\par

43. 我们于是将自身的安全与名誉托付给这样的人，因为如前所述，他既是义的又有智的。义德使我们不必担心受骗，智德则免去我们对失误的疑惧。不过，按一般人常说的方式，我们更乐于信靠一个正义的人，而非一个有智慧的人。然而，按哲学家们的定义，拥有一项德行便也拥有其他德行，而智德若无义德便不能存在。我们在我们的经卷中也见到这样的论述，因为\Person{达味}说：\BibleQuote{义人矜怜且施与人。}论到义人所施与的，他在别处又说：\BibleQuote{良善的人矜怜施与，他必以明哲指教自己的言语。}\par

44. \Person{撒罗满}那高贵无比的判决岂不是充满了智慧与正义吗？让我们来看看是否如此。经上说：\BibleQuote{有两个妇人站在撒罗满王面前，一个妇人说：\InnerQuote{我主，请听！我与这妇人同住一房，在房内我生了一个孩子，第三天她也生了一个孩子；我们住在一起，房内除我们二人外再没有别人。夜间，这妇人的孩子死了，因为她睡觉时压了自己的孩子。她半夜起来，乘我熟睡时，从我身边抱去我的孩子，放在她怀里，把她死了的孩子放在我怀里。清晨我起来给自己的孩子喂奶，竟发现他死了。及至天亮我细看，发觉那并不是我的孩子。}另一个妇人却说：\InnerQuote{不，活的是我的孩子，死的才是你的孩子。}}\par

45. 这就是她们的争执，双方都试图将活的婴儿据为己有，并否认死的是自己的。于是国王命人取来一把剑，吩咐将婴儿劈成两半，一半给这妇人，一半给那妇人。那活孩子的亲生母亲，因被怜子之情所动，遂喊道：\BibleQuote{我主，请不要杀害孩子；就将活的孩子给她吧，切莫伤害孩子！}但另一个却说：\BibleQuote{这孩子既不要归我，也不要归她，劈开吧！}于是国王下令将婴儿交给那说\Quote{不要杀他，要把他给那女人}的妇人；正如经上所记，\BibleQuote{因为她的心肠恋慕她的儿子。}\par

46. 猜测天主的心意在他内，并非错误；因为什么能瞒过天主呢？有什么比深藏于内的证据更隐秘的呢？那睿智的国王却深入其中，仿佛要审量一位母亲的感受，并引出了如同母亲心声的言语。因为当一位母亲宁愿自己的孩子与别人同住，也不愿他当着母亲的面被杀时，她的母怀便被揭示了出来。\par

47. 因此，这既是智慧的标记，能辨别隐秘的心念，从隐藏的泉源中引出真理，并如同用圣神之剑刺透人的五内及心思灵魂；也是义德的表现，使那杀死自己孩子的妇人不能夺走他人的孩子，而让真正的母亲重新得回自己的孩子。事实上，圣经对此已作了宣告：经上说，\BibleQuote{全\Place{以色列}听说了国王所判决的案件，便都敬畏国王，因为他们看出在他内有天主的智慧，以执行裁判。}\Person{撒罗满}自己也求得了智慧，为能秉公审断，领受一颗明智的心。\par

% structure: p0065 source=h2 target=subsection reason=relative-heading-rank; base=h2

\subsection{第九章}

\Emph{尽管义德与智德是不可分的，但我们必须顾念一般人常有的观念，因为他们将不同的\OriginalTerm{枢德}{cardinal virtues}加以区分。}\par

48. 同样，按照更为古老的圣经，也很清楚：智慧若无义德便不能存在，因为二者中若有一者存在，另一者必同在。\Person{达尼尔}以何等的智慧，通过彻底的盘问揭露了诬告者的谎言，使那些诬告者无言以对！凭他们自己的口供定罪，这是智德的标记；将罪人交付惩罚，并拯救无辜者脱离惩罚，则是义德的标记。\par

49. 因此，智慧与义德之间有不可分离的联系；然而，一般说来，这单一特殊形式的德行被划分为多种。例如，\OriginalTerm{节德}{temperance}在于轻视享乐，\OriginalTerm{勇德}{fortitude}见于承受劳苦与危险，智德在于选择善事，懂得如何分辨有益与无益之事；义德则在于妥善守护他人的权利并保护自己的权利，从而使人各得其所。为迁就通常接受的观念，我们可以作此四分法；因此，我们虽偏离了那些如同出自密室的、为探求真理而进行的精妙哲学讨论，却遵循了通行的应用和寻常的意义。就让我们按此划分，回到我们的主题吧。\par

% structure: p0069 source=h2 target=subsection reason=relative-heading-rank; base=h2

\subsection{第十章}

\Emph{人更愿意将自己的安全托付给正义之人，而非仅明智之人。但每个人通常都会寻求兼具正义与明智品质的人。撒罗满为此提供了榜样。（舍巴女王论及他的话将得到解释。）达尼尔和若瑟也是如此。}\par

50. 我们将自己的事务托付给所能找到的最明智的人，并且比其他任何人都更愿意向他寻求建议。然而，一个正义之人的忠告居于首位，且往往比最智慧者的卓越才能更具分量：\BibleQuote{因为朋友的创伤比他人的亲吻更好。}\BibleRef{箴 27:6} 正因这是一个正义之人的判断，所以它也是一个智慧之人的结论：前者在于争议事情的结果，后者在于发明的敏捷。\par

51. 如果一个人兼具二者，他所给出的建议便会十分可靠，所有人都因所展现的智慧而钦佩，并因其中的正义而喜爱。因此，众人都渴望聆听那兼具这两种美德之人的智慧，正如地上所有的君王都渴望见到撒罗满的面，聆听他的智慧。甚至舍巴女王也来见他，用问题试探他。她前来，述说了心中所有的事，聆听了撒罗满的一切智慧，无一言向她隐藏。\par

52. 在这段你听见她所说的话中，人啊，你要学习，这位没有任何事能瞒过她的女人是谁，以及热爱真理的撒罗满没有不告诉她的事：\BibleQuote{我在本国所听到关于你的言行和你的智慧的传闻，确是真的。我原不相信那些话，直到我来，亲眼见了，才知道人们告诉我的还不到一半；你所增添的善事，远超过我在本国所听到的。你的妻妾有福，你的臣仆有福，他们能时常侍立在你面前，聆听你的智慧。}你要认清真正撒罗满的筵席，以及谁被安排坐在这筵席上；要明智地认清，并思考在哪个国度里，万民将要听见真正的智慧和正义的声名，并以何种眼睛看见祂，注目那些不可见的事。 \BibleQuote{因为看得见的只是暂时的，看不见的才是永远的。}\BibleRef{格后 4:18}\par

53. 有福的妇女不是指那些听天主的话而结出果实的人吗？\BibleQuote{那听天主的话而遵行的人，更是有福的！}\BibleRef{路 11:28} 又说：\BibleQuote{凡遵行天主旨意的，就是我的兄弟、姊妹和母亲。}\BibleRef{玛 12:50} 那些站在祂面前的有福的仆人，不是保禄吗？他曾说：\BibleQuote{我直到今日仍站着向一切人作证。}\BibleRef{宗 26:22} 或者是那在圣殿中等待以色列安慰的西默盎？\BibleRef{路 2:25} 他怎能请求离去呢？除非是因为他站在主面前，无法仅凭己力离去，而只能按天主的旨意。撒罗满被立在我们面前，仅仅是作为榜样，人们曾热切期待聆听他的智慧。\par

54. 若瑟虽然在监里，却也未曾免于被咨询不确定之事。他的建议有益于整个埃及，使它不觉七年饥荒，甚至能解救其他民族脱离可怕的饥饿。\par

55. 达尼尔虽身为俘虏，却被立为王的谋士之首。他通过自己的建议改善了当下，并预言了未来。\EditorialNote{达尼尔 第2章} 众人凡事都信任他，因为他常常解释事物，并表明他说的是真理。\par

% structure: p0077 source=h2 target=subsection reason=relative-heading-rank; base=h2

\subsection{第11章}

\Emph{第三个使人获得信任的因素，在梅瑟、达尼尔和若瑟身上表现得尤为显著。}\par

56. 但第三个要点似乎也在那些堪受敬佩的人身上得到注意，即如以若瑟、撒罗满和达尼尔为榜样。我能如何评说梅瑟呢？全以色列总是等待他的建议，\BibleRef{出 18:13} 他的生活使他们信任他的明智，并增加对他的尊敬。谁不信任梅瑟的建议呢？长老们将一切他们自认为超出自己理解和能力的事，都留给他决断。\par

57. 谁会拒绝达尼尔的建议呢？天主亲口论及他说：\BibleQuote{谁比达尼尔更有智慧？}\BibleRef{则 28:3} 对于天主赐予如此恩宠的人，人们怎能怀疑他们的心智？藉由梅瑟的建议，战争得以结束，并因他的功劳，食物从天降下，饮水从磐石流出。\par

58. 达尼尔的灵魂何其纯净，竟能软化蛮族的性情，驯服狮子！他是何等节制，灵魂与肉身的自制能力何等强大！他毫不逊色地成了众人钦佩的对象，而——世人都钦佩这一点——他虽享有君王的友谊，却不求黄金，也不认为赐予自己的荣耀比他的信德更宝贵。因为他甘愿为天主的法律承受危险，也不肯为讨人的欢心而改变初衷。\par

59. 还有，我几乎略过的若瑟，他的贞洁和正义，我该怎样评说呢？他一方面拒绝了主母的诱惑和酬报，另一方面轻看死亡，压抑恐惧，选择了监狱。谁不认为他是个合适的人选，能在私下事务中提出建议呢？他那丰饶的精神和肥沃的心智，以他丰富的建议和心灵，充实了时间的荒芜。\par

% structure: p0083 source=h2 target=subsection reason=relative-heading-rank; base=h2

\subsection{第12章}

\Emph{没有人会向沾染恶习的人寻求建议，也不会向一个脾气暴躁或难以相处的人寻求，而是会向我们在圣经中拥有典范的人寻求。}\par

60. 因此，我们注意到，在寻求建议时，生活的正直、德行的卓越、乐善好施的习惯以及和善性情的魅力，都具有极大的分量。谁会在泥泞中寻找泉源？谁愿意饮用浑浊的水？因此，在奢华、放纵和各种恶习聚集的地方，谁会认为应当从那里汲取呢？谁不鄙弃污秽的生活？谁又会认为一个在自己的生活中无用的人，能对别人的事有益呢？再者，谁会不避开一个邪恶、心怀恶意、辱骂他人、时刻准备害人的人呢？谁不极力避开他呢？\par

61. 一个人即使最有资格提供最好的建议，却难以接近，谁会去找他呢？他就像水源被截断的泉源。如果一个人拒绝给予建议，拥有智慧又有什么益处呢？如果一个人断绝了提供建议的机会，就等于堵住了泉源，使它不再流向他人，也不再对自己有任何好处。\par

62. 我们可以很恰当地将这话用在那拥有智德，却以邪恶生活的污秽玷污了它，从而在源头污染了水的人身上。他的生活证实了他的精神堕落。一个人若看到另一个人品行邪恶，怎能断定他的建议是好的呢？如果我打算将自己托付给他，他就应当比我更优秀。难道我要认为一个从不为自己采纳建议的人适合给我建议吗？或者，难道我要相信，当他的心思充满享乐，被情欲制服，做贪吝的奴隶，被贪婪刺激，被恐惧吓倒时，他会有时间给我建议，却自己没有时间？这里哪有安静，又哪有建议的余地呢？\par

63. 那位我必须钦佩仰望的谋士，那位仁慈的主赐给我们祖先的，撇弃了一切令人反感的事。他应当成为追随他的人，因为这样的人能给予建议，并保护他人的智德免于邪恶；因为任何污秽都不能与之混杂。\par

% structure: p0089 source=h2 target=subsection reason=relative-heading-rank; base=h2

\subsection{第十三章}

\Emph{天主所作的见证显明了智慧之美。由此，他继续证明智慧与其他德行的联系。}\par

64. 有谁愿意自己面容美丽，却又因野兽般的身躯和可怕的利爪而破坏其魅力呢？其实德行的形态是如此奇妙而荣耀，尤其是智慧之美，正如整部圣经告诉我们的。因为智慧比太阳更辉煌，与星辰相比，远超任何星座。黑夜在其行列中夺去它们的光芒，但邪恶无法胜过智慧。\BibleRef{智 7:29}\par

65. 我们已经谈论了它的美，并以圣经的见证作了证明。剩下的，就是根据圣经的权威\BibleRef{智 7:22}来表明，智慧与邪恶之间不可能有任何相通，却与其他德行有着不可分离的结合。\BibleQuote{她拥有聪敏、无玷、稳固、圣洁、好善、敏捷的精神，从不禁止施恩，和蔼、坚定、无忧、全能、洞彻万有。}圣经又说：\BibleRef{智 8:7} \BibleQuote{她教授节制、正义和勇德。}\par

% structure: p0093 source=h2 target=subsection reason=relative-heading-rank; base=h2

\subsection{第十四章}

\Emph{智德与一切德行相结合，尤其与轻看财富相结合。}\par

66. 因此，智德贯穿万有，与一切美善相通。因为她若没有正义，怎能给出好的建议呢？这样她就能以恒常自持，不畏惧死亡，不被任何惊惶、任何恐惧所阻，也不认为有任何谄媚能使自己偏离真理是正当的，也不逃避放逐，知道世界就是智者的家乡。她不畏惧匮乏，因为她知道智者一无所缺，因为整个财富的世界都属于他。有谁比那不知如何为金钱的念头而激动，轻看财富，好似从某个高耸的视点俯视世人欲望的人更伟大呢？人们认为这样的人超出常人：\BibleQuote{他是谁？我们要赞美他，因为他在一生中行了奇妙的事。}\BibleRef{德 31:9} 的确，鄙视财富的人应当受人钦佩，因为大多数人甚至把财富置于自己的安全之上。\par

67. 节俭的规矩和自制的权威适合所有人，尤其适合那地位尊高的人；这样，对财宝的贪恋就不会抓住这样的人，而治理自由人的人也绝不会成为金钱的奴隶。更合宜的是，他的灵魂应超乎财宝之上，并甘愿服事朋友。因为谦逊会增加人对他的尊重。为人首领，不与\Place{叙利亚}商人和\Place{基肋阿得}商人一样贪图不义之财，不把一切幸福的希望寄托于金钱，不像雇工那样计算每日收益和积蓄，这是值得称赞和正确的。\par

% structure: p0097 source=h2 target=subsection reason=relative-heading-rank; base=h2

\subsection{第十五章}

\Emph{论慷慨。应主要对谁施予慷慨，以及资产微薄的人如何通过提供服务和建议来表现慷慨。}\par

68. 但是，如果灵魂摆脱这种缺点是值得称赞的，那么，以慷慨赢得人民的爱戴，既不对不配的人过于大方，也不对急需的人过于吝啬，该是多么光荣啊。\par

69. 慷慨有多种方式。我们不仅可以从自己的日常供应中分发和施予食物给那些需要的人，以维持他们的生命；我们还可以向那些羞于公开表露自己穷困的人提供建议和帮助，只要穷人们的公共供应没有枯竭。我现在说的是那些被委派担任某项职务的人。如果他是司铎或施赈员，就应当向主教报告这些人，不要隐瞒任何他所知道的需要帮助的人，或那些丧失了财富、现在陷于穷困的人；特别是如果他们陷入困境并非因为年轻时的挥霍，而是由于他人的偷窃，或并非因自己的过错而失去祖业，以致现在无法糊口。\par

慷慨的最高形式是赎回俘虏，将他们从敌人手中救出，使人免于死亡，尤其使妇女免于耻辱，使子女重归父母，父母重得子女，并使公民回归祖国。当\Place{色雷斯}和\Place{伊利里亚}遭到如此可怕的蹂躏时，这一点便得到了公认。当时全世界有多少俘虏被贩卖！若将他们集合起来，其数量会超过整整一个行省。然而，有些人却想把教会已赎回的人重新送回为奴。他们自己比奴隶制度更冷酷，竟对他人的慈悲侧目而视。他们（自己声称）若身为奴隶，便会甘愿为奴；若曾被出卖，也不会拒绝奴隶的服役。他们想要取消他人的自由，尽管他们不能取消自己为奴的身份，除非买主愿意收回赎价，如此，奴隶身份并非简单免除，而是被赎回了。\par

因此，慷慨的一个特质便是赎回俘虏，尤其从那些毫无人性火花、除非被贪欲用来勒索赎金否则绝无怜悯的蛮族敌人手中赎回。代他人偿债也是一件大事，倘若债务人无力偿还且被逼无奈，而该款项本是合法应偿、仅因贫穷而拖欠。同样，抚养子女、照顾孤儿，也是极大慷慨的标志。\par

还有的人将失去父母的少女嫁出，以保全其贞洁，不仅以美好祝愿相助，还赠予金钱。宗徒还教导了另一种慷慨：\BibleQuote{若有信友家中有寡妇，就当自行救济她们，不可累及教会，好使教会能救济那些真正守寡的妇女。} \BibleRef{弟前 5:16}\par

因此，这类慷慨是有用的，但并不对所有人通用。因为有许多善良的人资财微薄，自己的用度也仅够少量需求，无力济助减轻他人的贫困。然而，他们手边还有另一种善行，借此可以帮助更穷的人。因为慷慨有两种：一种提供实际援助，即以金钱；另一种则致力于积极的帮助，后者往往更加崇高、更加高贵。\par

\Person{亚巴郎}以得胜的武力夺回被掳的女婿\BibleRef{创 14:16}，比用金钱赎取他要伟大得多！圣\Person{若瑟}以他的谋划协助\Person{法郎}王备粮防饥，这比献给他金钱更加有益！因为金钱无法买回任何一个国家的丰产；而他却以自己的远见，使\Place{埃及}全地免于饥荒五年\BibleRef{创 41:53}。\par

金钱易耗；谋划却无穷竭。后者经恒常运用而愈发强劲。金钱渐少，转眼耗尽，甚至令善行本身受挫；因此，想要济助的人愈多，能获得帮助的反而愈少；而且往往当我们认为该当施与时，却已力不从心。然而，论到提供建议和积极帮助，开支的对象愈多，资源似乎愈发丰沛，且愈发回流其本源。明达智慧的丰沛之流常回注自身，它流向的范围愈广，所余留的便愈发具有活力。\par

% structure: p0107 source=h2 target=subsection reason=relative-heading-rank; base=h2

\subsection{第十六章}

\Emph{在施舍时，必须注意适度，不可将资财花费在无价值的人身上，而应当用于更值得的人。然而，施舍时也不可过于吝啬和犹豫。应效法圣\Person{若瑟}的榜样，圣经对他的明达大加赞赏。}\par

因此，显而易见，我们的慷慨必须持守适度，以免礼物变得无用。尤其司铎应当谨守节制，恐怕他们是为炫耀而施舍，而非为了正义。如今乞丐的贪婪前所未有地强烈。他们身强力壮地前来，除了游荡之外别无理由。他们想掏空穷人的钱囊——夺去他们的生活之资。小额不能满足，他们索求更多。他们以蔽体的衣衫为理由强索，又编造生活谎言以求更多钱财。若有人过于轻信其言，很快便会耗尽本应用于赡养穷人的资金。我们的施舍应当有条理，使贫者不致空手而去，也使急需者的生活之资不致耗尽而成为狡诈者的战利品。因此，须有适度，好使慈悲永不废弃，真正的急需永不遭忽视。\par

许多人假称负债。应查究真相。他们哀哭自己被强盗洗劫一空。遇此情况，只有确见祸患、或熟知其人，方可予以信任；然后慨然援助。对被教会拒斥者，若缺少食物，仍当供给。因此，施舍有度者，对任何人都非苛刻，却对所有人慷慨。我们不应仅仅侧耳聆听恳求者的呼声，更应注目细察他们的急需。善良的管家对软弱的召唤，比穷人的呼声更为敏锐。喧嚷不休的乞丐固然常能索得更多，但我们不应总向无礼让步。应当看见那看不见你的人。应当寻找那耻于见人的人。那在狱中的人，必须进入你的思念；那为疾病所困的人，虽不能达于你耳，却当呈现于你心。\par

78. 人们越看见你施予慈悲的热忱，就越爱慕你。我认识许多司铎，他们施予得越多，也就拥有得越多。因为那些看见一位善良的分施者的人，会给他一些东西，让他在履行职务时去分发，确信这慈悲之举必会达到穷人。若他们见他施予时过分大方或过于吝啬，便会鄙视这两种态度：前者因他在不必要的开销上浪费了他人劳苦的成果，后者则因他将之囤积于钱袋。因此，既需在慷慨中谨守节度，有时又似需加上鞭策。节度，是为使所施的恩惠能日复一日地延续，免致我们不得不从急需者那里收回我们任意浪费挥霍之物。鞭策，是因金钱用于穷人的食粮，胜于富人的钱包。我们务须提防，勿将贫困者的福祉封锁在我们的钱柜中，勿将穷人的生命如同埋葬于坟墓。\par

79. 若瑟本可送出埃及的一切财富，耗费王室的珍宝；但他甚至不愿显得浪费他人的财物。他宁愿出售谷物，而不愿送给饥饿者。因为若他只送给少数人，多数人便将一无所得。他充分证明了那种使众人都得饱足的慷慨。他开放仓库，让众人都可购买他们的粮饷，以免他们若白白领受，便会放弃耕种田地。因为那享用他人之物的人，常会忽视自己的事。\par

80. 首先，他收集了他们的金钱，然后是他们的器具，最后他为国王取得了他们对土地的全部权利。\BibleRef{创 47:14} 他并非要剥夺所有人的财产，而是要支持他们保守它。他也制定了一项普遍的税赋，好让他们能安然保有己物。这令那些被他取去土地的人大为悦纳，以至他们将此视为恢复福祉，而非出售权利。因此他们说：\BibleQuote{你救了我们的命，愿我们在主眼前蒙恩。} \BibleRef{创 47:25} 因为他们并未失去自己的任何东西，反而领受了一项新的权利；凡对他们有用之物无一缺失，因为他们已永远获得了它。\par

81. 啊，高贵的人！他不寻求无益施舍的虚浮光荣，却以其远见的恒久利益作为他的纪念。他做事的方式，使人民能借着缴纳来帮助自己，而在急需之时不必求助他人。因为放弃部分收成，无疑胜于丧失全部权利。他定税赋为全部出产的五分之一，从而显出他为将来筹谋的明达，并对他所征的税显出慷慨。此后埃及再未遭受这样的饥荒。\par

82. 他推断未来的能力何其卓越。首先，当他解释国王的梦时，何等明敏地道出了真理。这是国王的第一个梦：有七只母牛从河中上来，色美肉肥，在河岸上吃草。随后，又有七只公牛，色丑肉瘦，从河里上来，站在河岸上那些母牛旁边吃草。这些瘦弱丑陋的公牛，竟吞食了那些肥美好看的母牛。这是第二个梦：有七个饱满而好的麦穗从地上长出。随后，又有七个枯槁、被东风吹焦的麦穗，试图取而代之。这些枯槁细弱的麦穗，竟吞噬了那丰满佳美的麦穗。\par

83. 若瑟如此解释此梦：七只母牛代表七年，七个麦穗也代表七年——他借着牲畜与谷物的出产来解释时间。因为母牛产犊需时一年，而一季收成也满全一年。它们从河中上来，正如时日、年岁、时光流逝，如河水般急速流去。因此他指出，这地的前七年将丰饶多产，而后七年则将荒歉无产，其荒歉将吞噬先前的丰饶。所以他警告他们，务必在丰年积存粮秣，以备将来荒年之需。\par

84. 我们首先当赞叹什么呢？是他那深入真理之境的智识能力？还是他预见如此重大而长久之所需的谋略？或是他的醒寤与正义？借着醒寤，在身居高座之时，他积聚了如此庞大的储备；借着正义，他公平善待众人。至于他的心胸广大，我又当如何言说？虽被兄弟们卖为奴隶，\BibleRef{创 37:28} 他却未报复这项不义，反而终止了他们的匮乏。他的温和又当如何？他借着一种虔诚的计谋，谋求他所爱的兄弟之临在，他假借一场精心设计的窃案，宣称他偷了自己的财物，好将他视为爱的抵押？\par

85. 因此，他的父亲理当对他如此说：\BibleQuote{我的儿子若瑟长成壮大，我的儿子长成壮大了，我幼小的儿子，我所爱的。我的天主帮助了你，以高天的祝福、以普地的祝福祝福了你，这地拥有一切，因你父亲和母亲的祝福。这祝福超越了永世山岳的祝福，超越了永恒丘陵的愿望。} 且在\BibleBook{申命}{纪}中：\BibleQuote{你这曾在荆棘丛中显现的，愿你临于若瑟的头上，他的头顶上。他在兄弟中备受尊敬，他的光荣如头胎的公牛；他的角如独角兽的角，用以抵触万邦，直到地极。这是厄弗辣因的万军，是默纳协的千军。} \BibleRef{申 33:16}\par

% structure: p0119 source=h2 target=subsection reason=relative-heading-rank; base=h2

\subsection{第十七章}

\Emph{在我们所咨询的人身上应具有的德性，若瑟与保禄如何具备了这些德性。}\par

因此，那给人出谋划策的人，理当如此：他可以在一切善工、教导、品格真诚和庄重上，献上自己作为典范。这样，他的言语才是健全而无瑕的，他的劝告才有用，他的生活才有德行，他的见解也得体。\par

如此的人正是\Person{保禄}，他向贞女们提供了劝告\BibleRef{格前 7:25}，向司铎们给予了指导，以便将自己立为我们效法的榜样。因此他懂得如何谦卑，正如\Person{若瑟}所作的那样：他虽出身于圣祖们的尊贵家族，却不以自己卑下的奴隶身份为耻；反倒以甘心的服侍使其增辉，并以自己的德行使其荣耀。那曾经过买主和卖主之手，并称他们为\Quote{主人}的，他懂得如何谦卑。请听他谦卑自下的话：\BibleQuote{我主人因为我的缘故，不晓得他家里还有什么，他把所有的一切都交在我手中；除了你，因你是他的妻子，他并没有留下一样不给我；我怎能做这极大的恶事，得罪天主呢？}\BibleRef{创 39:8} 他的话语充满了谦卑，也充满了贞洁。谦卑，因为他顺服他的主人；精神尊贵，因为他感恩；贞洁，因为他认为被如此大罪玷污乃是可怕的罪过。\par

因此，那出谋划策的人理当如此。他必须没有任何阴暗、欺诈或虚假，以免给自己的生活和品格投下阴影；没有任何邪恶或罪恶，使前来寻求劝告的人退缩。因为有些事，人躲避；有些事，人鄙视。我们躲避那些能造成伤害的事，或是那些能背信弃义地、悄然滋长以伤害我们的事；比如，我们所求问的人其信誉可疑，或贪图钱财，以致一定数目的金钱就能使他改变主意。若有人行事不义，我们便躲避他，回避他。一个追求享乐又奢侈的人，即使他行事不虚，却贪婪且过分贪爱不义之财；这样的人会遭鄙视。一个沉溺于懒散怠惰生活的人，能给出什么勤劳的证据，能结出什么劳动的果实？他的心中，又能有过什么牵挂和忧虑呢？\par

因此，那出好主意的人说：\BibleQuote{我学会了，无论在什么景况下，都可以知足。}\BibleRef{斐 4:11} 因为他知道，贪爱钱财是万恶的根源\BibleRef{弟前 6:10}，所以他满足于自己所有的，不求他人的财物。他说，我所有的，为我足够了；无论多少，对我来说都是丰富。他似乎想尽可能清楚地说明这一点。他引用了这样的话：\Quote{我知足，}他说，\Quote{于我所有。} 这意味着：\Quote{我既无缺乏，也不过多。我没有缺乏，因为我不再寻求更多；我没有过多，因为这不是为我自己，而是为众人。} 这番话是针对钱财说的。\par

然而，他本可以对一切事都这样说，因为他当时所有的一切都使他知足；就是说，他不求更高的尊荣，不寻更多的服侍，不贪虚荣，也不指望那不应当的感激；而是忍耐劳苦，确信自己的功劳，等待他必须忍受的斗争终结。他说：\BibleQuote{我也知道}，他说，\BibleQuote{怎样受穷。}\BibleRef{斐 4:12} 未经教导的谦卑并不值得称赞，唯有那种具有节制和自我认识的谦卑才配得。因为有一种谦卑是基于恐惧，还有一种则是基于缺乏技巧和无知。故此圣经上说：\BibleQuote{他必拯救心灵谦卑的人。} 因此，他光荣地说：\BibleQuote{我知道怎样受穷；} 就是说，在何处，以何种节制，为何目的，在何种职责，在何种职分上。法利塞人不知道怎样受穷，因此被贬抑；税吏却知道，因而得以成义。\BibleRef{路 18:11}\par

\Person{保禄}也知道怎样富裕，因为他拥有富有的灵魂，尽管他不拥有富人的财宝。他知道怎样富裕，因为他不寻求金钱的馈赠，而是期待在恩宠中结出果实。我们也可以从另一层意思理解他所说的\Quote{知道怎样富裕}的话。因为他能再次说：\Quote{啊，格林多人！我们对你们口是开着的，我们的心是宽宏的。}\BibleRef{格后 6:14}\par

他在一切事上都习惯了饱饫和饥饿。那在基督内知道怎样饱饫的人是有福的。知识带来的饱足不是肉身的，而是属灵的。对知识的需求诚然是正当的：\BibleQuote{因为人生活不但靠食物，而且也靠天主口中所发的一切言语。}\BibleRef{申 8:3} 因为那知道怎样饱饫的人，也知道怎样饥饿，为能常常寻求新的事物，渴望天主，渴慕上主。他知道怎样饥饿，因为他知道饥渴者将得饱饫。\BibleRef{玛 5:6} 他也知道怎样富裕，并且能够富裕，因为他一无所有，却拥有一切。\BibleRef{格后 6:10}\par

% structure: p0128 source=h2 target=subsection reason=relative-heading-rank; base=h2

\subsection{第十八章}

\Emph{从十支派与国王\Person{勒哈贝罕}分裂的事实，我们可知邪恶的谋士所能造成的危害。}\par

因此，正义尤其使执掌任何职分的人显耀；反之，不义则使他们失败，并与他们为敌。圣经本身给了我们一个例子，其中说：以色列人在\Person{撒罗满}死后，请求他的儿子\Person{勒哈贝罕}使他们脱离残酷的轭，并减轻他父亲统治的严厉；他却轻看了老年人的劝告，顺从少年人的建议这样回答：\Quote{我父亲加重了你们的轭，我还要再加重；我父亲使你们肩负轻担，我要使你们肩负更重的。}\par

听到这话，人民被激怒了，说：\Quote{我们与\Person{达味}没有分，与\Person{叶瑟}的儿子也没有产业。以色列啊，回你的帐篷去吧！因为我们不要这人作我们的君王或领袖。}于是，他被人民离弃、抛弃，为了\Person{达味}的缘故，他仅能保留十个支派中的两个。\par

% structure: p0132 source=h2 target=subsection reason=relative-heading-rank; base=h2

\subsection{第十九章}

\Emph{许多人因正义、仁爱和谦恭而被赢得，但这一切必须出于真诚。}\par

95. 因此，显而易见，公正使帝国稳固，不义却使帝国覆亡。邪恶连一个家庭都不能治理，又如何能维系一个王国呢？因此，极需无比的慈爱，好使我们不仅能维护公共事务的管理，也能维护个人的权利。仁爱价值至高；因为它力求以恩惠包容众人，借着履行义务将人系于自己，并以其魅力使人对自己忠诚。\par

96. 我们也说过，言语的温良对赢得人心大有影响。但我们希望这温良真诚而明智，毫无谄媚，免得谄媚损害我们言谈的纯朴与纯洁。我们不仅要在行为上，也要在言语、纯洁和\OriginalTerm{信德}{faith}上，成为他人的楷模。我们愿别人怎样看我们，我们就当是怎样的人；并且要公开表露我们内在的情感。我们不可在心里说不义的话，以为能藏在静默中，因为那创造隐秘之物的主听见隐秘之言，知道人心中的秘密，并在人内植入了情感。因此，我们应如同在审判者眼前一般，看待我们的一切行为，如同摆在光明中，使其显明给众人。\par

% structure: p0136 source=h2 target=subsection reason=relative-heading-rank; base=h2

\subsection{第二十章}

\Emph{与善人交往对所有人，尤其对年轻人，极为有益，正如\Person{若苏厄}与\Person{梅瑟}等人的榜样所证明的。再者，年龄相异的人往往在德行上相似，\Person{伯多禄}与\Person{若望}便是明证。}\par

97. 与善人结交，实为美事。年轻人追随伟大而明智者的指导，也极有益处。因为与明智者同住，自己便成为明智者；依恋愚昧者的人，也被视为愚昧。这种与明智者的友谊，既大有助益于教导我们，也给予我们的正直一个确切的明证。年轻人会很快显示出，他们效法那些他们所亲近的人。这一见解基于以下事实：在他们全部的日常生活中，他们逐渐与那些他们充分交往的人相似。\par

98. \Person{农}的儿子\Person{若苏厄}成为如此伟大，因为与\Person{梅瑟}的结合不仅教导他认识法律，而且\OriginalTerm{圣化}{sanctifying}他，使他领受\OriginalTerm{恩宠}{grace}。当在主的\OriginalTerm{帐幕}{tabernacle}中，主的尊威以其\OriginalTerm{神圣临在}{divine Presence}闪耀显现时，惟独\Person{若苏厄}在帐幕内。当\Person{梅瑟}与天主交谈时，\Person{若苏厄}也被圣云遮蔽。司祭和民众站在下面，\Person{若苏厄}与\Person{梅瑟}上山接受法律。全体民众都在营内；\Person{若苏厄}却在营外，在约柜的帐幕内。当云柱降下，天主与\Person{梅瑟}说话时，他作为忠信的仆人侍立一旁；他，一个年轻人，并未走出帐幕，尽管远远站着的长老们因这些神圣奇事而战栗。\par

99. 因此，在所有这些奇妙作为和可畏奥秘中，惟独他紧随圣\Person{梅瑟}身旁。因此，那位在此与天主的交往中作为他同伴的人，继承了他的权柄。\BibleRef{申 34:9}他实在堪当挺身而出，成为能遏止河流者，并能说：\Quote{太阳，停住！}使夜晚延迟，白日延长，好似为见证他的胜利。\BibleRef{苏 10:12}为何？——一个\Person{梅瑟}未蒙受的祝福——惟独他被拣选，引导民众进入\Place{福地}。他是一个人，因\OriginalTerm{信德}{faith}所行的奇事而伟大，因他的凯旋而伟大。\Person{梅瑟}的作为属于更高的型式，而他带来的成功更为显著。这两人，因蒙受天主\OriginalTerm{恩宠}{grace}的扶助，都超越了凡人的地位。一位统治海洋，另一位统治天空。\par

100. 因此，老幼之间的结合是美好的。一位作见证，另一位给予安慰；一位提供指导，另一位带来喜乐。我不提\Person{罗特}，他年轻时依恋着正要出发的\Person{亚巴郎}。\BibleRef{创 12:5}因为有些人或许会说，这主要是因他们的亲属关系，而非出于他的自愿行为。那么，对于\Person{厄里亚}和\Person{厄里叟}，我们该说什么呢？尽管\WorkTitle{圣经}没有明确说明\Person{厄里叟}是年轻人，但我们可从中推断他较为年轻。在\BibleBook{宗徒}{大事录}中，\Person{巴尔纳伯}带着\Person{马尔谷}，\Person{保禄}带着\Person{息拉}\BibleRef{宗 15:39}和\Person{弟茂德}\BibleRef{宗 16:3}以及\Person{弟铎}。\BibleRef{铎 1:5}\par

101. 我们也看到，职责按照他们在任何方面的优长而分配。长老们率先提供建议，年轻人则表现出积极性。通常，那些德行相似而年龄不同的人因他们的结合而极其喜乐，正如\Person{伯多禄}和\Person{若望}一样。我们在\WorkTitle{福音}中读到，\Person{若望}是年轻人，甚至在他自己的言语中，尽管他在功劳和智慧上不逊于任何长老。在他身上，有着可敬的品格成熟和如同白发老者般的明智。无玷的生活配得上美好的老年。\par

% structure: p0143 source=h2 target=subsection reason=relative-heading-rank; base=h2

\subsection{第二十一章}

\Emph{保护弱者，帮助异乡人，或履行类似的职责，极大地增添一个人的价值，尤其对于经受过考验的人而言。同时，贪爱钱财会招致重责；对于\OriginalTerm{司铎}{priests}而言，浪费也深受谴责。}\par

102. 当一个人从权势者手中拯救穷人，或使被判死刑的罪犯免于一死时，他所受的敬重也大大增加；只要这事做得不引起骚动，以免我们似乎是为炫耀而非出于怜悯而行，免得我们在试图医治轻微创伤时，反而造成更重的伤害。但是，如果一个人解救了一个被权势者的资源和党羽所压垮，而非因自身邪恶所应得之罪而沦陷的人，那么一种伟大而高尚的舆论见证便变得强而有力。\par

103. 款待也有助于使许多人得人称赞。因为这是一种公开显示仁慈情感的方式：使异乡人不会缺乏款待，反而受到有礼的接待，并且当他到来时，门为他敞开。在全世界眼中，最合适的乃是：异乡人应受尊敬地接待；款待的魅力在我们的餐桌前不应缺乏；我们应以随时和自如的服务迎接客人，并期待他的到来。\par

104. 这尤其是\Person{亚巴郎}受称赞之处，因为他守在自己的帐幕门口，惟恐有旅客偶然经过。他留心观察，为遇见旅客，并抢先上前，请他不要走过，说： \BibleQuote{我主，如果我蒙你垂爱，请不要由你仆人这里走过去。} \BibleRef{创 18:3}因此，他因款待而得报偿，获得了后裔的恩赐。\par

105. 他的侄子\Person{罗特}\BibleRef{创 19:20}，不但在血统上，而且在德行上与他相近，由于他乐于款待，使自己和全家免于\Place{索多玛}的惩罚。\par

106. 因此，人应当好客、和善、正直，不贪图他人财物，若受攻击，甘愿放弃自己的一些权利，而不夺取他人的。他应避免争讼，憎恶争吵。他应恢复合一与宁静的恩宠。一个善人放弃自己的一些权利，不仅是\OriginalTerm{慷慨}{liberalitas}的标志，也会带来巨大的利益。首先，免于诉讼之费，这并非小利。其次，友谊的增进带来众多好处，日后对那当时能轻视小利的人极为有用。\par

107. 在款待的一切职责中，应向众人表达善意，但对于正直的人更当尊重。因为主说：\BibleQuote{谁若接纳一位义人，因他是义人，将要领受义人的赏报。}\BibleRef{玛 10:41}款待在主前如此蒙恩，甚至一杯凉水也不会得不到赏报\BibleRef{玛 10:42}。你看，\Person{亚巴郎}在寻找客人时，竟接待了天主本身。你看，\Person{罗特}接待了天使\BibleRef{创 19:3}。而你又怎能知道，当你接待人时，你不是接待\Person{基督}呢？基督可能就在那来访的旅客身上，因为基督就在穷人身上，正如祂自己说的：\BibleQuote{我在监里，你们来探望了我；我赤身露体，你们给了我穿的。}\BibleRef{玛 25:36}。\par

108. 所以，追求恩宠而不追求金钱，乃是美事。诚然，这种邪恶早已侵入人心，以致金钱取代了尊荣，人心充满对财富的艳羡。于是\OriginalTerm{贪爱金钱}{philargyria}潜入，仿佛使一切善行枯竭；人们将超常的开销视为损失。然而，就在此处，圣经早已警诫贪爱金钱，以免它成为阻碍，说： \BibleQuote{有情吃蔬菜，胜于无情食肥牛。}\BibleRef{箴 15:17}又说：\BibleQuote{干饼一张而平安共食，胜于满屋佳肴而互相争吵。}\BibleRef{箴 17:1}因为圣经教训我们，不当挥霍，而当慷慨。\par

109. 施予有两种：一来自\OriginalTerm{慷慨}{liberalitas}，一来自\OriginalTerm{挥霍浪费}{prodigalitas}。接纳旅客、衣裸者、赎俘虏、济穷乏，是慷慨的标志。耗费金钱于奢华筵宴与美酒，便是挥霍。因此经上说：\BibleQuote{醇酒致挥霍，酗酒致侮慢。}\BibleRef{箴 20:1}仅为博取民众欢心而耗费自己的财富，便是挥霍。有些人耗尽祖产于竞技场的赛会，或戏剧表演和角斗士展演，甚至野兽搏斗，只为在这些事上超越祖先的名声。他们所做的这一切尽是愚妄，因为即使在善工上过度花费金钱也是不合宜的。\par

110. 真正的慷慨在于对穷人本身保持适度，以便有足够的资助更多的人；不为求宠而超越合理的界限。凡出于纯洁真诚的意向所行的，便是合宜的。不轻率从事不必要的计划，也不忽略当尽的本分，这也是合宜的。\par

111. 但对于\OriginalTerm{司铎}{sacerdos}来说，尤应以相称的光辉装饰天主的圣殿，使上主的庭院因他的努力而显赫。他应始终按慈悲的要求花费金钱。他理当给予旅客所应得的。这不可过多，只要足够；不可超过，只需符合善意所要求的，如此他绝不至为博取他人欢心而损害穷人，也不至对\OriginalTerm{圣职人员}{clerus}显得过于吝啬或过于慷慨。前者不近人情，后者则是挥霍。若为救助那些应从其可怜营生中挽回的人所必要的金钱匮乏，便是不近人情。若有过多金钱用于享乐，便是挥霍。\par

% structure: p0155 source=h2 target=subsection reason=relative-heading-rank; base=h2

\subsection{第22章}

\Emph{我们必须在过分的温和与过度的严厉之间持守正确的分寸。那些以虚假的温和试图钻入他人心中的人，不会获得任何实质或持久的东西。\Person{阿贝沙隆}的例子清楚地表明了这一点。}\par

112. 再者，甚至我们的言语和教导也应有适当的分寸，以免显得过于温和或过于严厉。许多人宁愿过于温和，以显得善良。但确定的是，任何虚假或伪装都无法具有真德行的形式；不但如此，它甚至无法持久。起初它茂盛，随后，随着时间流逝，便如小花般凋谢消逝，但真实诚恳的却有深根。\par

113. 为用实例证明我们所言属实——即虚伪的不能持久，一时茂盛却迅速失败——我们将从那个家族中取一个伪装与虚假的例子，我们已从那里引用了许多实例，以显示他们在德行上的成长。\par

114. \Person{阿贝沙隆}是\Person{达味}王的儿子，以美貌著称，容貌俊美，青春正盛；在\Place{以色列}全境，没有像他这样俊美的人。他从脚掌到头顶，毫无瑕疵。他为自己预备了马车和马匹，并派五十人在他前面奔跑。阿贝沙隆清早起来，站在城门的路旁，凡有争讼要去求王判断的人，阿贝沙隆就叫住他，问说：\Quote{你是从哪一座城来的？}那人回答：\Quote{你的仆人是\Place{以色列}一个支派的人。}阿贝沙隆就对他说：\Quote{你的话很好，而且合情合理；但是王没有派人来听你的案件吗？唉，谁立我作判官呢？凡有争讼来到我这里的人，我必给他公正的裁判。}他用这些话笼络他们。当他们前来向他致敬时，他便伸手拉住他们，与他们亲吻。他便这样夺取了众人的心。因为这种奉承的话很快就渗透到人的内心深处。\par

115. 那些被宠坏而又野心勃勃的人，选择了那暂时看来是荣耀、令人喜悦和享受的东西。但就在这位最具智慧的先知认为应当暂缓行事之时，他们却再也无法忍耐和承受了。后来，\Person{达味}对胜利深信不疑，便嘱咐那些出去作战的人，要饶恕他的儿子。他本人不愿亲自参与战斗，以免让人觉得他拿起武器对抗那虽是儿子、却企图杀害父亲的人。\par

116. 由此可见，那些真实、出自真诚而非虚假之心的事物，才是持久稳固的；而那些由伪装和奉承所带来的事物，却绝不能长久。\par

% structure: p0162 source=h2 target=subsection reason=relative-heading-rank; base=h2

\subsection{第二十三章}

\Emph{那些容易被金钱收买或奉承讨好的人的忠诚，是靠不住的。}\par

117. 谁会认为那些被金钱收买而服从的人，或是被奉承引诱的人，会永远对人忠心呢？因为前者随时准备出卖自己，而后者则无法忍受严格的管束。他们很容易被一点奉承所笼络，但若有人用一句话责备他们，他们就会抱怨，离弃人，心怀敌意地离开，愤怒地抛弃人。他们宁愿统治而不愿服从。他们认为那些本应由他们置于自己之上的人，反倒应当服从自己，仿佛这些人欠了他们恩情似的。\par

118. 有谁会认为那些他以为必须用金钱或奉承才能笼络住的人，会对自己忠诚呢？因为拿了你的钱的人，会觉得他被廉价地看待和轻视，除非一再付钱给他。因此，他常常期待着他的报酬；而那个被祈求和奉承的人，则总是想要别人再来求他。\par

% structure: p0166 source=h2 target=subsection reason=relative-heading-rank; base=h2

\subsection{第二十四章}

\Emph{我们必须只以正当的方式谋求高升。担任职务必须明智且有节制地履行。下级神职人员不应以虚假的德性贬损主教的声誉；同样，主教也不应嫉妒神职人员，而应在一切事上保持公正，尤其是在作出判断时。}\par

119. 因此，我认为，尤其是在教会中，一个人应努力只以良好的行为和正当的目标去谋求高升；这样，就不会有骄傲自负，不会有漫不经心的懈怠，不会有可耻的内心倾向，不会有不当的野心。一颗纯朴率直的心便足以应对一切，也足够自荐了。\par

120. 再者，担任职务时，不应苛刻严厉，也不可过于宽纵；免得一方面让人觉得我们在施行专制权力，另一方面又让人觉得我们根本没有履行所担任的职务。\par

121. 我们也必须努力以我们所能做的善行和职责去赢得许多人，并维护已有的恩惠。因为如果他们现在因某种严重的冤屈而恼怒，他们自然会忘记从前的恩情。的确常有这样的事：一个人若以某种不恰当的方式决定把别人置于他们之前，那些曾蒙受恩待并被一步步提升的人，就会被赶走。然而，一位司铎在施恩和决断时，理应表现出公正，并像对待父母一样尊重其他的神职人员。\par

122. 那些曾蒙悦纳的人，如今不应变得傲慢，反而应记念所领受的恩宠，在谦卑中坚定不移。一位司铎，若因某位神职人员、侍从或任何教会成员，通过施行怜悯、守斋、品行正直或教导与诵读而赢得了重视，就不应感到不悦。因为教会的恩宠便是对教导者的赞誉。赞许他人的工作是件好事，只要其中毫无夸耀的意图。因为每人都应从他人口中而非自己口中获得赞誉，并且每人都应因他所行的事，而非仅因他所存的愿望而受到称赞。\par

123. 但若有人不服从他的主教，想要自我抬举、自我高升，并以虚假的学识、谦卑或怜悯的外表来掩盖主教的功绩，他便因骄傲而偏离了真理；因为真理的准则是：不做任何推进自己而损害他人的事，也不利用自己所拥有的一切美善去羞辱或指责他人。\par

124. 决不可庇护恶人，也不可容许将圣物交予不配之人；另一方面，也不要去骚扰和逼迫那些过错尚未清楚证实的人。在任何事情上，不公正都会很快引起反感，尤其在教会中更当如此，因为在教会里应当有公平，人人应受同等的对待，以免有权势的人索取更多，富人也占据更多。因为无论贫穷还是富有，我们在基督内都是一体。度更圣洁生活的人，不应因此为自己多求什么；他反而应当为此更加谦卑。\par

125. 施行审判时，我们不可徇情面。必须撇开偏袒，按案情本身的是非曲直断案。最破坏别人对我们的好感或信任的，莫过于在受理的任何案件中，把弱者的官司判给强者。同样，如果我们对穷人严苛，却为犯罪的富人开脱，也是如此。人很惯于奉承居高位者，以免他们自以为受到伤害，或感到恼怒，好像被推翻了似的。但如果你怕得罪人，那就不要担任审判之职。如果你是司铎或某位圣职人员，就不要坚持去审判。若是钱财上的事，你虽可保持缄默，然而为求表里一致，总应站在公道一边。但在涉及天主的事上，那里有危及整个教会的危险，若装出一无所见的样子，便是犯下不小的罪。\par

% structure: p0175 source=h2 target=subsection reason=relative-heading-rank; base=h2

\subsection{第二十五章。}

\Emph{恩惠更应施与穷人而非富人，因为后者要么认为别人期望他们的回报，要么因似乎欠下如此举动的债务而恼怒。但穷人让天主代替自己成为债主，并诚心承认所领受的恩惠。在这些议论之外还加上一项警告，要人轻视财富。}\par

126. 但是，你对富人施恩又有什么益处呢？难道是因为他更愿意报答爱他的人吗？因为我们通常施恩给那些我们指望从他们得到回报的人。然而，我们更该顾念的，却是那些弱者和无力的人，因为我们希望，为了那没有回报能力的人，从主耶稣那里获得酬报。主耶稣曾借着一场婚宴的比喻\BibleRef{路 14:12}，概括地展现了德行。借此比喻，祂吩咐我们施恩，宁可给那些不能回报我们的人，教导我们要请赴我们宴席和饭局的，不是富人，而是穷人。因为富人被邀请，似乎是为了让他们反过来为我们预备一场盛宴；穷人既一无所有可作回报，当他们领受任何东西时，便使主成了我们的酬报者，因祂把自己献作穷人的担保。\par

127. 就常情而言，施恩于穷人比施恩于富人更有益处。富人轻视恩惠和以欠下人情为耻。更有甚者，凡所给予他的，他都视为按其功劳当得的，仿佛只是偿还一笔正当的债款而已；或者，他认为给予者只是希望从富人那里获得更大的回报。这样，在接受恩惠时，富人便基于此理，认为自己给予的，比他从前领受的更多。然而穷人虽无金钱可以回报，但至少会表露感激之情。在这一点上，可以肯定他回报的比领受的更多。因为金钱是用硬币偿付的，感激却永不枯竭；金钱越付出越减少，感激若囤积不用便会枯竭，若施予他人反能得以保全。再者——这是富人避之不及的——穷人承认自己感到受恩之债的束缚。他确实认为别人是给予了帮助，而非为了换取他的尊敬才予以提供。他认为自己的儿女再次被赐予，自己的生命得以恢复，家庭得以保全。那么，施恩于善人，比施恩于忘恩负义之辈，要好多少呢！\par

128. 因此，主对祂的门徒说：\BibleQuote{你们不要带金，也不要带银，也不要带铜钱。}\BibleRef{玛 10:9} 祂借此如同用镰刀，割除人心中不断滋长的爱财之心。伯多禄也对那生来就被人抬着的瘸子说：\BibleQuote{金子和银子，我没有；但把我所有的给你：因纳匝肋人耶稣基督的名字，起来行走罢！}\BibleRef{宗 3:6} 这样，他给予的不是金钱，而是健康。拥有无金钱的健康，比没有健康的金钱要好多少啊！那瘸子站起来了，这是他原先未曾希望的；他没有得到金钱，尽管他原曾希望得到。但在主的圣人中，财富几乎难觅踪影，因此财富成了他们轻看的对象。\par

% structure: p0180 source=h2 target=subsection reason=relative-heading-rank; base=h2

\subsection{第二十六章。}

\Emph{贪爱钱财之为恶，历时何等久远，从旧约的许多事例中可清楚看出。然而，拥有钱财是何等虚妄的事，也是一目了然。}\par

129. 然而，人的习性如此长久地倾注于这种对金钱的崇拜，以至于除非他是富人，便无人认为他配受尊荣。这并不是新近的风气。不，这种恶习（更糟的是）在许多年前就已滋长在人们心中。当耶里哥城在司铎们的号角声中倒塌，农的儿子若苏厄取得胜利时，他知道百姓的英勇因贪财爱金而被削弱。因为阿干从毁灭之城的战利品中取了一件金衣、二百银协刻耳和一根金条，他被带到主面前，既不能否认偷窃之举，便承认了。\BibleRef{苏 7:21}\par

130. 所以，贪爱钱财是一种自古即有、根深蒂固的恶习，甚至在天主法律颁布之时就已显露出来；因为曾颁下一条法律来加以遏制。\BibleRef{出 20:17} 由于贪爱钱财，巴拉克以为巴郎可因酬报而被引诱去诅咒我们祖先的百姓。\BibleRef{户 22:17} 若非天主命令他住口不诅咒，贪爱钱财本当获胜。阿干被贪爱钱财所胜，带领祖先的全族百姓走向灭亡。如此，那能使太阳停止不落的农的儿子若苏厄，却无法制止人心中的贪爱钱财悄悄蔓延。他的声音一出，太阳停住了，但贪爱钱财却不停歇。太阳停住时，若苏厄完成了他的凯旋，但当贪爱钱财蔓延时，他几乎失掉了胜利。\par

131. 为什么？难道那妇人\Person{德里拉}的贪爱金钱没有欺骗了众人中最勇敢的\Person{三松}吗？\BibleRef{民 16:6} 他，曾徒手撕裂吼叫的狮子；\BibleRef{民 14:6} 曾被绑住交在敌人手中，独自无助，却挣断绳索，击杀了一千人；\BibleRef{民 15:14} 曾扯断筋绳交织的绳索，就像扯断细网线一般；我说，他把头枕在那妇人膝上，就被夺去了那带来胜利之头发的美饰，那正是他的力量所在。金钱流入那妇人的怀中，天主的恩宠便离开了那人。\BibleRef{民 16:20}\par

132. 那么，贪爱金钱是致命的。金钱诱人，却也玷污拥有它的人，并不帮助没有它的人。即使金钱有时是一种帮助，那也只对表明自己需要的穷人才是帮助。对于不渴望它、不寻求它、不需要它的帮助，也不因追求它而偏离正道的人，它有何益处呢？对于别人，若拥有者独自因它而富有，它又有何益处呢？难道他因为拥有了那常令人丧失荣誉的东西，因为拥有了必须看守而非真正拥有的东西，就更可敬了吗？我们所使用的，我们才拥有；但我们使用的之外，带给我们的不是拥有的果实，只是看守的危险。\par

% structure: p0186 source=h2 target=subsection reason=relative-heading-rank; base=h2

\subsection{第二十七章}

\Emph{在鄙视金钱中，有正义的典范，主教和圣职人员应与其他人一起追求这一德行。另外略说几句：不应过快地施行绝罚。}\par

133. 最后，我们知道鄙视财富是正义的一种形式，因此我们应避免贪爱金钱，并尽全力切勿做任何违背正义的事，而要在一切言行中守护正义。\par

134. 如果我们愿意悦乐天主，就必须有爱德，必须同心合意，必须随从谦逊，各人看别人比自己强。这才是真正的谦逊：绝不骄傲地为自己要求任何事物，却自视卑微。主教应待圣职人员和侍从——他们实是他的子女——如自己的肢体，并给各人分配他所见适合的职责。\par

135. 肢体腐坏了，割除它并非没有痛苦。要用各种疗法治疗许久，看能否治愈。若无法治愈，良医才会割除它。同样，一位好主教的愿望是：愿意医治软弱的，除去扩散的溃疡，烧灼某些部分而不割除；最终，当它们无法治愈时，才怀着自身的痛苦割除之。因此，宗徒那美好的规则鲜明地呈现出来：我们各人不要单顾自己的事，也要顾别人的事。\BibleRef{斐 2:4} 这样，就决不会发生我们因愤怒而放纵自己的情感，或因偏私自己的愿望而让步逾矩的事。\par

% structure: p0191 source=h2 target=subsection reason=relative-heading-rank; base=h2

\subsection{第二十八章}

\Emph{即使慈悲会招致憎恶，也应慷慨施行。关于这一点，提到了众所周知的\Person{安博罗削}为赎回俘虏而打碎圣器的故事；并就教会所拥有金银的正当用途，给出了非常美好的劝告。接着，通过圣\Person{老楞佐}的行动指出教会的真正宝藏后，制定了一些在熔化和为这类用途使用祝圣过的圣器时应遵守的规则。}\par

136. 分担他人的不幸，量力而行、有时甚至超力而行地帮助他人的需要，这是激发慈悲心的极大动因。因为，为了慈悲，担当起事情或忍受憎恶，比显露冷酷之心要好。所以我曾因打碎圣器以赎回俘虏而招致憎恶——此事可能令亚略派不悦。并非他们不喜悦这行为本身，而是他们可借此抓住一些把柄来指责我。谁能如此刚硬、残忍、心如铁石，竟会因为一个人免于死亡，或一个女人免于野蛮人的污辱——那比死亡更可怕的事，或男童、女孩和婴儿免于偶像的污染——他们因惧怕死亡而受玷污——而不悦呢？\par

137. 我们这样做并非没有正当理由，然而我们也在百姓中继续这事，以便承认并一再表明：为主保存灵魂远胜于保存黄金。因为，那位不携带黄金派遣宗徒的主\BibleRef{玛 10:9}，也无需黄金地聚集了教会。教会拥有黄金，不是为了储存，而是为了使用，并花费在需要者身上。守护无用之物有何必要呢？我们岂不知亚述人从主的殿中夺走了多少金银吗？如果其他供应不足，司铎将它熔化来赡养穷人，岂不比被亵渎圣物的敌人抢走玷污更好？难道主自己不会说：\Quote{你为什么任由这么多穷乏人饿死呢？你肯定有黄金吧？你本当给他们食物。为什么这么多俘虏被带到奴隶市场，为什么这么多未被赎回的人被留下任由敌人杀害？} 保存活器皿比保存金器皿更好。\par

138. 对此无人能给出回答。因为你还能说什么呢？\Quote{我担心天主的殿会需要它的饰品？}祂会回答说：\Quote{圣事不需要黄金，也不专属于黄金——因为它们不是用黄金买来的。圣事的光荣在于赎回俘虏。的确，它们才是珍贵的器皿，因为它们救人脱离死亡。那才是主的真宝藏，能产生祂圣血所产生的效果。那时，主的圣血之器皿才确实被辨认出来，因为我们在这任何一次救赎中都看到，圣爵从敌人手中赎回了那些祂的圣血从罪恶中赎回的人。}当教会赎出一长串俘虏时，说得多么美妙：\Quote{这些是基督所赎回的。}请看那可通过火炼的黄金，请看那有用的黄金，请看那基督的黄金，使人免于死亡；请看那黄金，藉此，廉耻得以赎回，贞洁得以保存。\par

139. 于是，我情愿将他们作为自由人交给你们，而不是积攒金子。这一群被掳之人，这一队人，确实比金杯的景象更为荣耀。\OriginalTerm{救赎主}{Redeemer}的金子应当用于这项工作，好救赎那些处于危险中的人。我承认这一事实：\OriginalTerm{基督圣血}{blood of Christ}不仅闪烁在金杯中，而且还借救赎之功，赋予了它们天主作为的能力。\par

140. 这样的金子，圣\Person{老楞佐}这位殉道者为\OriginalTerm{主}{Lord}保存了。因为当人们向他索要教会的珍宝时，他许诺会将它们展示出来。第二天，他把穷人聚集起来。在被问及他所许诺的珍宝在哪里时，他指着穷人说：\Quote{这些就是教会的珍宝。}他们的确是珍宝，因为基督活在他们内，因为他们对祂怀有信德。同样，宗徒也说：\BibleQuote{我们是在瓦器中存有这珍宝。}\BibleRef{格后 4:7} 基督还能有比那些祂自称生活在其中的人更大的珍宝吗？因为经上记载：\BibleQuote{我饿了，你们给了我吃的；我渴了，你们给了我喝的；我作客，你们收留了我。}\BibleRef{玛 25:35} 又说：\BibleQuote{凡你们对我这些最小兄弟中的一个所做的，就是对我做的。}\BibleRef{玛 25:40} 耶稣还能有比那些祂乐意在其中显现的人更好的珍宝吗？\par

141. 圣\Person{老楞佐}指出了这些珍宝，并且获胜了，因为迫害者无法将它们夺走。\Person{耶苛尼雅}在围城时期保存了他的金子，没有用来购买食物，结果眼睁睁看着他的金子被掳走，自己也被掳去。而圣\Person{老楞佐}宁愿将教会的金子用在穷人身上，也不愿为迫害者保留，他因此获得了殉道的荣冠，这是对他独到而深远的洞见的奖赏。难道曾有人对圣\Person{老楞佐}说：\Quote{你不该花费教会的珍宝，也不该变卖圣器}吗？\par

142. 每个人都必须怀着真诚的信实和明达的先见来履行这一职责。若有人从中为自己谋取利益，便是罪行；但他若将珍宝用于穷人，或赎回被掳之人，便是行哀矜。因为没有人能说：为什么穷人活着？没有人能抱怨俘虏被赎回，没有人能因为主的圣殿得以建造而挑剔，没有人能因为为信友安葬遗体而扩地而恼怒，没有人能因为基督信徒的坟墓中死者得享安息而心烦。基于这三种方式，甚至可以拆卸、熔炼或变卖教会的圣器。\par

143. 必须注意，勿使奥秘之杯离开教会，以免圣爵的用途被转用于卑贱之事。因此，在教会中，首先寻找那些未曾被祝圣为如此神圣用途的器皿。然后，将它们拆毁，随之熔炼，分成小份分给穷人，也用于赎回被掳之人。但若没有新的器皿，或那些似乎从未被用于如此神圣用途的器皿，那么，正如我已说过的，我认为，所有这些都可被用于此种用途而不致失敬。\par

% structure: p0201 source=h2 target=subsection reason=relative-heading-rank; base=h2

\subsection{第二十九章}

\Emph{寡妇或所有信友托付给教会的财产，应当被保护，即使这会给自身带来危险。\Person{敖尼雅}司祭，及\Place{提契努姆}的\Person{安博罗削}主教的榜样，说明了这一点。}\par

144. 必须十分谨慎，好使寡妇所托付的财产保持不受侵犯。不仅要保护寡妇的财产，任何人的财产都应保护，且不引起抱怨。因为信实必须向所有人显示，尽管寡妇和孤儿的事应优先考虑。\par

145. 因此，受托于圣殿的一切财物，只以寡妇的名义保存着，正如我们在\BibleBook{玛加伯}{}书中所读到的。\EditorialNote{见：\BibleBook{玛加伯}{下} 第三章}因为当有人报告说，\Person{西满}曾奸诈地告诉\Person{安提约古}王，在\Place{耶路撒冷}的圣殿里有许多钱财，\Person{赫略多洛}就被派去调查此事。他来到圣殿，向大司祭说明了他的恶意情报及其来意。\par

146. 于是大司祭说，那里存放的只是供养寡妇和孤儿的生活费。当\Person{赫略多洛}正要前去夺取，并代表国王加以没收时，司祭们穿上了司祭长袍，俯伏在祭台前，含泪呼求生活的天主，因为祂曾经赐给他们关于托管财物的律法，求祂彰显自己为祂诫命的守护者。大司祭的面容与神色改变，显示出他内心的忧伤、焦虑和紧张。众人都哭了，因为如果连在天主的圣殿里也不能保存安全而又忠信的守护，这圣地必将受人轻蔑。妇女们束起胸衣，平素深居简出的贞女们也来敲门。有人奔向城墙，有人从窗户向外张望，所有人都向天举起双手祈祷，求天主维护祂的律法。\par

147. 但\Person{赫略多洛}并未因此退缩，执意要实施他的意图，并且已经带着他的随从包围了宝库。突然，有一位骑士显现给他，周身金光灿烂的甲胄，极其可畏，他的坐骑也佩着贵重的装饰。同时还有两位少年显现，他们英武神力，俊美非凡，辉煌荣耀，华丽盛装。他们站在他周围，从两边击打这亵圣的恶徒，不停手地鞭打他。我还需多说什么？他被黑暗笼罩，跌倒在地，几乎吓死，亲身经历了天主大能的明证，内心不再存有任何安全的希望。原先惧怕的人重获喜乐，先前骄傲的人反生恐惧。\Person{赫略多洛}的一些朋友见此不幸，便恳求\Person{敖尼雅}，为他求命，因他已奄奄一息。\par

148. 因此，当大司祭为此祈求时，那两位青年再次显现给赫略多洛，穿着同样的衣服，对他说：你该感激大司祭敖尼雅，因为看在他的份上，主才使你得以存活。至于你，既然遭受了天主的鞭笞，就去向你的朋友传报，你在这里学得了多少关于圣殿的圣洁和天主的大能。说完这话，他们便隐没了。于是赫略多洛恢复了生命，向主奉献了祭献，感谢了司祭敖尼雅，便率领军队回到国王那里，说：\Quote{如果你有一个仇敌，或是一个图谋不轨的人，就派他到那里去，你会看到他遍体鳞伤地回来。}\par

149. 所以，我的儿子们，对于托管金必须保持诚信，也必须显明关怀。如果你们藉着教会的援助，阻止了某些寡妇孤儿无法抵抗的强权压迫，并且表明主的命令在你们看来比富人的青睐更为重要，你们服务的光芒将更加闪耀。\par

150. 你们也记得，为了寡妇的托管金，甚而也为了别人的，我们曾多少次与皇室的攻击斗争。你和我曾共同分担此事。我还要提及提西农教会最近的事件，它险些失去了所收存的寡妇托管金。因为有人想凭皇帝的敕令要求这些款项，当时神职人员并没有维护他们的权利。因为他们自己，一旦被召担任职务并被派去调停，就认为不能违抗皇帝的命令。敕令的明文被宣读，法庭首席官员的命令就在那里，经办此事的人就在眼前。还有什么好说的呢？款项就被交出了。\par

151. 然而，在与我商议后，圣善的主教占有了他明知寡妇财产已被运入的那些房间。既然那些财产无法被运走，便全部记录在案。后来，凭着文件证明，又有人来索要。皇帝重申了命令，并要亲自接见我们。我们拒绝了。当我们向他解释了神圣法律的力量、一系列经文以及赫略多洛的危险后，皇帝终于变得通情达理。此后，又有人企图没收，但那位良善的主教赶在行动之前，将他所接收的全部归还给了寡妇。这样，诚信得以保全，而压迫不再令人恐惧；因为如今处于危险中的，是事情的本身，而非诚信。\par

% structure: p0211 source=h2 target=subsection reason=relative-heading-rank; base=h2

\subsection{第三十章}

\Emph{本书结尾劝戒避免恶意，追求明智、信德及其他德行。}\par

152. 我的儿子们，要躲避恶人，提防嫉妒者。恶人与嫉妒者之间的区别是：恶人因自己的好运而高兴，嫉妒者却因想到别人的好运而受折磨。前者爱邪恶，后者恨善良。因此，只愿自己得好处的人，几乎比愿人人都遭殃的人还可容忍。\par

153. 我的儿子们，行动之前要思考，经过长时间思考后，再去做你们认为正确的事。当值得称颂的死亡机会出现时，要立刻抓住。迟延的荣耀会飞走，并且不容易再次把握。\par

154. 要爱信德。因为约史雅就是因他的虔诚和信德，从敌人那里赢得了极大的爱戴。他在十八岁时举行了主的逾越节，是前人所未曾举行过的。既然他在热诚上超越了前人，那么，我的儿子们，你们也要表现出对天主的热诚。让对天主的热诚渗透你们，吞噬你们，好使你们每人都能说：\BibleQuote{为你的殿宇所怀的热诚把我耗尽。}基督的一位宗徒被称为热诚者\BibleRef{路 6:15}。但我何必说到宗徒呢？主亲自说过：\BibleQuote{为你的殿宇所怀的热诚把我耗尽。}那么，这热诚应是对天主的真实热诚，而不是卑劣的属世热诚，因为后者会招致嫉妒。\par

155. 愿你们中间有超越一切理解的平安。你们要彼此相爱。没有什么比爱德更甜蜜，没有什么比平安更蒙福。你们自己知道，我一向爱你们，如今也爱你们胜过一切。如同同一父亲的子女，你们已在兄弟情谊的纽带下团结一致。\par

156. 凡是好的，就要持守；愿平安与爱的天主在耶稣基督内与你们同在，愿尊崇和荣耀，权能和力量，归于祂，并归于圣神，直到永远。阿们。\par

\SourceRecord{Translated by H. de Romestin, E. de Romestin and H.T.F. Duckworth.；Nicene and Post-Nicene Fathers, Second Series；Vol. 10.；Edited by Philip Schaff and Henry Wace.；Buffalo, NY: Christian Literature Publishing Co.,；1896.；<http://www.newadvent.org/fathers/34012.htm>.}

% structure: p0001 source=h1 target=section reason=page-title

\section{论神职的义务（第三卷）}

% structure: p0002 source=h2 target=subsection reason=relative-heading-rank; base=h2

\subsection{第一章}

\Emph{我们由\Person{达味}和\Person{撒罗满}获得教导，应如何与自己的心商议。\Person{西庇阿}不应被视为归功于他的那句名言的原创者。作者证明圣洁的先知们在安闲之时成就了何等荣耀的事，并借他们和他人闲暇时刻的榜样表明，义人在困苦中绝非孤身一人。}\par

1. 先知\Person{达味}教导我们，应在心中行走，如同置身广厦；应与它交谈，如同与可信的伴侣对话。他自言自语，并与自己谈心，正如这些话所表明的：\BibleQuote{我说：我要谨守我的行径。}他的儿子\Person{撒罗满}也说：\BibleQuote{你应饮你自己井里的活水，由你自己泉源中汲取的流水！}\BibleRef{箴 5:15}意思是：运用你自己的谋略。因为：\BibleQuote{人心中的谋略，有如深水。}\BibleRef{箴 20:5}经上说：\BibleQuote{不许外人分享。你的泉源专为你而淌，应使你的井水专属你；当同你自幼结发的妻子共享其乐。愿妩媚的母鹿，娇柔的母羊，常与你亲昵。}\BibleRef{箴 5:17}\par

2. 因此，\Person{西庇阿}并非第一个知道自己在独处时并不孤单，或在闲暇时最不闲暇的人。在他之前，\Person{梅瑟}已知此事，他虽沉默，却在呐喊\BibleRef{出 14:16}；他虽安然站立，却在战斗，不，不仅战斗，且已战胜未曾靠近的敌人。他如此安闲，以致别人托着他的手；然而他并不比他人少行动，因为他以安闲的双手战胜了敌人，而那些在战场上的人却未能征服。\BibleRef{出 17:11}就这样，\Person{梅瑟}在沉默中发言，在安闲中辛劳。他的辛劳岂不胜于他的寂静之时？他在山上四十日，领受了全部律法？\BibleRef{出 24:17}在那孤寂中，有一位离他不远，与他交谈。因此\Person{达味}也说：\BibleQuote{我要听主天主在我内所说的话语。}天主向人说话，比人自言自语，要伟大多少倍！\par

3. 宗徒们经过时，他们的影子治愈了病人。\BibleRef{宗 5:15}他们的衣服被触摸，便赐予了健康。\par

4. \Person{厄里亚}发言，雨便停止，三年零六个月未降于地。他再次发言，面缸未空，油瓶未竭，直至整个长久的饥荒时期。\par

5. 然而——许多人以战争为乐——何种事更为荣耀？是以强大军队的力量结束一场战斗，还是仅凭在天主面前的功德？\Person{厄里叟}安居一处，而\Place{叙利亚}王向我们先祖的百姓发动大战，并以各种诡诈计谋增加恐怖，企图设伏捉拿他们。但先知洞悉他们的一切准备，藉天主的恩宠，其心神遍在一切处，他将敌人的心思告知同胞，并警告他们提防何处。\Place{叙利亚}王得知此事，便派遣军队围困先知。\Person{厄里叟}祈祷，使所有敌军被击打以致失明，并使那些前来围困他的人以俘虏身份进入\Place{撒玛黎雅}。\par

6. 让我们将他这种安闲与他人的相比较。他人为求休息，惯于将心思从俗务中抽离，远离人群与交往，寻求乡间僻静或田野孤寂，或在城中让心灵休息，享受平安与宁静。但\Person{厄里叟}始终活跃。在孤寂中，他分开\Place{约旦河}，使下游流去，上游倒流归源。在\Place{加尔默耳山}，他向那位迄今无子的妇女许诺，一个意想不到的儿子将为她诞生。他使死人复活，他纠正食物的苦涩，掺入面粉使之变甜。他分给百姓十个饼作食物，待他们吃饱后，拾起剩下的碎块。他使掉落并沉入\Place{约旦河}深处的斧头铁刃，因将木柄投入水中而浮起。他变麻风为洁净，变干旱为甘霖，变饥荒为丰裕。\par

7. 义人何时能孤单，因为他常与天主同在？他何时被离弃，因为他从不与基督分离？经上说：\BibleQuote{谁能使我们与基督的爱隔绝？我深信，无论是死亡，是生命，是天使，都不能。}\BibleRef{罗 8:35}劳累何时能夺去他的劳绩？因他绝不被夺去那劳绩冠冕的功德。何处能限制那拥有全世财富为产业的人？何种审判能困住那从不被任何人责难的人？因为他是\BibleQuote{像是人所不知的，却是人所共知的；像是待死的，看，他却活着；像是忧苦的，却常常喜乐；像是贫穷的，却使许多人富足；像是一无所有的，却拥有一切。}因为义人只看重合宜与美德的事。所以他虽在别人眼中看似贫穷，对自己却是富足的，因他的价值不凭短暂之物的价值衡量，而凭永恒之物的价值衡量。\par

% structure: p0011 source=h2 target=subsection reason=relative-heading-rank; base=h2

\subsection{第二章}

\Emph{哲学家们就德行与利益何者为重的讨论，与基督徒毫无关系。因为对基督徒而言，凡不公义之事便无利益。什么是全德的本分，什么是寻常的本分？同样的言辞常以不同方式适用于不同事物。最后，义人从不损人以利己，反常常寻求有益于他人的事。}\par

8. 我们既已论及前述两个主题，即讨论了何谓德行，何谓利益，现在随之而来的问题是：我们应否将德行与利益比较，并追问该追随何者。因为，我们已探讨过一事究竟是德行抑或邪恶，又在另一处讨论了有用抑或无用，所以此处有人以为应当探究一事究竟是德行还是利益。\par

9. 我之所以如此做，是为了避免让人以为我容许这两者彼此对立，而我已表明它们本为一体。因为我曾说过，凡德行之事必是有益的，凡有益之事必是德行的。因为我们不追随肉性的智慧，那种智慧认为财富充裕便是最大的利益；我们追随的是天主的智慧，依此智慧，世间极其看重的东西，反倒被视为损失。\par

10. 因为这一\OriginalTerm{完美义务}{\GreekText{κατόρθωμα}}，即完全而完美地履行的义务，始于德行的真正泉源。随之而来的是另一种，即普通的义务。从其名称即可看出，它不涉及严苛或非凡的德行实践，因为可为极多人所共有。许多人惯于积攒钱财；享用精心准备的宴席和可口的饭菜是一般习性；但守斋或自制却是少数人的操行，而不贪图他人财物更是一种罕见的德行。反之，欲剥夺他人的财产，且不安于己分——在此不乏有人为伴。那些（哲学家会如此说）是首要的义务——这些则是普通的。首要的只见于少数人，普通的则见于多数人。\par

11. 再者，同样的词语常有不同的含义。例如，我们称天主为善，也称人为善；但两者含义截然不同。我们在一种意义上称天主为义，在另一种意义上称人为义。同样，当我们称天主为智、称人为智时，含义亦不相同。福音中这样教导我们：\BibleQuote{你们应当是成全的，如同你们的圣父是成全的一样。}\BibleRef{玛 5:48} 我又读到，保禄是成全的，却又不是成全的。因为他说：\BibleQuote{这并不是说我已经达到这地步，或已成为成全的人；我还在竭力追求，好能夺得其所以将我夺得的。}\BibleRef{斐 3:12} 接着他立刻又说：\BibleQuote{所以，我们凡是成全的人……}\BibleRef{斐 3:15} 由此可见，成全有两种形式：一种仅是普通的成全，另一种则具有最高的价值。一种适用于今世，另一种适用于来世。一种符合人的能力，另一种则符合未来世界的完美。然而，天主在所有方面都是公义的，超越一切的智慧，在一切事上都是完美的。\par

12. 在人群中间也存在着差异。达尼尔，经上论到他说：\BibleQuote{谁比达尼尔更有智慧呢？}\BibleRef{则 28:3}，他所具有的智慧却与他人不同。撒罗满也是如此，他充满智慧，超越古人一切的智慧，且胜过埃及所有的智者。在一般意义上成为智者，与真正成为智者，这两者全然不同。普通的智者只为现世的事情考虑，并为自己着想，以便从他人那里夺取些什么，据为己有。真正智慧的人却不知如何顾念自己的利益，而是全心向往那永恒的事，向往那合宜而德行的事，不求一己之利，但求众人之益。\par

13. 那么，就让这成为我们的准则，以便我们绝不在德行的与有益的两者之间迷失方向。正直的人绝不可存心剥夺他人的任何东西，也绝不可希望损人利己以自肥。这准则是宗徒给予你们的，他说：\BibleQuote{凡事都可行，但不全是有益的；凡事都可行，但不全建树人。人哪，不要只求自己的事，也该求别人的事。}\BibleRef{格前 10:23} 这就是说：人不可求自己的利益，而该求别人的利益；人不可求自己的光荣，而该求别人的光荣。因此他在另一处说：\BibleQuote{各人不可只顾自己的事，也该顾别人的事，存心谦下，彼此该想自己不如人。}\BibleRef{斐 2:3}\par

14. 并且，谁也不可寻求一己的欢心或称赞，而该寻求他人的。我们在\BibleBook{}{箴言}一书中可以清楚地看到这一点，圣神藉着撒罗满说：\BibleQuote{我儿，你若智慧，对你自己和你的近人都有益；你若变为邪恶，你独自承担。}\BibleRef{箴 9:12} 智者像正直人一样，给人以劝告，并在拥有任何一种德行的形式上与他分担。\par

% structure: p0020 source=h2 target=subsection reason=relative-heading-rank; base=h2

\subsection{第三章}

\Emph{所给出的不可寻求一己私利的准则，首先由基督的榜样得以确立，其次由人这一词的涵义，最后则由我们肢体的构形与功能来确立。因此，作者指出，剥夺他人有益之物是何等罪恶，因为这种恶行既违犯了自然律，也违犯天主的法律。此外，藉此我们还会失落那使我们高于其他受造物的恩赐；最后，由于它，国家的法律也被妄用并遭到极度轻蔑。}\par

15. 所以，如果有人愿意取悦众人，他就必须凡事努力，不求一己之利，而求众人之益，如同保禄也曾努力去做的那样。因为这就是\BibleQuote{与基督的肖像相似}，\BibleRef{罗 8:29} 即，一个人不图谋他人之物，也不剥夺他人之物以归己有。因为我们的主基督，\BibleRef{斐 2:6} 祂虽具有天主的形体，却空虚自己，取了奴仆的形体，祂愿以自己的事工之德行来充实这形体。那么，你岂可抢夺基督所穿戴的人呢？你岂可剥去基督所装束的人呢？因为当你试图损人利己之时，你所做的正是此事。\par

16. 人啊，请想一想，你从何处得名——甚至是从土，它不索取任何人的东西，却慷慨地给予一切，并为一切生物提供多样的出产以供使用。因此，人道被称为人特有的一种内在德行，因为它辅助其同伴。\par

17. 你们身体的形态和肢体的用途，正教导你们这一点。一个肢体能要求另一个肢体的职责吗？眼睛能为自己要求耳朵的职责，或口要求眼睛的职责，或手要求脚的服事，或脚要求手的服事吗？不，即使是双手，左右手也有不同的职责，以至于如果任何人改变其中任何一个的用途，他就是违反自然行事。我们必须把整个人搁置一边，才能改变各个肢体的服事：例如，如果我们试图用左手取食，或用右手做左手的职责，比如清除食物残渣——当然，除非需要如此。\par

18. 暂且想象一下，赋予眼睛能力，从头部抽取理解力，从耳朵抽取听觉，从心灵抽取思考力，从鼻子抽取嗅觉，从口抽取味觉，然后让眼睛自己承担这些，这岂不立刻摧毁整个自然秩序吗？因此宗徒说得好：\BibleQuote{若全身是眼，哪里还有听觉？若全身是听觉，哪里还有嗅觉？}\BibleRef{格前 12:17}。所以，我们虽有许多肢体，却是一个身体，所有肢体对身体都是必要的。因为没有肢体可以对另一个说：\Quote{我不需要你。}因为那些似乎更软弱的肢体，反而更为必要，且需要更大的关心和照顾。\BibleQuote{若一个肢体受苦，所有的肢体都一同受苦。}\BibleRef{格前 12:26}\par

19. 如此，我们看到，剥夺他人——我们本应与之同甘共苦的人——的任何东西，或对我们应与其分享服事的人行不公或伤害，是何等严重的事。这是一条真正的自然律法，它约束我们向所有人表达一切善意，好使我们众人作为同一身体的各部分，轮流互相帮助，永不可想剥夺他人的任何东西，因为连拒绝给予帮助都是违反自然律法的。我们生来如此：肢体与肢体结合，一个与另一个协同工作，在互相服事中彼此协助。但若一个肢体不履行其职责，其余的就会受阻碍。例如，若手挖出眼睛，岂不是阻碍了眼睛工作的功用？若手伤害了脚，岂不阻止了多少行动？然而，整个人偏离其职责，比单独一个肢体偏离，要坏得多！若整个身体因一个肢体而受伤，整个人类团体也会因一个人而受扰乱。人类的本性受损，圣教会的团体也受损，这团体升为一个合一的整体，在信德和爱德的合一中紧密相连。主基督也为众人而死，祂将因祂宝血的代价徒然付出而悲伤。\par

20. 因为，正是主的律法教导我们必须遵守这条规则，好使我们永不可为自己的利益而剥夺他人的任何东西。因为律法上说：\BibleQuote{不可挪移你先人所立的地界。}\BibleRef{箴 22:28}。它吩咐将迷失的邻人牛牵回。\BibleRef{出 23:4}。它命盗贼被处死。\BibleRef{出 22:2}。它禁止克扣工人的工价，\BibleRef{肋 19:13}，并命人归还款项，不可取利。\BibleRef{申 23:19}。帮助一无所有的人，是善意的标记；但索取超过所给予的，则是冷酷的本性。如果一个人需要你的援助，因为他没有足够的自己的财物来偿还债务，你竟在善意的伪装下，要求那无力偿还较少款项的人支付更大数额，这难道不是邪恶的事吗？你只不过是将他从别人的债务中解救出来，却将他置于你自己的手掌之下；而你竟称那是人道的善意，其实不过是更大的邪恶。\par

21. 正是在这事上，我们高于所有其他生物，因为它们不懂得如何行善。野兽抢掠，人却与他人分享。因此圣咏作者说：\BibleQuote{义人常好施乐借。}然而，也有一些受野兽善待的：它们用所得喂养幼崽，鸟类用食物使雏鸟饱足；但唯独赐给人去喂养众人，如同他们自己的一样。这完全符合自然的要求。如果拒绝给予都不允许，那么剥夺他人怎么可能允许呢？而且，我们的法律岂不也教导我们同样的事吗？它们命令，那些以伤害人身或财产的方式从他人夺走的东西，必须附加赔偿归还；以便通过刑罚遏止盗贼偷窃，并通过罚款使他悔改。\par

22. 然而，假设有人不惧怕刑罚，或嘲笑罚款，难道这就使剥夺他人之物成为一件值得的事吗？那将是卑鄙的恶习，只配属最卑下之人。这与自然如此相悖，以至于虽然贫困似可逼使人如此行，但自然绝不能催迫这事。然而我们竟发现奴隶中有暗中偷窃，富人中则有公然抢劫。\par

23. 但是，有什么如此违反自然，胜于为了自己的利益而伤害他人呢？我们自己心中自然的情感，催迫我们为众人保持警觉，承受辛劳，为众人工作。而且，每个人不避风险地寻求众人的安宁，并认为拯救国家免于毁灭远比自己躲开危险更值得称谢，这也被视为光荣的事。我们必须认为，为国家劳苦，远较安乐悠闲、充分享受闲暇地度过平静生活，更为高贵。\par

% structure: p0031 source=h2 target=subsection reason=relative-heading-rank; base=h2

\subsection{第四章。}

\Emph{既然已指出，凡为了自己的利益而伤害他人者，必将受到自己良心的可怕惩罚，因此可以推论，对一个人无益的事，对所有人也同样无益。所以在基督徒中间，没有哲学家所提出的关于两个遭遇船难者的问题的余地，因为他们必须对所有人表现出爱德与谦逊。}\par

24. 由此我们推断，一个随从\OriginalTerm{本性}{nature}的指引、顺服本性的人，决不会伤害他人。如果他伤害了他人，就是违背了本性，也不会以为自己所得到的，与其说是利益，倒不如说是损失。有什么惩罚比内心的\OriginalTerm{良心}{conscience}创伤更痛苦呢？有什么审判比我们心灵的审判更严厉呢？每个人在其中被定罪，并因自己曾错误地加于弟兄的伤害而自责。圣经对此说得非常明白：\BibleQuote{愚人的口里，有招致恶行的棍杖。}\BibleRef{箴 14:3} 所以说，愚昧被定罪，因为它导致恶行。我们岂不应当避免愚昧，甚于死亡、损失、贫穷、流放或疾病吗？谁会不认为身体的瑕疵或遗产的损失，远远轻于灵魂的瑕疵或名誉的损失呢？\par

25. 因此，显而易见，众人都当思量并持守，个人的益处与众人的益处相同，除了对公共利益的益处之外，没有什么可被视为有益的。因为一个人怎能单独受益呢？那对众人无用的，便是有害的。我确实不能认为，一个对众人无用的人，能对他自己有用。因为，如果对众人有一条共同的\OriginalTerm{本性}{nature}律法，那么对众人也就有一种共同的利益处境。我们受本性律法的约束，要为众人的益处行事。所以，谁若希望按本性去顾念他人的利益，却违背本性律法去伤害他，那是不对的。\par

26. 因为，我们听说，在赛跑中，赛跑者被教导和告诫，人人都要靠脚步迅捷而非任何卑劣手段去赢得比赛，要尽全力奔跑，以争得胜利，而不可胆敢绊倒他人，或用手推开他。那么，在人生的赛程中，我们岂不更当不以虚谎和欺诈对待他人，而赢得胜利吗？\par

27. 有人问，在遭遇海难时，智慧人是否应该从一个无知的船员那里夺走一块木板？尽管从公共利益看来，一个智慧人逃脱海难似乎比一个愚昧人逃生更好，但我并不认为一个\OriginalTerm{基督徒}{Christian}、一个正义且智慧的人，应当以他人的死亡来保全自己的生命；正如当他遇到一个持械的强盗时，他不应回击，以免在自卫时玷污了对近人的爱德。关于这一点，福音书中的定论清楚明白。\BibleQuote{把你的剑放回原处；因为凡持剑的，必死在剑下。}\BibleRef{玛 26:52} 有什么强盗比那前来杀害基督的迫害者更可恨呢？但基督不愿受保护免于迫害者的创伤，因为祂愿意以自己的创伤医治万民。\par

28. 你为什么视自己高于他人呢？身为基督徒，本当看别人比自己强，不为自己图谋什么，不篡夺荣誉，不因自己的功劳要求赏报。再者，你为什么不愿承受自己的困苦，反倒要损毁他人的利益呢？因为，还有什么比不知足于自己所有的、或是寻求属于他人的并试图以可耻的方式获取，更违背本性的呢？如果\OriginalTerm{有德行的}{virtuous}生活符合本性——因为天主创造的一切本是十分美好——那么可耻的生活必然与之相背。有德行的生活与可耻的生活不能相容，因为它们已被本性律法绝对地分开。\par

% structure: p0038 source=h2 target=subsection reason=relative-heading-rank; base=h2

\subsection{第五章}

\Emph{正直的人，即使有保密不为人知的希望，也绝不行违背职分的事。为说明这一点，哲学家们编造了\Person{古革斯}指环的故事。在揭穿这个故事后，他引用了\Person{达味}和\Person{洗者若翰}生平中为人熟知且真实的例子。}\par

29. 我们在此预先宣布讨论的结论，仿佛已经结束一样：我们声明一条固定不变的准则，就是我们只应追求那有德性的事。智者除了能光明磊落地、毫无虚假地做成的事以外，什么也不做；即使在可能不被人察觉的情况下，他也不做任何会使自己卷入恶行的事。因为，他在自己眼中已是有罪的，然后才在别人眼中显为有罪；罪行的公开，并不比他自己意识到的更令他羞耻。我们可以用善良人士的真实榜样，而不是哲学家们所用的虚构故事，来证实这一点。\par

30. 因此，我不必想象大地因暴雨而松动，随后裂开一个大深渊，如\Person{柏拉图}所描述的那样。因为他让\Person{古革斯}下到那个深渊里，在那里遇到寓言中的铁马，马身两侧有门。这些门打开后，他发现一个死者手指上有一枚金指环，尸体躺在那里，毫无生命。他贪求黄金，便取走了指环。当他回到他所属的国王的牧羊人那里时，偶然把宝石转向掌心，他便看见一切，却无人看见他。然后当他把指环转回原位，又为众人所见。他意识到这种奇异的能力后，便利用这指环与王后通奸，杀了国王，并屠杀所有他认为应当除掉的人，免得他们成为他的障碍，然后夺取了王国。\par

31. \Person{柏拉图}说：把这个指环给一个智慧人，当他犯下错误时，可以借助它而不被人察觉；但他并不会比无法藏匿时更少受罪的玷污。智慧人的藏身处，不在于可免罚的希望，而在于他自己的清白。最后，法律并不是为义人立的，而是为不义的人立的。\BibleRef{弟前 1:9} 因为义人内心有他心中的法律，以及公平正义的规范。因此，他远离罪恶，不是由于害怕惩罚，而是由于有德行的生活规范。\par

因此，回到我们的主题，我现在要举出的不是虚假的例子来代替真实的，而是真实的例子来代替虚假的。我何须想象地上的一道裂缝、一匹铁马，以及从死者手指上发现的一枚金戒指；并说这戒指有如此力量，佩戴者可随心所欲地显现，若不想被人看见，便可从旁观者眼前消失，仿佛不在场。这个故事显然是要回答这样一个问题：一位明智的人若有机会使用这枚戒指以隐藏自己的罪行，并获得一个王国——我说，他是否会不愿犯罪，并认为罪的污点远比惩罚的痛苦更糟，还是会用它作恶，指望不被发现？我说，我又何须借用戒指的假设，既然我能用事实说明：一位明智的人，明知犯罪不仅不会被发现，还可因放纵获得王国，反之若不犯罪则有自身安危之虞，却宁愿选择冒险保全自身以远离罪行，而不愿犯罪以获得王国？\par

达味逃避撒乌耳王的面容时，因为王带着三千精兵在旷野追杀他，要将他处死，他进入王的营地，发现王正在睡觉。在那里，他不但没有伤害王，反而保护王免遭与他同来者的杀害。因为阿彼瑟对他说：\BibleQuote{今天主将你的仇敌交在你手中，现在就让我杀了他。}达味却回答：\BibleQuote{不可杀他！因为谁伸手加害主的受傅者，而能无罪呢？}达味又说：\BibleQuote{我指着永生的主起誓，若非主击倒他，或他的日子到了而死，或在战场阵亡，愿主禁止我伸手加害主的受傅者。}\par

因此，他不让人杀害他，只拿走了他头旁的长矛和水壶。随后，当所有人都熟睡时，他离开营地，走到山丘顶上，开始斥责侍卫，特别是他们的将领阿贝乃尔，没有忠实地看守他们的主人和君王。接着，他指示他们王的长矛和水壶原来在头旁的位置。当王呼唤他时，他便归还了长矛，并说：\BibleQuote{愿主按各人的正义和忠诚报答各人，因为主将你交在我手里，我却不愿报复主的受傅者。}即使他说这话时，仍害怕他的谋害而逃亡，在流亡中辗转。然而，他从未将安全置于无罪之前，因为当第二次有机会杀死君王时，他仍不使用那带给他的机会，这机会给他提供了确实的安全而非恐惧，以及王国而非流亡。\par

在若翰的事情上\BibleRef{玛 14:3}，用戒指有什么益处呢？他若保持沉默，便不会被黑落德处死。他本可以在黑落德面前保持沉默，既能被人看见，又不会被杀。但因为他不但无法容忍自己犯罪以保全自身，甚至无法忍受和容忍他人的罪，便为自己带来了死亡的原因。当然，没有人能否认他本可以保持沉默，正如在盖吉斯的故事中，人们无法否认他本可藉戒指的帮助而隐形。\par

然而，尽管那个寓言没有真理的力量，却有可取之处：如果一个正直的人能隐藏自己，他仍会避免犯罪，如同他无法隐藏自己一样；他不会藉戴上戒指来隐藏自己的身体，而是藉穿上基督来隐藏自己的生命。正如宗徒所说：\BibleQuote{因为你们已经死了，你们的生命已与基督一同藏在天主内了。}\BibleRef{哥 3:3}所以，不要有人在此力求显扬，不要有人骄傲，不要有人夸耀。基督不愿在此世间被认识，他活在世上时，不愿他的名在福音中传扬。他来是为隐藏于这个世界。因此，我们也要效法基督的榜样，隐藏我们的生命，逃避虚荣，不求被人知晓。更好是活在世上的谦卑中，而活在天上的光荣里。\BibleQuote{当基督——我们的生命显现时，那时，你们也要与祂一同出现在光荣之中。}\BibleRef{哥 3:4}\par

% structure: p0048 source=h2 target=subsection reason=relative-heading-rank; base=h2

\subsection{第六章}

\Emph{我们不应让利益观念控制我们。那些通过贩卖谷物获利的人提出什么借口，又应如何回应他们。在此摆在我们面前的是福音中的一些比喻和撒罗满的一些话语。}\par

因此，不要让利益胜过德行，而是德行胜过利益。我所说的利益是指一般人所认为的利益。要根除贪爱钱财，消灭私欲。圣人说他从未曾经商。因为谋求高价不是纯朴而是狡诈的标志。经上说：\BibleQuote{抬高粮价的人，必为民众所咒骂。}\BibleRef{箴 11:26}\par

这段陈述清楚明确，没有留下辩论余地，不像好辩的说话方式常会有的，因为有人主张农业为人人所称道；土地的出产容易种植；人播种愈多，获得的称许也愈大；再者，他勤劳工作所得丰厚的回报并非由欺诈得来，而对荒芜土地的疏忽怠慢通常应受责备。\par

39. 他说：\Quote{我辛勤地耕地；我慷慨地播种；我努力耕耘；我收获了丰厚的增产；我焦急地储藏，忠实地保存，小心地守护。如今在饥荒时期我出售，帮助饥饿的人。我卖的是自己的粮食，不是别人的。而且不比别人贵，甚至更便宜。这里有什么欺诈呢？当许多人如果没东西可买就会陷入巨大危险时，勤劳要被定罪吗？勤勉要受指责吗？远见要受辱骂吗？} 或许他还会说：\InnerQuote{\Person{若瑟}在丰年积蓄粮食，在歉收时出售。有谁被迫以高价去买？对买主使用了暴力吗？购买的机会给予所有人，没有人受到伤害。}\par

40. 当这人说完，一个人的想法把他带到如此地步时，另一个人站起来说：\Quote{农业确实是好的，因为它为所有人提供果实，以单纯的勤劳增添大地的富饶，不带任何欺骗或欺诈。若有任何错误，损失也更大，因为人播种越好，收成就越好。若他播种了纯净的麦种，他就收获更纯净、更清洁的收成。肥沃的土地以多倍的产量回报她所领受的。一块好田地连本带利地回报它的出产。}\par

41. 你必须期望从肥沃土地的庄稼中获得你劳动的报酬，必须期望从丰饶大地的出产满心中获得公正的回报。你为何利用自然的勤勉却使之成为欺骗？你为何吝于将从众人生产的供人使用？你为何减少人民的丰富供应？你为何以匮乏为目标？你为何使穷人渴望荒年？因为当他们感受不到丰收季节的益处时，由于你抬高价格，储藏谷物，他们宁愿什么也不生长，也好过你以他人饥饿为代价做生意。你把缺粮、食物供应不足看得太重。你叹息土壤丰产；你哀叹普遍丰收，悲叹仓廪充盈；你留意寻找什么时候收成差、歉收。你因灾祸对你的愿望微笑而欢喜，以便无人能有出产。然后你因你的收成来到而欢喜。那时你从所有人的苦难中积聚财富，却称此为勤劳和勤勉，而它只不过是狡诈的机敏和精明的生意诀窍。你称之为补救，然而它只是邪恶的设计。我该称之为抢劫还是仅仅获利？这些机会被抓住，仿佛劫掠的时机，在其中，你如同某个残忍的拦路强盗，扑向人们的肚腹。价格抬得更高，好像凭空加上利息，但生命的危险也增加了。因为那时储藏谷物的利息增高。你像高利贷者一样藏起谷物，又像拍卖人一样把它拿去拍卖。你为何愿所有人遭殃，因为饥荒将更加严重，仿佛不应剩下粮食，仿佛随后的年份会更不丰收？你的获利是公众的损失。\par

42. 圣\Person{若瑟}向所有人开放粮仓；他没有关闭粮仓。他没有试图获取全年出产的足价，而是定为按年支付的款项。他没有为自己取任何东西，而是在可能遏制未来饥荒的范围内，以谨慎的远见作出了安排。\par

43. 你们读过主耶稣在福音中怎样论及那个寻求高价的粮商，他的财产为他带来丰盛的果实，但他却好像仍处匮乏，说：\BibleQuote{我要怎么办呢？我没有地方收藏我的产物。我要拆毁我的仓房，另建更大的。}\BibleRef{路 12:17}尽管他不知道自己当晚可能被索命。他不知如何是好，似乎处在疑惑中，好像缺粮似的。他的仓房容纳不下一年的收成，他却以为自己有所匮乏。\par

44. 因此，\Person{撒罗满}说得对：\Quote{那屯积粮食的，必将粮食留给万民，} 不是留给他的后嗣，因为贪婪所得的利益与继承权毫不相干。凡非正当地聚集的，必如被风刮走一般，被外人夺取。他又加上一句：\Quote{那贪婪丰收年产物的人，在人民中是受诅咒的；但慷慨分施的人却蒙祝福。} 那么你看到，所说的是关乎施舍粮食的人，而不是寻求高价的人。因此，真正的利益并不存在于德行所失大于利益所得之处。\par

% structure: p0058 source=h2 target=subsection reason=relative-heading-rank; base=h2

\subsection{第七章。}

\Emph{在饥荒时期，绝不可将外方人驱逐出城。这里引述一位基督徒贤人的高贵忠告，与此相对比的是在\Place{罗马}发生的可耻行为。通过比较二者，显示前者既有德又有益，而后者则两者皆无。}\par

45. 然而，那些想要禁止外方人入城的人，也不能得到我们的赞同。他们将在该提供帮助之时驱逐他们，并且将他们隔离于共同母城的交往之外。他们拒绝让他们分享为众人准备的出产，并回避已经开始的往来；他们不愿在急需时分，将他们自己曾共同享有权利的东西，与他们分享。走兽不驱赶走兽，然而人却排斥人。野兽和动物认为大地供给的食物是众人共有的。它们都帮助同类；而人，本该视一切与人类有关之事为己任，却与自己的同类为敌。\par

46. 有一个人，年事已高，当城市遭受饥荒时，如同常有的情况，民众要求禁止外乡人入城。他担任市府最高长官之职，召集了官员和富户，要求他们为公众福祉商议。他说，驱逐外乡人，犹如一人被另一人抛弃，垂死时得不到食物，同样残忍。我们不让自己的狗上桌而任其挨饿，却把人拒之门外。再者，如此多人死亡，为世界是何等的损失，致命的瘟疫夺去他们的生命。为他们的城市，如此众多的人丧亡，又是何等的损失，他们本可以帮助纳税或经营生意。他人的饥饿对任何人都无益，也不可尽量拖延援助的时日，而对解决匮乏无所作为。更有甚者，当如此多的土地耕种者离去，如此多的劳动者死去，未来的粮食供应必将断绝。难道我们要驱逐那些惯常供给我们食物的人吗？在急需的时候，我们不愿意供养那些一直供养我们的人吗？即使现在，他们提供的帮助是何等的大。\BibleQuote{人生活不只靠饼。}\BibleRef{申 8:3} 他们甚至是我们自己的家人，其中许多人甚至是我们的亲属。让我们对所领受的有所回报。\par

47. 但也许我们害怕贫困会增加。首先，我回答，慈惠绝不会落空，总能找到援助的方法。其次，我们可以通过认捐来弥补将给予他们的粮食供应。我们可以用黄金来补足。再者，如果我们失去这些耕种者，难道我们不必购买别的耕种者吗？供养一个工人比购买一个工人便宜多少呢？而且，从哪里能够找到接替前者的工人呢？即使找到一个，不要忘记，对于一个习惯于不同方式的生手，我们可能在人数上补足，但在需要完成的工作上却不行。\par

48. 我还需要多说什么呢？当款项募集起来时，粮食就运来了。因此城市的丰裕并未减少，而外乡人也得到了救助。这件事使那位圣人从天主那里获得了何等的嘉许！在人间又获得了何等的荣耀！他确实赢得了显赫的名声，他能指着全行省的人民，向皇帝据实禀报：\Quote{所有这些人都是我为你保全下来的；他们因元老院的仁慈而活着；你的议会把他们从死亡中抢救了出来！}\par

49. 这种作法比起最近在罗马发生的事，要有利得多。在那里，那些一生大部分时间已在城中度过的人，竟被从那辽阔的城市驱逐。他们流着泪，带着儿女离去；他们为儿女作为公民而哀叹流亡，他们说这种流亡本应避免；他们同样因爱的联结断绝，亲缘关系被割裂而悲伤。然而，丰收的年成曾向我们微笑。只有这座城市需要运进粮食。它本可以获得帮助，如果它向那些它正将他们的子女赶走的意大利人寻求粮食的话。没有什么比把一个人当作外国人驱逐，却又好像他属于我们一样索取他的服务，更为可耻。你怎能驱逐一个靠自己的出产为生的人呢？你怎能驱逐那供养你食物的人呢？你留下你的仆人，却赶走你的亲属！你夺走粮食，却不表示好感！你强行取食，却不表示感谢！\par

50. 这是何等的可悲，何等的无益！因为不合宜的，怎能有利益呢？罗马有时缺乏来自其社团的何等巨大的供应，然而她不能解散它们而仍能避免饥荒，同时等待有利的风向和期望中的船只运来的粮食。\par

51. 前述的管理方式是多么更有德、更有益啊！因为有什么比当贫乏者得到富人的赠予而得助，饥饿者得到食物，每日的口粮没有缺乏，更为合宜或有德呢？有什么比当耕种者得以留住于土地，乡村居民不致丧亡，更为有利呢？\par

52. 凡合乎德行的，也是有利的；凡有利的，也是合乎德行的。反之，凡不利的，便是不合宜的；凡不合宜的，也是不利的。\par

% structure: p0068 source=h2 target=subsection reason=relative-heading-rank; base=h2

\subsection{第八章}

\Emph{那些将德行置于利益之上的人蒙天主悦纳，可由若苏厄、加肋布和其他探子的榜样得到证明。}\par

53. 我们的祖先何时能摆脱奴役，除非他们相信服事埃及王不但是可耻的，而且是无益的呢？\par

54. 若苏厄和加肋布，奉命去侦探那地时，带回消息说那地确实肥沃，但住着极为凶猛的民族。\BibleRef{户 13:27} 民众害怕战争，拒绝占领那地。若苏厄和加肋布，被派作探子，力劝他们那地是肥沃的。他们认为在异民面前退缩是不合宜的；他们宁愿被石头砸死，这正是民众所威胁的，也不愿放弃他们合乎德行的立场。其余的人不断劝阻，民众反对，说他们要与残忍可怕的民族作战；他们将在战斗中阵亡，而他们的妻子儿女将沦为战利品。\BibleRef{户 14:3}\par

55. 主的怒气发作，要击杀所有人，但藉着\Person{梅瑟}的祈祷，祂缓和了自己的判决，推迟了报复，因为祂知道，祂已经充份惩罚了那些没有信德的人，即使祂暂时饶恕他们，不杀死那些不信的人。然而，祂说，他们不得进入他们所拒绝的那地，作为他们不信的惩罚\BibleRef{户 14:29}；但他们的子女和妻子，他们没有抱怨，并且因着性别和年龄是无辜的，应领受那应许之地的产业。因此，二十岁以上的人，他们的尸体倒在旷野。其余人的惩罚被搁置。但是，那些与\Person{若苏厄}一同上去，并且认为应当劝阻百姓的人，立刻死于一场大瘟疫\BibleRef{户 14:37}。\Person{若苏厄}和\Person{加肋布}\BibleRef{苏 14:6}与那些因年龄或性别而无罪的人一同进入了应许之地。\par

56. 因此，较好的一方选择光荣胜于安全；较差的一方选择安全胜于德行。但天主的判断赞许那些认为德行高于利益的人，而谴责那些宁愿选择看似更符合安全而非德行之事的人。\par

% structure: p0074 source=h2 target=subsection reason=relative-heading-rank; base=h2

\subsection{9}

\Emph{欺骗和不诚实的赚钱方式，对于职责是服务众人的圣职人员是完全不适宜的。他们绝不应涉足金钱事务，除非是关系到人的生命。他们被给予\Person{达味}的榜样，即他们不应伤害任何人，即使在受到挑衅时；还有\Person{纳波特}的死，为阻止他们选择生存而舍弃德行。}\par

57. 一个人若不喜爱有德行的生活，反而被卑劣的生意所刺激，从事低下的贸易，或因贪婪之心燃烧，日夜急切地损害他人的财产，不将灵魂提升至有德生活的光辉，也不重视真正赞美的美，这是最为可憎的。\par

58. 由此便产生了藉巧言寻求继承权，以及佯装节制和庄重而取得的继承。但这与基督徒的观念完全相悖。因为一切藉诡计获得、由欺骗积聚的东西，都丧失了坦诚的优点。即使在未担任圣职的人员中，寻求他人的继承也被视为不妥。那些生命行将结束的人，应该运用自己的判断，使他们能按自己认为最好的方式自由地立遗嘱，因为他们以后将无法修改。因为将属于他人或为他们积攒的财物据为己有，是极不光彩的。此外，司铎或圣职人员的职责是尽可能对所有人有益，对任何人无害。\par

59. 若帮助一个人就必然伤害另一个人，那么宁可都不帮助，也不要压迫一方。因此，干预金钱事务并非司铎的职责。因为在此情形下，往往败诉的一方会遭受损害，然后他会认为自己的失败是由于中间人的干预。司铎的职责是不伤害任何人，随时准备帮助所有人。但能够做到这一点，唯独属于天主的能力。在生死攸关的事上，若伤害了那在危险中本应帮助的人，无疑是严重的罪过。但为了人的安全，虽常需承担重大的辛劳，而介入金钱事务却招致他人的憎恨，则是愚昧的。在这种情形下，冒险是光荣的。所以，司铎的职责应坚定地持守这一点：不伤害任何人，即使在受到挑衅和因某些伤害而痛苦时。那说\BibleQuote{如果我以恶报善}的人，真是善良的。因为如果我们不伤害那没有伤害过我们的人，这算得上什么光荣呢？但受了伤害而能宽恕，才是真正的德行。\par

60. 这是多么有德行的行为，当时\Person{达味}宁愿宽恕那作他仇敌的君王，尽管他本可以伤害他！这又是多么有益，因为它在他继承王位时帮助了他。因为所有人都学会了忠于他们的君王，不篡夺王位，反而敬畏尊崇他。这样，德行就比利益更受重视，然后利益也随着德行而来。\par

61. 但他宽恕了他，这还算小事；他更为他在战争中被杀而悲伤，流泪哀悼说：\BibleQuote{基耳波亚山啊，愿你既无雨露，也无甘霖；死亡的山啊，因为英雄的盾牌在那里被弃，\Person{撒乌耳}的盾牌，不以油敷抹，却以伤者的血和勇士的脂油。\Person{约纳堂}的弓绝不后退，\Person{撒乌耳}的剑不空空而回。\Person{撒乌耳}和\Person{约纳堂}，可爱可亲，生不离，死不散。他们比鹰更迅速，比狮还强健。\Place{以色列}的女儿们，为\Person{撒乌耳}悲哭吧！他曾使你们衣锦穿红，在你们的衣裳上缀满金饰。英雄们怎会仆倒在战斗中！\Person{约纳堂}竟在至高之处被杀！我的兄弟\Person{约纳堂}，我为你悲痛万分，你为我异常可爱。你的爱对我宛如妇女的爱。英雄们怎会仆倒，可珍的武器怎能丧亡！}\par

62. 有哪一位母亲能为她的独生子如此哭泣，就像他在这里为他的仇敌哭泣一样？有谁能以如此赞美来追随他的恩人，就像他追随那谋害他性命的人一样？他是多么深情地悲伤，又是何等深切地哀悼他！山岭因先知的诅咒而干旱，一种天主的能力充满了他所说出的判决。因此，自然界本身因目睹了君王的死亡而受到了惩罚。\par

63. 在圣善的\Person{纳波特}身上，除了他对圣善生活的重视之外，是什么导致了他的死亡呢？因为当国王向他要葡萄园，并答应给他钱时，他拒绝了为他祖业而设的价格，认为这是不体面的，宁愿以死来逃避这种耻辱。\BibleQuote{上主决不许我做这事，将我先人的产业给你！}也就是说，这样的耻辱不要临到我身上，愿天主不允许这种邪恶用暴力得逞。他并不是在谈论葡萄树——天主并不关心葡萄树或田地——他是在谈论他祖先的权利。他本可以接受另一个葡萄园或国王的葡萄园，并成为国王的朋友，这在人们看来，就现世而言，并非没有不小的益处。但因为这行为卑鄙，他认为这不可能是真正有益的，于是宁愿在荣誉无损的情况下承受危险，也不愿通过使自己蒙羞来获取所谓的利益。我在此再次谈论通常所理解的利益，而不是那存在于圣善生活恩宠中的利益。\par

64. 国王本可以自己强行夺取，但他觉得那样太过无耻；后来\Person{纳波特}死了，他感到悲伤。主也宣布那女人的残忍将受到相称的惩罚，因为她不顾德行，宁愿要可耻的获利。\par

65. 每一种不公正的行为都是可耻的。即使在日常事务中，虚假的砝码和不公道的量器也是可咒骂的。如果市场中或生意上的欺诈都会受到惩罚，那么当它出现在履行德行之职责中时，难道能免于指责吗？\Person{撒罗满}说：\BibleQuote{不同的两样砝码，不同的两样升斗，二者同为上主所憎恶。}\BibleRef{箴 20:10}在此之前也说：\BibleQuote{骗人的天秤，上主憎恶；法码准确，他才喜悦。}\BibleRef{箴 11:1}\par

% structure: p0085 source=h2 target=subsection reason=relative-heading-rank; base=h2

\subsection{第十章}

\Emph{不仅在世俗法律中，而且在圣经中，我们都被警告在一切协议中避免欺诈，正如\Person{若苏厄}和基贝红人的例子所清楚表明的。}\par

66. 因此，在一切事上，忠信是合宜的，正义是可悦的，公平的适度是令人喜乐的。但对于合同，尤其是土地出售、协议或盟约，我该说什么呢？难道不是有规则正是为了杜绝一切虚假的欺骗，并使那被发现的欺骗者受到加倍的惩罚吗？因此，在一切地方，对德行的重视都居首位；它排除欺骗，驱除欺诈。所以先知\Person{达味}恰当地概述了自己的判断，说：\BibleQuote{他对近人没有作过恶事。}因此，欺诈不仅应该从合同中消失（在合同中，出售物品的缺陷必须予以记录；除非卖方已提及缺陷，否则即使已将货物完全转让给买方，也会因欺诈之诉而被宣布无效），而且在其他一切事上也应如此。必须显出坦率，必须显明真相。\par

67. 圣经在归于\Person{农的儿子若苏厄}名下的那部旧约\BibleBook{若苏厄}{书}中，已经清楚陈述了（其实并非法学家的法律规则，而是）古圣祖们关于欺骗的判断。当消息在各民族中传开，说希伯来人过海时海水干涸了；水从岩石中流出；每天有从天上降下的食物，数量足够如此成千上万的人民食用；\Place{耶里哥}的城墙在圣号角声中倒塌，被民众的呼喊声所摧毁；此外，\Place{哈依}王被打败，并被挂在树上直到晚上；那时，基贝红人因惧怕他的强力，便带着诡计前来，假装他们来自极远之地，因长途跋涉而鞋破衣烂，还出示了衣物变旧的证据。他们说，他们之所以忍受如此的劳苦，是因为渴望与希伯来人达成和平并建立友谊，于是便开始恳请\Person{若苏厄}与他们结盟。而他，由于对各地尚不了解，对居民也一无所知，未能识破他们的欺骗，也没有求问天主，竟轻易相信了他们。\par

68. 在那个时代，人们如此看重自己许下的诺言，以至于无人相信别人会试图欺骗。谁又能因这一点责难圣人们呢？——他们以为别人与自己怀有同样的情操，认为无人会说谎，因为真理是他们自己的伴侣。他们不知道什么是欺骗，他们乐于相信别人也如自己一样，因为他们无法怀疑别人会是自己所不是的那种人。因此\Person{撒罗满}说：\BibleQuote{幼稚的人，有话必信。}\BibleRef{箴 14:15}我们不应指责他轻易相信，而应称赞他的善意。对任何可能伤害他人的事一无所知，这正是无辜者的特征。虽然他被别人欺骗，但他仍然对一切人都怀着善意，因为他以为一切人都有忠信。\par

69. 因此，受这些考虑所诱导，他相信了他们，与他们立了约，赐予他们和平，并与他们联合。但当他们到了他们的地方，欺骗被发现了——因为他们虽然住得很近，却假装是陌生人——我们祖先的民众便开始因受欺骗而发怒。然而，\Person{若苏厄}认为已订立的和平不能破坏（因为它已由誓言所确认），以免在惩罚别人的背信时，自己变成了违背承诺之人。不过，他让他们付出代价，强迫他们从事最低贱的劳作。这个判决固然温和，却是持久的，因为在他们世代相传的劳役中，一直到今天\BibleRef{苏 9:27}，都存留着对他们古代诡诈的惩罚。\par

% structure: p0091 source=h2 target=subsection reason=relative-heading-rank; base=h2

\subsection{第十一章}

\Emph{在引述了若干存在于某些修辞学家作品中关于欺诈的例子之后，他指出这些例子以及所有其他例子，在圣经中都得到了更充分且更明确的谴责。}\par

70. 对于继承人接受遗产时弹指作响或赤身跳舞之事，我将闭口不提。这些都是众所周知之事。我也不会谈论为激起买主欲望而假借一次捕捞聚集大量鱼群的事情。他为何如此贪恋奢华与美味，以至于纵容这种欺诈行径呢？\par

71. 我何必述说那桩众所周知的、关于\Place{锡拉库萨}一处惬意幽静别墅和一位西西里人狡诈的故事呢？他遇到一个外地人，知道他想购置地产，便请他到自己的庄园赴宴。外地人答应了，次日便去。在那里，他看到众多渔人，以及一场摆满珍馐的盛宴。当着宾客的面，渔人被安置在花园空地上，那里此前从未撒过网。每人都依次向宾客呈上所获，鱼被摆上桌，进入就座者的眼帘。外地人对鱼的数量众多和船只之多感到惊讶。得到的答复是，此处水源丰沛，因水味甘甜，故有大量鱼群来此。简言之，他诱使外地人急切地想得到这块地产，自己则假装被说服，似乎不情愿地收下了款项。\par

72. 次日，买主与朋友们来到庄园，却发现一条船也没有。他问是否渔人们那天在过节，却被回答说，除了昨日，他们从来不在那里打鱼；但他又能凭什么去追究这种欺诈呢？他曾经那般可耻地贪图奢华。因为指控别人有过错的人，自己应该无可指摘。因此，我不会将这类琐事纳入教会的惩戒权限，因为教会彻底谴责一切贪图不义之财的欲望，并用寥寥数语简短扼要地禁止一切狡猾奸诈的行为。\par

73. 我还要说什么呢？有人凭借一份遗嘱声称自己是继承人或受遗赠人，而这份遗嘱虽然是被他人伪造的，但他却明知此事，并试图通过他人的罪行获利。即使世俗法律也判定，明知故犯地使用假遗嘱的人犯有恶行。但正义的法则却是清楚明白的：善人不应该偏离真理，不应该使任何人蒙受不义之损，也不应以任何方式行事诡诈或参与任何欺诈。\par

74. 然而，论及此事，还有什么比\Person{阿纳尼雅}的例子更清楚呢？他对自己出卖土地所得的价格行为不实，因为他在出卖后，把价银的一部分放在宗徒们脚前，却假称是全部。\BibleRef{宗 5:2}为这欺诈之罪，他灭亡了。他本可以什么都不奉献，那样做并不会犯下欺诈之罪。但既然他的行为中搀进了欺骗，他非但没有因慷慨而获得恩宠，反而为自己的诡计付出了代价。\par

75. 主在福音中也曾拒绝那些心怀诡诈来到祂面前的人，说：\BibleQuote{狐狸有穴，}\BibleRef{玛 8:20}因为祂命令我们生活在简朴和内心的清白之中。\Person{达味}也说：\BibleQuote{你以欺诈为利刃，}以此指出奸诈之辈，正如这种器具虽常用于助人装扮，却往往割伤人。如果有人像叛徒那样假装示好而策划阴谋，以致将本该保护的人置于死地，那么就当视他如同那类因心醉神迷、手抖失控而易于伤人的器具。因此，那被邪恶之酒灌醉的人，通过可怖的背信弃义行为，使大司祭\Person{阿希默肋客}丧命，因为当君王被嫉妒的刺所激怒而追击先知时，阿希默肋客曾善意款待了先知。\par

% structure: p0099 source=h2 target=subsection reason=relative-heading-rank; base=h2

\subsection{第十二章}

\Emph{我们不可许下错误的诺言；而且如果发了不义的誓言，也不可遵守。可证\Person{黑落德}在此事上犯了罪。\Person{依弗大}所作的誓愿应受谴责，同样，凡天主不愿我们从祂那里领受的任何其他誓愿，都应受谴责。最后，将依弗大的女儿与两位毕达哥拉斯派人士比较，且将她置于他们之上。}\par

76. 人的心性应当无玷而健全，使他能出言无诡诈，并守身于圣化之中；\BibleRef{得前 4:6}使他不以虚伪之言欺骗兄弟，也不许下任何可耻之事。若他既已许下这样的诺言，那么宁可不去履行，也远比履行那羞耻之事要好。\par

77. 人们常常用庄严的誓言约束自己，虽然后来知道不该许下那诺言，却因顾及誓言而履行。这正是我们先前所提及的\Person{黑落德}的做法。他为了一个舞女而许下可耻的赏赐承诺——并残酷地履行了。说这可耻，是因为竟为一个舞蹈而许下王国；说它残酷，是因为一位先知的首级竟为了一个誓言而被牺牲。若论信守这样的誓言，还不如背誓来得好——如果这真能称作背誓的话，因为那只是一个醉汉在酒宴中、或一个娇纵放荡之徒在舞蹈进行中所许下的诺言。先知的首级被盛在盘子里端来，\BibleRef{谷 6:28}而这事竟被视为守信之举，实则却是疯狂之举！\par

78. 我绝不会相信，领袖\Person{依弗大}许下誓愿时是经过深思熟虑的；他许愿说，当他平安归来时，无论谁先从他的家门出来迎接他，他就要将那人献于天主。因为他后来懊悔了，当他的女儿出来迎接他时。他撕裂衣服，说：\BibleQuote{唉！我的女儿，你使我十分愁苦，你成了我的祸根。}\BibleRef{民 11:35} 虽然他怀着虔诚的敬畏，承担起残酷任务的痛苦履行，但他仍规定并留传给后世，每年要有一段哀悼和悲痛的日子。这是一个严苛的誓愿，但它的履行更加痛苦，那履行的人最有理由悲伤。从此，这在\Place{以色列}中成了年年遵循的规矩和定例，正如经上所说：\BibleQuote{以色列的女子每年去哀悼基肋阿得人\Person{依弗大}的女儿四天。}\BibleRef{民 11:40} 我不能责怪这人认为必须履行他的誓愿，但这确实是一个可悲的必须，只能以他孩子的死亡来解决。\par

79. 与其许下天主不愿接受的愿，不如不许愿。在\Person{依撒格}的事上，我们有一个例子，因为主指定了一只公绵羊代替他献上。\BibleRef{创 22:13} 因此，并非所有的许诺都必须履行。实际上，主自己常常改变祂的决定，正如圣经所指出的。在称为\BibleBook{}{户籍纪}的书中，祂曾宣告要处以死刑惩罚百姓并消灭他们，\BibleRef{户 14:12} 但后来，因\Person{梅瑟}的恳求，祂又与百姓和好。又一次，祂对\Person{梅瑟}和\Person{亚郎}说：\BibleQuote{你们离开这会众，我好在顷刻之间消灭他们。}\BibleRef{户 16:21} 当他们离开会众时，大地忽然裂开，开口吞下了\Person{达堂}和\Person{阿彼兰}。\par

80. \Person{依弗大}女儿的例子，远比那两个\OriginalTerm{毕达哥拉斯学派}{Pythagoreans}的例子更光荣、更古老，后者在哲学家中间被视为如此著名。其中一人，被暴君\Person{狄奥尼修斯}判处死刑，并确定了行刑的日子，他请求准许回家，以便安顿家人。但为了担心他可能失信不归，他提供了自己死亡的担保，条件是他若在指定日期缺席，担保人就准备好代他受死。另一人没有拒绝提出的担保条件，并以平静的心等待那死亡的日子。因此，一个人没有逃避，另一个人则在指定日期返回。这一切显得如此奇妙，以致那暴君寻求与他们为友，尽管他曾急于要毁灭他们。\par

81. 那么，在受人尊敬和有学问的人身上被视为充满奇迹的事，在一名贞女身上却显得更加光辉、更加荣耀，正如她对悲伤的父亲所说：\BibleQuote{就照你口中所出的对待我吧！}\BibleRef{民 11:36} 但她请求延迟两个月，好让她和同伴们上山去，合宜而虔诚地哀悼她那将奉献于死的贞洁。同伴们的哭泣没有打动她，她们的悲伤没有说服她，她们的哀号没有挽留住她。她不让日子过去，也不让时辰错过。她回到父亲那里，仿佛是随自己心愿回来的，并在父亲犹豫时用自己的意志催促他，如此出于自由选择而行动，因此起初那可怕的事变成了虔诚的祭献。\par

% structure: p0107 source=h2 target=subsection reason=relative-heading-rank; base=h2

\subsection{第十三章}

\Emph{\Person{友弟德}，在为美德之故忍受了许多危险之后，获得了非常多且巨大的益处。}\par

82. 看！\Person{友弟德}向你们展现自己，值得钦佩。她接近\Person{敖罗斐乃}——一个被人民畏惧的人，被胜利的\Place{亚述}军队所围绕。她首先以形体的优雅和面容的美丽打动了他。然后以谈吐的优雅诱使他入彀。她的第一次胜利是：她从未受玷污的洁净中从敌人的营帐返回。\BibleRef{友 12:20} 她的第二次胜利是：她战胜了一个男人，并以她的计谋使那人民溃散。\par

83. \Place{波斯}人对她的胆量感到惊恐。因此，那两个\OriginalTerm{毕达哥拉斯学派}{Pythagoreans}中人身上为人所钦佩的，在她身上也值得我们的钦佩，因为她不怕死亡的危险，甚至也不怕她的贞洁所面临的危险，这对好女人来说是更为关切的事。她不怕一个恶徒的击打，甚至也不怕整支军队的武器。她，一个女人，站在战斗者行列之间——就在胜利的武器之中——不顾死亡。人若看她那巨大的危险，会说她是出去送死；人若看她的信德，又会说她是出去作战。\par

84. \Person{友弟德}当时追随了美德的召唤，而随着她追随美德，她赢得了巨大的益处。阻止主的子民将自己交给异教徒，是有德行的；阻止他们背叛本族的礼仪与奥秘，或阻止他们将奉献的贞女、可敬的寡妇和端庄的妇女交给蛮族的污秽，或阻止他们以投降来结束围城，也是有德行的。她甘愿为所有人遭遇危险，以便将所有人从危险中解救出来，这更是有德行的。\par

85. 她的美德多么强大有力，以致她，一个女人，竟能在最重大的事情上提议出谋划策，而不将其留给人民的领袖！同样，她的美德多么强大有力，竟能确定无疑地指望天主助佑她！她的恩宠多么宏大，竟能得到天主的助佑！\par

% structure: p0113 source=h2 target=subsection reason=relative-heading-rank; base=h2

\subsection{第十四章}

\Emph{\Person{厄里叟}所做的如何有德行且有益处。这与\Place{希腊}人常提及的事迹相比。\Person{若翰}为了美德献出了生命，\Person{苏撒纳}为了同样的理由将自己置于死亡的危险之中。}\par

86. 厄里叟所追随的，舍德行其谁？他使来捉拿他的叙利亚军队眼目昏迷，将他们领进撒玛黎雅。然后他说：\Quote{主啊，求你开启他们的眼睛，使他们看见。}他们就看见了。但以色列王想要杀掉这些进来的人，并请求先知准许他如此行，他却回答说：这些人并非因手力或兵器被俘，故不可杀戮，反而应当供食相助。于是他们饱餐足食。那些叙利亚强盗此后以为，他们绝不可再踏上以色列的土地。\par

87. 此事比希腊人昔日所行，高贵何止千万！当时两国为争荣耀与霸权相互争斗，其中一国有机会暗中烧毁对方战船，他们却以此为耻，宁愿以荣誉获取较小优势，也不以可耻方式获取较大优势。他们果真不能如此行事而不自取其辱，也不能以此计诱骗那些为结束波斯战争而结盟的人。即便可以口头否认，他们却总难免面红耳赤。然而，厄里叟却愿意拯救，而非毁灭那些被迷惑的人——他们固然受骗，却非因任何卑鄙行为——他们是被主的大能击打而致盲的。因为放过敌人、赐予对手性命实为合宜，尤其当他本可不放过而取其性命之时。\par

88. 由此可见，凡是合宜的，时常有益。圣善的友弟德因顾及体面而将自身安危置之度外，终结了围城之危，并以自己的德行赢得了对众人共同的利益。厄里叟因宽恕而获得的名声，远胜于杀戮所能取得的，而且他保全被俘之敌，使之用于更大的利益。\par

89. 若翰心中所怀的，无非是德行，因此他不容忍国王的邪恶结合，说道：\BibleQuote{你不可娶她为妻。}\BibleRef{玛 14:4}他本可沉默不语，但他觉得，若因怕死而不说实话，或使先知职务屈从于国王，或陷于谄媚，都是不合宜的。他明知自己反对国王必将丧命，却宁要德行而不要安全。然而，有什么比这苦难更为有益呢？这苦难给圣人带来了荣耀。\par

90. 圣善的苏撒纳也一样，当她受到假见证的威胁，看到自己一面面临危险，一面面临羞辱时，她宁愿以有德的死来避免羞辱，也不愿在求生的欲望中忍辱偷生。因此，当她定志于德行时，她也保全了性命。然而，她若选择那看似有益的自保之路，就绝不会获得如此大的名声，不仅如此，那可能不仅无益，甚至有害——她可能甚至难逃其罪行的惩罚。因此可见，凡是可耻的，都不能是有益的；反之，凡是有德的，也不能是无益的。因为利益始终是德行的双倍，德行亦然。\par

% structure: p0120 source=h2 target=subsection reason=relative-heading-rank; base=h2

\subsection{第十五章}

\Emph{在提及罗马人的一项高贵的行动之后，作者从梅瑟的事迹表明，他极其重视德行。}\par

91. 据说一位罗马将军曾有过令人难忘的事迹：当敌国国王的医生前来，答应给他毒药时，他却将医生捆绑，遣返给敌人。诚然，在投身权力角逐之后，拒绝以卑鄙手段获取胜利，实为高贵之举。他不认为德行在于胜利，反而宣称，若非以荣誉赢得的胜利，便是可耻的。\par

92. 让我们回到我们的英雄梅瑟，以及更为崇高的事迹上，以显示它们既较高，也较早。埃及王不肯放我们祖先的百姓离去。于是梅瑟命司祭亚郎将手杖伸向全埃及的水上。亚郎伸出手杖，河水立时变成了血。\BibleRef{出 7:19}无人能饮这水，所有埃及人都渴得要死；然而祖先们却有纯净的水丰涌而出。他们向天撒灰，人身和牲畜身上便起了疮疖和起泡的疮。\BibleRef{出 9:10}他们降下夹杂着火焰的冰雹，地上的一切都被摧毁。\BibleRef{出 9:23}梅瑟祈祷，一切便恢复了原先的美丽。冰雹止住，疮疖痊愈，河流恢复了平常的供水。\BibleRef{出 9:29}\par

93. 然后，再次，浓密的黑暗笼罩大地三日之久，因为梅瑟举手展开了黑暗。\BibleRef{出 10:22}埃及所有的长子都死了，而希伯来人的所有后代却毫发无损。\BibleRef{出 12:29}人们请求梅瑟结束这些恐怖之事，他祈祷，并蒙应允。一方面，他克制自己，不参与欺诈，这值得称赞；另一方面，他本其内在的良善，使天命的惩罚远离敌人，这也值得注意。他实在是，如经上所载，温柔和谦卑的。\BibleRef{户 12:3}\par

94. 他将手杖掷于地上，杖便变成了蛇，吞噬了埃及的众蛇；\BibleRef{出 7:12}这预表\OriginalTerm{圣言}{Word}应成为肉身，以罪过的宽恕和赦免，来消灭可怖之蛇的毒液。因为手杖代表真实、尊贵、充满能力和统治荣耀的圣言。杖变成了蛇；同样，由圣父所生的圣子，成为由女人所生的人子，并如同蛇被高举在十字架上，将祂的良药倾注在人类的创伤上。因此主亲口说：\BibleQuote{正如梅瑟在旷野里举起了蛇，人子也当被举起来。}\BibleRef{若 3:14}\par

95. 再者，\Person{梅瑟}给出的另一个征兆也指向我们的\Person{主耶稣基督}。他把手放入怀中，再抽出来，手便变得如同雪一样白。他再次把手放入怀中再抽出来，手又恢复了人肉的样子。\BibleRef{出 4:6} 这首先预示了\Person{主耶稣}天主性的原始光荣，然后预示了祂摄取我们的肉躯；在这真理上，天下万民万邦都必须信从。所以他放入自己的手，因为\Person{基督}就是天主的右手；凡不信祂的天主性与降生成人的人，必像罪人一样受罚；就像那位国王，他虽不信明显而清楚的征兆，但后来受罚时，却祈求蒙受怜悯。由此可见，\Person{梅瑟}对德行的重视是何等之大，尤其从他甘愿为人民舍己，祈求天主饶恕人民，否则就把自己从生命册上抹去这一事实，更可显明。\BibleRef{出 32:32}\par

% structure: p0127 source=h2 target=subsection reason=relative-heading-rank; base=h2

\subsection{第十六章}

\Emph{在略述\Person{托彼特}之后，他证明\Person{辣古耳}在德行上胜过哲学家。}\par

96. \Person{托彼特}也在生活中清晰地彰显了真正的德行：他离开宴席去埋葬死者，\BibleRef{多 2:4} 并邀请穷人到自家简陋的餐桌上共享膳食。而\Person{辣古耳}是更为光辉的榜样。他珍视德行，当有人请求娶他女儿时，他并未隐瞒女儿的缺陷，以免以沉默骗取求婚者。因此，当\Person{托彼特}——\Person{多俾亚}的儿子——请求将女儿嫁给他时，他回答说，按法律她应嫁给近亲，但他已经将她许过六个男人，全都死了。\BibleRef{多 7:11} 这位义人，真是为他人担忧胜过为自己，宁愿女儿不嫁，也不愿他人因与她结合而遭遇危难。\par

97. 他何等简明地解决了哲学家们的一切难题！他们讨论房屋的瑕疵，是应该隐瞒还是由卖主声明。\Person{辣古耳}却十分确定，女儿的缺点不应隐瞒。而且，他本非急着要把女儿嫁出——是别人来求娶的。毫无疑问，他的行为远比那些哲学家高尚，因为我们知道，比起单纯的金钱交易，女儿的未来要远为重要得多。\par

% structure: p0131 source=h2 target=subsection reason=relative-heading-rank; base=h2

\subsection{第十七章}

\Emph{古时的先祖们在即将被掳时，怀着何等的美德情感藏匿了圣火。}\par

98. 让我们再来思考那个发生在被掳时期的举动，它达到了德行与光荣的顶峰。美德不受任何逆境的阻挡，它在逆境中兴起，在逆境中胜过在顺境中。镣铐或刀剑之中，烈火或奴役之中（奴役比任何刑罚更让自由人痛苦），濒死者的痛苦之中，祖国的灭亡之中，生者的恐惧之中，被杀者的鲜血之中——在这一切之中，我们的先祖从未停止对美德的关切与深思。在已覆灭的祖国的灰烬与尘土中，它以虔诚的努力发出炽热的光芒，灿烂闪耀。\par

99. 当我们的祖先被掳往\Place{波斯}时，\BibleRef{加下 1:19} 有些当时在全能天主前服务的司铎，悄悄把从主的祭台上取来的火藏在山谷里。那里有一口敞开的深井，里面没有水，而且人迹罕至，无人知晓，也没有人闯入。他们在那里以神圣的印记秘密地封存了隐匿的圣火。他们并不急于埋藏金子或藏匿银子以留给子孙，而是在本身巨大的危难中，心怀一切美德，认为必须保全圣火，免得不洁之人亵渎它，免得被杀者的鲜血熄灭它，免得悲惨废墟的瓦砾掩埋它。\par

100. 于是他们前往\Place{波斯}，仅在信仰上是自由的；因为那是唯一不能被敌人掳去的东西。过了许久之后，天主出于祂的美意，感动\Place{波斯}王的心，命他重建\Place{犹太}的圣殿，并在\Place{耶路撒冷}恢复常规的礼仪。为执行此事，\Place{波斯}王委派了司铎\Person{乃赫米雅}。他带上了那些司铎的孙子们，这些司铎就是在离开故土时曾藏匿圣火以保全它不灭的那些人。然而，当他们抵达时——据列祖的历史记载——他们找到的不是火，而是水。当需要火在祭台上燃烧时，司铎\Person{乃赫米雅}命他们打水来，把水交给他，并将水洒在木柴上。这时，啊，何等奇妙的景象！虽然天空原本密布乌云，忽然太阳照耀，一团大火燃起，众人见到如此清晰的主的恩惠标记，无不充满喜乐。\Person{乃赫米雅}祈祷；当祭献完成，司铎们向天主咏唱赞美诗。\Person{乃赫米雅}又命人将剩下的水浇在大石上。水一浇上，火焰便迸发，而从祭台上发出的光芒比先前更灿烂辉耀。\par

101. 这个征兆传开后，\Place{波斯}王下令，在当初藏火后来发现水的地方修建一座圣殿，并有许多人向那里奉献礼物。与圣\Person{乃赫米雅}在一起的人称那地方为\OriginalTerm{乃弗塔尔}{Naphthar}，\BibleRef{加下 1:36} ——意为\Quote{洁净}——也有人称之为\OriginalTerm{乃弗泰}{Nephi}。在先知\Person{耶肋米亚}的历史中也可以看到，他曾吩咐后来的人取用这火。这火就是落在\Person{梅瑟}的祭献上并吞噬了祭品的那个火，正如经上所载：\BibleQuote{有火从主那里降下，吞噬了祭坛上的全燔祭。} \BibleRef{肋 9:24} 祭献必须只用这火来祝圣。因此，也有火从主那里降在\Person{亚郎}的儿子们身上，他们想要献上凡火，火便吞噬了他们，他们的尸首被扔到营外。\BibleRef{肋 10:2}\par

\Person{耶肋米亚}来到一个地方，发现一所像洞穴的屋子，就把帐幕、约柜和香坛搬了进去，并封闭了入口。那些跟来的人仔细查看那地方，要作记号，却无法找着。\Person{耶肋米亚}知道他们的意图，就说：\BibleQuote{这地方将一直不为人所知，直到天主再次聚集祂的子民，恩待他们。那时，天主就会显示这些事物，主的尊威也要显现。} \BibleRef{加下 2:5}\par

% structure: p0138 source=h2 target=subsection reason=relative-heading-rank; base=h2

\subsection{第十八章}

\Emph{在对已提及事件，尤其是\Person{乃赫米雅}所献祭献的叙述中，圣神和基督徒的圣洗被\OriginalTerm{预表}{type}。\Person{梅瑟}与\Person{厄里亚}的祭献以及\Person{诺厄}的历史也同样指向这同一奥迹。}\par

我们就是主的会众。我们承认我们主天主的赎罪之祭，就是我们的赎罪者\Person{耶稣基督}在祂的苦难中所完成的。我觉得，当我们读到主耶稣用圣神和火施洗时（正如\Person{若望}在他的福音中所说的），我们也不能忽略那火。\BibleRef{若 1:33} 祭品被焚尽，实在合宜，因为它原是为罪的。但那火是圣神的预像，圣神将在主升天后降临，赦免众人的罪，并像火一样点燃人的心灵和信者的心。因此，\Person{耶肋米亚}领受圣神后说：\BibleQuote{它在我心中犹如烈火，燃烧在我的骨骸中，我疲惫不堪，无法再忍。} \BibleRef{耶 20:9} 在\BibleBook{宗徒}{大事录}中，当圣神降临于宗徒和那些等候父所许诺的人身上时，我们也读到有火舌分别停留在他们头上。\BibleRef{宗 2:3} 他们的灵魂被圣神的影响所充满，以至于有人认为他们喝醉了新酒，而他们其实是领受了语言多样化的恩赐。\BibleRef{宗 2:13}\par

这还能有什么别的意思呢——就是说，火变成了水，水又引出了火——岂不是指\OriginalTerm{灵性恩宠}{spiritual grace}用火焚尽我们的罪，又用水洗净它们吗？因为罪被洗去，也被焚去。因此，宗徒说：\BibleQuote{火要试验各人的工程怎样。} \BibleRef{格前 3:13} 又说：\BibleQuote{谁的工程若被焚毁，他就要受亏损；他自己固然可以得救，却是像从火中经过一样。} \BibleRef{格前 3:15}\par

所以，我们这样论述，是为了证明罪是借着火被焚尽的。现在我们知道，这确实是当时那神圣的火，它作为未来赦罪的\OriginalTerm{预像}{type}，降在祭品之上。\par

这火在流徙期间（那时罪恶为王）是隐藏的，但在自由之时就被带了出来。它虽然变成了水的形状，却仍保持火的性质以焚尽祭品。当你读到天主父说：\BibleQuote{我是吞灭的火。} \BibleRef{申 4:24} 又说：\BibleQuote{他们离弃了我这活水的泉源。} \BibleRef{耶 2:13} 时，不要惊讶。主耶稣也像火一样炽热了听众的心，又像水泉一样清凉他们。因为祂在福音中亲口说过，祂来是把火投在地上 \BibleRef{路 12:49}，又要给口渴的人喝活水。 \BibleRef{若 7:37}\par

在\Person{厄里亚}时代，当他挑战异教先知们不用火就点燃祭台时，火也降下来了。他们做不到，他就三次把水倒在祭品上，水流满了祭台四周的水沟；然后他呼求，主的火就从天上降下，焚尽了全燔祭。\par

你就是那祭品。静默地默观其中的每一点。圣神的气息降在你身上，当祂焚尽你的罪时，祂似乎焚烧了你。\Person{梅瑟}时代所焚尽的祭品是赎罪祭，因此，正如\BibleBook{玛加伯}{下}所记载的，\Person{梅瑟}说：\BibleQuote{因为赎罪祭不应被吃掉，所以就烧了。} \BibleRef{加下 2:11} 它岂不是为你被焚尽吗？因为在圣洗圣事中，整个外在的人死去。宗徒大声说：\BibleQuote{我们的旧人已与祂同钉在十字架上了。} \BibleRef{罗 6:6} 这样，正如先祖的榜样教导我们的，埃及人就被吞没——希伯来人则因圣神而更新，如同他也干脚渡过红海——在那里我们的先祖都曾在云中和海中受了洗。 \BibleRef{格前 10:1}\par

在\Person{诺厄}时代的洪水中，凡有血肉的都死了，只有义人\Person{诺厄}和他全家得以保全。\BibleRef{创 7:23} 当一切必朽坏的生命被剪除时，一个人岂不是被焚尽了吗？外在的人被毁坏，内在的人却得以更新。不仅在圣洗中，而且在悔改中，肉身的毁坏也导向精神的增长，正如我们因宗徒的权威所领教的。圣\Person{保禄}说：\Quote{我以为我好像亲自在场一样，已判断了行这事的人，将这样的人交与撒殚，为使他的肉身毁灭，好使他的灵魂在主耶稣的日子可以得救。}\par

我们为了更充分地揭示那已启示给我们的圣事，并为思考这奇妙\OriginalTerm{奥迹}{mystery}，似乎进行了一段有些冗长的离题，因为这圣事确实充满德能与虔诚敬畏。\par

% structure: p0148 source=h2 target=subsection reason=relative-heading-rank; base=h2

\subsection{第十九章}

\Emph{\Place{基贝亚}居民对某\Person{肋未人}之妻所犯的罪行被叙述出来，而从他们采取的报复行动，可以推知古时人们心中必定充满了怎样的德性观念。}\par

我们的祖先何等注重德性，竟为一位妇女因受淫乱之辈的羞辱，而发动战争来报复！甚至当百姓战败后，他们发誓不将自己的女儿嫁给\Person{本雅明}支派！那个支派本来没有希望延续后裔，若非他们因急需而获准使用欺骗手段。并且这种许可似乎不缺乏对抢夺罪行施以相应的惩罚，因为他们只被允许通过抢夺而非\OriginalTerm{婚姻圣事}{sacrament of marriage}来成婚。的确，那些曾经破坏他人婚姻交合的人，理应失去自己的婚姻礼仪。\par

111. 这故事充满多少悲惨的特点！经上说，\BibleRef{民 19:1}，有一个肋未人，为自己娶了一个妻；我猜想她被称为妾，是来自\Quote{\OriginalTerm{交合}{concubitus}}一词。后来过了些日子，如常会发生的那样，她因某些事生气，回到父亲那里，与他同住了四个月。然后她丈夫起来，往岳父家去，要与妻子和好，赢回她，带她回家。那女人跑来迎接他，领丈夫进了父亲的家。\par

112. 少女\EditorialNote{民 4-9}的父亲欢欢喜喜地出来迎接他，那人与他同住了三天，一同吃喝休息。次日，肋未人天一亮就起来，却被岳父挽留，不想这么快就失去与他作伴的乐趣。再过了一天又一天，少女的父亲不容女婿动身，直到他们的欢乐和彼此的情谊达到圆满。但到了第七天，天色将晚，在愉快地用过饭后，老人以夜幕将临为由，劝他宁可在朋友中间安歇，而不在陌生人中间；但他无法再留他，便让他带着女儿一同离去。\par

113. 他们走了一小段路程，\BibleRef{民 19:10}，虽夜幕即将降临，他们已靠近\Place{耶步斯}人的城，仆人请求主人转向那里，但他拒绝了，因为那不是以色列子民的城。他希望赶到\Place{基贝亚}，那里居住着本雅明支派的人。但他们到了那里，却没有人接待他们以示好客，只有一个年纪很老的旅客——他看见他们，问肋未人说：\Quote{你往哪里去？从哪里来？} 肋未人回答说他正在赶路，往\Place{厄弗辣因山}去，而且没有人接待他，老人便款待了他，预备了饭食。\par

114. 他们吃饱了，\BibleRef{民 19:22}，桌子撤去之后，有些卑鄙的人冲上来围住了房子。老人便将他的女儿，一个处女，和与她同床的妾，交给这些恶人，只求他们不要向客人施暴。但是道理无用，强暴占了上风，肋未人便把妻子推出去，他们便强认了她，整夜凌辱她。她被这残忍或所受冤屈的悲伤压倒，倒在主人门口——她丈夫曾从那里进去——就断了气，以生命的最后努力守护着贤妻的情感，至少为丈夫保全了她的遗体。\par

115. 这事传开后（长话短说），几乎全体以色列人，都奋起作战。战事悬而未决，胜负未定，但是在第三次交战中，本雅明人被交在以色列人手中，\BibleRef{民 20:48}，他们被天主的审判定罪，为自己的放荡受到了惩罚。此外，判决规定，父家中的任何人都不得将女儿嫁给本雅明人为妻。这判决以庄严的誓言加以确认。但后来因为对弟兄们作出了如此严厉的判决而心生懊悔，他们便缓和了严厉态度，把那些失去父母、父亲因自己的罪而被杀的少女嫁给他们，或者让他们通过抢夺的办法找到妻子。由于这种如此卑鄙恶行的罪过，那些侵犯他人婚姻权利的人，被证明不配求娶婚姻。但为了怕一个支派从百姓中消灭，他们便默许了这种欺骗。\par

116. 我们的祖先如何高度尊重德行，以下事实可以证明：四万人拔剑攻击自己本雅明支派的弟兄，渴望为所受的羞耻报仇，因为他们不能容忍贞洁受到侵犯。于是，在那场战争中，双方共倒下六万五千勇士，他们的城池也被焚毁。当初以色列人虽然战败，却因对战争逆转毫无惧色，他们不顾为贞洁报仇所需的代价。他们冲入战场，准备用自己的鲜血洗去所犯罪行的污点。\par

% structure: p0157 source=h2 target=subsection reason=relative-heading-rank; base=h2

\subsection{第二十章}

\Emph{在\Place{撒玛黎雅}的可怕围城按照\Person{厄里叟}的预言结束后，他叙述了四个癞病人对德行的重视。}\par

117. 既然连癞病人——正如我们在\WorkTitle{列王纪}中所读到的——都关心什么是德行，我们又何必惊奇，主的子民竟也关注合宜与德行呢？\par

118. \Place{撒玛黎雅}发生了大饥荒，因为\Place{叙利亚}人的军队围困了那城。国王焦虑地在城墙上巡视卫兵，这时有一个妇人向他喊叫说：\Quote{这妇人劝我把我的儿子交出来——我就把他交出来，我们把他煮了吃了。她还答应我以后会把她的儿子带来，让我们一起来吃他的肉；可是现在她却把自己的儿子藏起来，不带来了。} 国王烦恼，因为这些妇人似乎不仅以人肉为食，而且还吃了自己儿女的肉；他为这种可怕悲惨的例证所动，便威胁要杀死先知\Person{厄里叟}。因为他相信，他本有能力解围免荒；或者他生气，因为先知没有允许国王击杀那些被他弄瞎眼的叙利亚人。\par

119. \Person{厄里叟}与长老们坐在\Place{贝特耳}，国王的使者还没有到他那里，他就对长老们说：\BibleQuote{你们看，这凶手之子竟派人来斩我的头！} 然后使者进来了，带来了国王的命令，威胁他立即处死。先知回答他说：\BibleQuote{明天大约这时候，在\Place{撒玛黎雅}城门口，一细亚细面要卖一协刻耳，两细亚大麦也要卖一协刻耳。} 这时，国王派来的使者不信，说：\BibleQuote{即使主打开天上的闸门，降下谷雨，也决不会发生这样的事！} \Person{厄里叟}就对他说：\BibleQuote{你既不相信，你必亲眼看见，却不得吃。}\par

120. 忽然，在\Place{叙利亚}人的营盘中，听到了仿佛战车的声音、战马嘶鸣和大军的喧嚣，以及一场大战的骚动；\Place{叙利亚}人以为\Place{以色列}王召集了\Place{埃及}王和\Place{阿摩黎人}的王来帮助他作战，便在黎明时逃走，舍下了营帐，因为他们害怕被这些新来的敌人突然袭击而溃败，无法抵挡诸王联合的力量。这事在\Place{撒玛黎雅}无人知晓，因为他们因恐惧和饥饿而身体虚弱，不敢出城。\par

121. 但在城门口有四个\OriginalTerm{癞病人}{lepers}，他们活着是痛苦，死了倒有益。他们彼此说：\Quote{看，我们坐在这里等死。如果我们进城，我们必死于饥荒；如果我们留在这里，也毫无活路。不如去\Place{叙利亚}人的营盘，他们或者会立刻杀死我们，或者给我们一条生路。}于是他们去了，进入\Place{叙利亚}人的营盘，看，全被敌人遗弃了。他们进入帐幕，先是找到食物，吃饱之后，又把所能找到的金银都拿起来。但他们虽然只顾着夺取战利品，却商议派人报告国王说\Place{叙利亚}人已经逃走了，因为他们认为这样比隐瞒消息、把欺骗得来的财物据为己有更有德行。\par

122. 人民听到这消息，就出去掠夺\Place{叙利亚}人的营盘。敌人的物资带来了丰收，使粮食变得极其便宜，正如先知的话：\BibleQuote{一细亚细面值一\OriginalTerm{协刻耳}{shekel}，二细亚大麦也值一协刻耳。}在人民这般欢乐之中，那个国王扶着手肘的军官死了，他在众人争先恐后地出城或满载欢笑而归时，被踩踏而死。\par

% structure: p0165 source=h2 target=subsection reason=relative-heading-rank; base=h2

\subsection{第二十一章}

\Emph{王后\Person{艾斯德尔}在其生命危险中追随了美德的恩宠；不，甚至一位异教国王也是如此，当死亡威胁到他最亲密的朋友时。因为友谊必须永远与美德结合，正如\Person{约纳堂}和\Person{阿希默肋客}的榜样所显示的。}\par

123. 为什么王后\Person{艾斯德尔}将自己置于死地，而不怕凶残国王的愤怒？\BibleRef{艾 4:16}这不正是为了拯救她的人民免于死亡，一件既合宜又有德行的事吗？\Place{波斯}王本人虽然凶残而骄傲，却也认为向那个曾告发谋害他的阴谋的人表示敬意，是合宜的，\BibleRef{艾 6:10}以解救一个自由民族脱离奴役，从死亡中夺回他们，而不宽恕那个推行如此不合宜计划的人。所以他最后将那个地位仅次于他、并在所有朋友中视为首脑的人，送上了绞架，\BibleRef{艾 7:9}因为他认为此人用邪恶的计谋羞辱了他。\par

124. 因为那种维护美德的值得称赞的友谊，无疑应该胜于财富、荣誉或权力。诚然，它不应超过美德，但应紧随其后。\Person{约纳堂}就是这样，他为了友情，不避讳他父亲的不悦，也不顾自身的安危。\Person{阿希默肋客}也是如此，他为了履行款待的义务，认为宁可忍受死亡，也不愿背叛自己逃亡中的朋友。\par

% structure: p0169 source=h2 target=subsection reason=relative-heading-rank; base=h2

\subsection{第二十二章}

\Emph{美德绝不可为朋友而放弃。然而，若必须为朋友作证反对他，则应谨慎行事。朋友之间，开诚相见需要何等诚恳，受苦需要何等宽宏大量，规过需要何等自由！友谊是诸美德的守护者，这些美德只能在品格相似的人身上找到。友谊在责善时应温和，且不追求自己的利益；因此，真正的朋友在富人中间少有。友谊的尊严是什么？朋友的背信弃义，较之他人更为恶劣，也更为可憎，这从\Person{犹达斯}和\Person{约伯}的朋友们的例子中可以看出来。}\par

125. 因此，绝不能将任何事物置于美德之前；为了使友谊的渴望不至于将美德抛弃，\WorkTitle{圣经}也就友谊的问题给我们提出了警告。事实上，哲学家们提出了各种问题；例如，一个人是否应当为了朋友而图谋反对自己的国家，以便为朋友服务？背弃信义以帮助并维护朋友的利益，是否正当？\par

126. \WorkTitle{圣经}也这样说：\BibleQuote{棍棒、刀剑和利箭，怎样对待人，作假见证陷害朋友的人，对他也是如此。}\BibleRef{箴 25:18}但请注意它补充的内容。它责备的不是为朋友作的见证，而是虚假的见证。因为如果为了天主的事业或者自己国家的事务，不得不作见证呢？难道友谊应该高过我们的宗教，或者我们对同胞的爱吗？然而，在这些事情中，需要真实的见证，这样朋友才不至被朋友的背信所害，而凭着朋友的诚信，他本应得到开脱。因此，人决不应该取悦渴望行恶的朋友，或者图谋陷害无辜的人。\par

127. 当然，如果必须作证，那么，当知道朋友的任何过失时，应当私下规劝他——如果他不听，就必须公开规劝。因为规劝是好的，往往比沉默的友谊更好。即使朋友觉得受了伤害，仍要规劝他；即使严厉的纠正伤了他的心，仍要规劝他，不要害怕。\BibleQuote{朋友的伤痕，胜于仇人的亲吻。}\BibleRef{箴 27:6}所以，规劝你犯错的朋友；不可抛弃无辜的人。因为友谊应当是恒常的，建立在真诚感情之上。我们不应出于幼稚的爱好或某种无谓的想法，就更换自己的朋友。\par

向朋友敞开胸怀，好使他忠于你，并使你从他那里获得生命的愉悦。\BibleQuote{忠实的朋友，是生命的妙药，不朽的恩宠。}\BibleRef{德 6:16} 对待朋友应如对待平等之人，要谦让，不要羞于在朋友之前履行友善的职责。因为友谊不知骄傲为何物。因此智者说：\BibleQuote{不要羞于向朋友致意。}\BibleRef{德 22:25} 在患难之时，不可离弃朋友，不可抛弃他，也不可辜负他，因为友谊是生命的支撑。那么，让我们按宗徒的教导彼此分担重担：\BibleRef{迦 6:2} 因为他这话是对那些被同一身体的\OriginalTerm{爱德}{charity}所拥抱的人说的。如果朋友在顺境中互相帮助，为何在逆境中不提供支持呢？让我们以劝告相助，献出我们最大的努力，全心与他们同情。\par

如有必要，我们甚至应为朋友忍受艰难。为了朋友的清白，常须承受敌意；若为受责难和控告的朋友辩护和答辩，也常会遭受辱骂。不要怕这样的不悦，因为义人的声音说：\BibleQuote{纵有灾祸临于我身，为朋友的原故，我也要忍受。}\BibleRef{德 22:26} 在逆境中方能考验朋友，因为在顺境中，人人似乎都是朋友。然而，正如逆境中需要忍耐和坚忍，在顺境中也需要有力的影响，以约束并驳斥那变得傲慢的朋友的骄横。\par

\Person{约伯}在逆境中多么高贵地说：\BibleQuote{我的朋友，你们可怜可怜我罢！}\BibleRef{约 19:21} 这并非可怜的哀号，而是责备。因为当他受到朋友们不公正的责难时，他回答说：\BibleQuote{我的朋友，你们可怜可怜我罢！} 意思是，你们本应表示怜悯，但你们反而攻击并压倒一个你们本该因友谊而同情其苦难的人。\par

所以，我的孩子们，要持守你们与弟兄们开始的友谊，因为在世上没有什么比这更美好了。今生能有一个你可以敞开心扉、可以分享信任、可以托付心中秘密的人，确实是一种安慰。有一个可靠的人在身边，在顺境中与你一同喜乐，在患难中同情你，在迫害中鼓励你，这是一种安慰。那些希伯来少年是多么好的朋友啊，火窑的烈焰也不能将他们彼此的爱分离！关于他们，我们已经谈过。圣\Person{达味}说得好：\Quote{撒乌耳与约纳堂，相亲相爱，生时不离，死时也不分离。}\par

这就是友谊的果实；因此，不可为了友谊而将信德弃置一旁。对天主不忠的人，不能成为人的朋友。友谊是怜悯的守护者，平等的导师，它使尊高者与卑微者平等，卑微者与尊高者平等。因为性格迥异的人之间不可能有友谊，因此彼此的好意应当相互契合。若事情需要，不要让卑微者缺少权威，也不要让尊高者缺少谦卑。让他倾听对方，如同倾听同等地位的人——一个平等的人，也让对方如同朋友般劝诫和责备，不是出于炫耀之心，而是出于深厚的爱。\par

你的劝诫不可严厉，你的责备不可尖刻，因为友谊既应避免谄媚，也应远离傲慢。朋友是什么？岂不是爱的伴侣？你与他结合，使灵魂与他相连，与他交融，渴望二人成为一体；你把自己托付给他，如同托付给第二个自我，对他无所畏惧，也不为自己的利益而要求任何不名誉的事。友谊不是用来牟利的，而是充满端庄，充满\OriginalTerm{恩宠}{grace}。友谊是一种德行，而非生财之道。它不因金钱而生，而因尊重而生；不因酬报的许诺，而因双方争相施惠的竞逐。\par

最后，贫穷者的友谊通常比富有者的更好，而且常有富人无友，而穷人朋友众多。因为真正的友谊不能存在于虚伪的谄媚之处。许多人谄媚地试图取悦富人，却没有人愿意假装讨好穷人。对穷人所说的一切都是真实的，他的友谊没有嫉妒。\par

有什么比既为天使又为人类共享的友谊更宝贵呢？因此，主耶稣说：\BibleQuote{要用不义的钱财结交朋友，好叫他们在你们匮乏的时候，收留你们到永远的居所。}\BibleRef{路 16:9} 天主亲自使我们成为朋友，而非仆人，正如祂自己所说：\BibleQuote{你们如果实行我所命令你们的，就是我的朋友。}\BibleRef{若 15:14} 祂给了我们一个可效法的友谊典范。我们要满足朋友的愿望，将自己心中的秘密向他敞开，也不可忽视他的信任。让我们向他敞开心扉，他也会向我们敞开。因此祂说：\BibleQuote{我称你们为朋友，因为凡由我父听来的一切，我都显示给你们了。}\BibleRef{若 15:15} 那么，一个真正的朋友，无所隐藏；他倾吐自己的灵魂，如同主耶稣倾吐了祂父的奥秘。\par

136. 所以，遵行天主圣意的人，便是祂的朋友，并以此名受尊荣。与祂同心合意的人，也是祂的朋友。因为在朋友之间，心意相通，而伤害友谊的人，最为可憎。因此，对于那位背叛者，主发现谴责其背叛的最严重之处在于，他毫无感激之情，竟在友谊的餐桌上掺入恶意的毒液。所以祂说：\BibleQuote{是你，与我同心的人，我的向导，我的知交，曾与我一同享用佳肴。}这意思是：这是不可容忍的，因为你攻击了那曾赐你\OriginalTerm{恩宠}{grace}的主。\BibleQuote{如果我的仇敌辱骂我，我还可以忍受，我也会躲避那恨我的人。}仇敌可以躲避；朋友却不可，倘他图谋陷害。对于不托付计划的人，我们可以防备；对于已然托付了计划的人，我们便无法防备了。因此，为彰显这罪的可憎，祂没有说：\Quote{你，我的仆人，我的宗徒}；却说：\Quote{你，与我同心的人}；即是说：你并非我的，而是你自己的背叛者，因为你背叛了一个与你同心的人。\par

137. 主自己，当祂对那三位不敬重圣约伯的首领发怒时，却愿意通过他们的朋友来宽恕他们，为使友谊的祈祷能赢得罪过的赦免。因此，约伯祈求，天主就宽恕了。友谊帮助了那些为傲慢所害的人。\BibleRef{约 42:7}\par

138. 我的孩子们，我把这些事留给你们，你们要将它们铭记在心——你们自己将会证实它们是否有益。同时，它们为你们提供了大量的榜样，因为几乎所有的榜样，都取自我们的先祖，连同他们许多的言语，都包含在这三卷书中；因此，尽管文辞或许不够优美，然而在如此狭小的篇幅中，排列着古老的榜样，仍能提供丰富的教诲。\par

\SourceRecord{Translated by H. de Romestin, E. de Romestin and H.T.F. Duckworth.；Nicene and Post-Nicene Fathers, Second Series；Vol. 10.；Edited by Philip Schaff and Henry Wace.；Buffalo, NY: Christian Literature Publishing Co.,；1896.；<http://www.newadvent.org/fathers/34013.htm>.}
